% Options for packages loaded elsewhere
\PassOptionsToPackage{unicode}{hyperref}
\PassOptionsToPackage{hyphens}{url}
\documentclass[a4paper]{extarticle}
\usepackage{fix-cm}
\usepackage[fontsize=13pt]{scrextend}
\usepackage{xcolor}
\usepackage[margin=2.5cm,top=2.5cm,bottom=2.5cm,left=2.5cm,right=2.5cm,headheight=1.25cm,headsep=0.5cm,footskip=1.25cm]{geometry}
\usepackage{amsmath,amssymb}
\setcounter{secnumdepth}{-\maxdimen} % remove section numbering
\usepackage{iftex}
\ifPDFTeX
  \usepackage[T5,T1]{fontenc}
  \usepackage[utf8]{inputenc}
  \usepackage{textcomp} % provide euro and other symbols
  \IfFileExists{newtxtext.sty}{\usepackage{newtxtext,newtxmath}}{\usepackage{lmodern}}
\else % if luatex or xetex
  \usepackage{fontspec}
  \IfFontExistsTF{Times New Roman}{\setmainfont{Times New Roman}}{\setmainfont{TeX Gyre Termes}}
  \setmonofont{Consolas}[Scale=MatchLowercase]
\fi
% Use upquote if available, for straight quotes in verbatim environments
\IfFileExists{upquote.sty}{\usepackage{upquote}}{}
\IfFileExists{microtype.sty}{% use microtype if available
  \usepackage[]{microtype}
  \UseMicrotypeSet[protrusion]{basicmath} % disable protrusion for tt fonts
}{}
\makeatletter
\@ifundefined{KOMAClassName}{% if non-KOMA class
  \IfFileExists{parskip.sty}{%
    \usepackage{parskip}
  }{% else
    \setlength{\parindent}{0pt}
    \setlength{\parskip}{6pt plus 2pt minus 1pt}}
}{% if KOMA class
  \KOMAoptions{parskip=half}}
\makeatother
\usepackage{color}
\usepackage{fancyvrb}
\newcommand{\VerbBar}{|}
\newcommand{\VERB}{\Verb[commandchars=\\\{\}]}
\DefineVerbatimEnvironment{Highlighting}{Verbatim}{commandchars=\\\{\}}
\RecustomVerbatimEnvironment{verbatim}{Verbatim}{fontsize=\small}
% Add ',fontsize=\small' for more characters per line
\newenvironment{Shaded}{}{}
\newcommand{\AlertTok}[1]{\textcolor[rgb]{1.00,0.00,0.00}{\textbf{#1}}}
\newcommand{\AnnotationTok}[1]{\textcolor[rgb]{0.38,0.63,0.69}{\textbf{\textit{#1}}}}
\newcommand{\AttributeTok}[1]{\textcolor[rgb]{0.49,0.56,0.16}{#1}}
\newcommand{\BaseNTok}[1]{\textcolor[rgb]{0.25,0.63,0.44}{#1}}
\newcommand{\BuiltInTok}[1]{\textcolor[rgb]{0.00,0.50,0.00}{#1}}
\newcommand{\CharTok}[1]{\textcolor[rgb]{0.25,0.44,0.63}{#1}}
\newcommand{\CommentTok}[1]{\textcolor[rgb]{0.38,0.63,0.69}{\textit{#1}}}
\newcommand{\CommentVarTok}[1]{\textcolor[rgb]{0.38,0.63,0.69}{\textbf{\textit{#1}}}}
\newcommand{\ConstantTok}[1]{\textcolor[rgb]{0.53,0.00,0.00}{#1}}
\newcommand{\ControlFlowTok}[1]{\textcolor[rgb]{0.00,0.44,0.13}{\textbf{#1}}}
\newcommand{\DataTypeTok}[1]{\textcolor[rgb]{0.56,0.13,0.00}{#1}}
\newcommand{\DecValTok}[1]{\textcolor[rgb]{0.25,0.63,0.44}{#1}}
\newcommand{\DocumentationTok}[1]{\textcolor[rgb]{0.73,0.13,0.13}{\textit{#1}}}
\newcommand{\ErrorTok}[1]{\textcolor[rgb]{1.00,0.00,0.00}{\textbf{#1}}}
\newcommand{\ExtensionTok}[1]{#1}
\newcommand{\FloatTok}[1]{\textcolor[rgb]{0.25,0.63,0.44}{#1}}
\newcommand{\FunctionTok}[1]{\textcolor[rgb]{0.02,0.16,0.49}{#1}}
\newcommand{\ImportTok}[1]{\textcolor[rgb]{0.00,0.50,0.00}{\textbf{#1}}}
\newcommand{\InformationTok}[1]{\textcolor[rgb]{0.38,0.63,0.69}{\textbf{\textit{#1}}}}
\newcommand{\KeywordTok}[1]{\textcolor[rgb]{0.00,0.44,0.13}{\textbf{#1}}}
\newcommand{\NormalTok}[1]{#1}
\newcommand{\OperatorTok}[1]{\textcolor[rgb]{0.40,0.40,0.40}{#1}}
\newcommand{\OtherTok}[1]{\textcolor[rgb]{0.00,0.44,0.13}{#1}}
\newcommand{\PreprocessorTok}[1]{\textcolor[rgb]{0.74,0.48,0.00}{#1}}
\newcommand{\RegionMarkerTok}[1]{#1}
\newcommand{\SpecialCharTok}[1]{\textcolor[rgb]{0.25,0.44,0.63}{#1}}
\newcommand{\SpecialStringTok}[1]{\textcolor[rgb]{0.73,0.40,0.53}{#1}}
\newcommand{\StringTok}[1]{\textcolor[rgb]{0.25,0.44,0.63}{#1}}
\newcommand{\VariableTok}[1]{\textcolor[rgb]{0.10,0.09,0.49}{#1}}
\newcommand{\VerbatimStringTok}[1]{\textcolor[rgb]{0.25,0.44,0.63}{#1}}
\newcommand{\WarningTok}[1]{\textcolor[rgb]{0.38,0.63,0.69}{\textbf{\textit{#1}}}}
\usepackage{longtable,booktabs,array}
\usepackage{caption}
\newcounter{none} % for unnumbered tables
\newcolumntype{L}[1]{>{\raggedright\arraybackslash\hyphenpenalty=10000\exhyphenpenalty=50}p{#1}}
\newcolumntype{C}[1]{>{\centering\arraybackslash}p{#1}}
\usepackage{calc} % for calculating minipage widths
% Correct order of tables after \paragraph or \subparagraph
\usepackage{etoolbox}
\makeatletter
\patchcmd\longtable{\par}{\if@noskipsec\mbox{}\fi\par}{}{}
\makeatother
% Allow footnotes in longtable head/foot
\IfFileExists{footnotehyper.sty}{\usepackage{footnotehyper}}{\usepackage{footnote}}
\makesavenoteenv{longtable}
\providecommand{\tabularnewline}{\\}
\newlength{\thesisdefaulttabcolsep}
\setlength{\thesisdefaulttabcolsep}{2pt}
\usepackage{graphicx}
\usepackage{tikz}
\usetikzlibrary{calc}
\makeatletter
\newsavebox\pandoc@box
\newcommand*\pandocbounded[1]{% scales image to fit in text height/width
  \sbox\pandoc@box{#1}%
  \Gscale@div\@tempa{\textheight}{\dimexpr\ht\pandoc@box+\dp\pandoc@box\relax}%
  \Gscale@div\@tempb{\linewidth}{\wd\pandoc@box}%
  \ifdim\@tempb\p@<\@tempa\p@\let\@tempa\@tempb\fi% select the smaller of both
  \ifdim\@tempa\p@<\p@\makebox[\linewidth][c]{\scalebox{\@tempa}{\usebox\pandoc@box}}%
  \else\makebox[\linewidth][c]{\usebox\pandoc@box}%
  \fi%
}
% Set default figure placement to htbp
\def\fps@figure{htbp}
\makeatother
\usepackage{xstring}
\newcommand{\safeincludesvg}[1]{%
  \StrSubstitute{#1}{.svg}{.png}[\safeincludepngpath]%
  \StrSubstitute{#1}{./assets/figures/}{./svg-inkscape/}[\safeincludepdfbasepath]%
  \StrSubstitute{\safeincludepdfbasepath}{.svg}{_svg-raw.pdf}[\safeincludepdfpath]%
  \IfFileExists{\safeincludepngpath}{%
    \includegraphics[width=0.96\linewidth,height=0.78\textheight,keepaspectratio]{\safeincludepngpath}%
  }{%
    \IfFileExists{\safeincludepdfpath}{%
      \includegraphics[width=0.96\linewidth,height=0.78\textheight,keepaspectratio]{\safeincludepdfpath}%
    }{%
      \fbox{\parbox{0.9\linewidth}{\centering Missing figure assets for:\\ \texttt{\detokenize{#1}}}}%
    }%
  }%
}
\newcommand{\safeincludesvgsmall}[1]{%
  \StrSubstitute{#1}{.svg}{.png}[\safeincludepngpath]%
  \StrSubstitute{#1}{./assets/figures/}{./svg-inkscape/}[\safeincludepdfbasepath]%
  \StrSubstitute{\safeincludepdfbasepath}{.svg}{_svg-raw.pdf}[\safeincludepdfpath]%
  \IfFileExists{\safeincludepngpath}{%
    \includegraphics[width=0.76\linewidth,height=0.56\textheight,keepaspectratio]{\safeincludepngpath}%
  }{%
    \IfFileExists{\safeincludepdfpath}{%
      \includegraphics[width=0.76\linewidth,height=0.56\textheight,keepaspectratio]{\safeincludepdfpath}%
    }{%
      \fbox{\parbox{0.9\linewidth}{\centering Missing figure assets for:\\ \texttt{\detokenize{#1}}}}%
    }%
  }%
}
\newcommand{\safeincludesvgcompact}[1]{%
  \StrSubstitute{#1}{.svg}{.png}[\safeincludepngpath]%
  \StrSubstitute{#1}{./assets/figures/}{./svg-inkscape/}[\safeincludepdfbasepath]%
  \StrSubstitute{\safeincludepdfbasepath}{.svg}{_svg-raw.pdf}[\safeincludepdfpath]%
  \IfFileExists{\safeincludepngpath}{%
    \includegraphics[width=0.68\linewidth,height=0.46\textheight,keepaspectratio]{\safeincludepngpath}%
  }{%
    \IfFileExists{\safeincludepdfpath}{%
      \includegraphics[width=0.68\linewidth,height=0.46\textheight,keepaspectratio]{\safeincludepdfpath}%
    }{%
      \fbox{\parbox{0.9\linewidth}{\centering Missing figure assets for:\\ \texttt{\detokenize{#1}}}}%
    }%
  }%
}
\newcommand{\safeincludesvgnarrow}[1]{%
  \StrSubstitute{#1}{.svg}{.png}[\safeincludepngpath]%
  \StrSubstitute{#1}{./assets/figures/}{./svg-inkscape/}[\safeincludepdfbasepath]%
  \StrSubstitute{\safeincludepdfbasepath}{.svg}{_svg-raw.pdf}[\safeincludepdfpath]%
  \IfFileExists{\safeincludepngpath}{%
    \includegraphics[width=0.56\linewidth,height=0.34\textheight,keepaspectratio]{\safeincludepngpath}%
  }{%
    \IfFileExists{\safeincludepdfpath}{%
      \includegraphics[width=0.56\linewidth,height=0.34\textheight,keepaspectratio]{\safeincludepdfpath}%
    }{%
      \fbox{\parbox{0.9\linewidth}{\centering Missing figure assets for:\\ \texttt{\detokenize{#1}}}}%
    }%
  }%
}
\newcommand{\safeincludesvgchart}[1]{%
  \StrSubstitute{#1}{.svg}{.png}[\safeincludepngpath]%
  \StrSubstitute{#1}{./assets/figures/}{./svg-inkscape/}[\safeincludepdfbasepath]%
  \StrSubstitute{\safeincludepdfbasepath}{.svg}{_svg-raw.pdf}[\safeincludepdfpath]%
  \IfFileExists{\safeincludepngpath}{%
    \includegraphics[width=0.72\linewidth,height=0.28\textheight,keepaspectratio]{\safeincludepngpath}%
  }{%
    \IfFileExists{\safeincludepdfpath}{%
      \includegraphics[width=0.72\linewidth,height=0.28\textheight,keepaspectratio]{\safeincludepdfpath}%
    }{%
      \fbox{\parbox{0.9\linewidth}{\centering Missing figure assets for:\\ \texttt{\detokenize{#1}}}}%
    }%
  }%
}
\setlength{\emergencystretch}{10em} % reduce overfull lines
\setlength{\tabcolsep}{\thesisdefaulttabcolsep}
\setlength{\arrayrulewidth}{0.5pt}
\captionsetup[table]{justification=centering,singlelinecheck=false,labelfont=it,textfont=it}
\captionsetup[longtable]{justification=centering,singlelinecheck=false,labelfont=it,textfont=it}
\setlength{\LTcapwidth}{\linewidth}
\setlength{\LTleft}{0pt}
\setlength{\LTright}{0pt}
% Pandoc-generated longtables in this thesis render correctly but emit repeated
% box warnings during the output routine; raise fuzz thresholds to silence this
% known-noise without changing the visible layout.
\setlength{\hfuzz}{460pt}
\setlength{\vfuzz}{360pt}
\hbadness=10000
\setcounter{topnumber}{2}
\setcounter{bottomnumber}{2}
\setcounter{totalnumber}{4}
\renewcommand{\topfraction}{0.85}
\renewcommand{\bottomfraction}{0.75}
\renewcommand{\textfraction}{0.1}
\renewcommand{\floatpagefraction}{0.75}
\usepackage{setspace}
\setstretch{1.15}
\setlength{\parindent}{0pt}
\setlength{\parskip}{6pt}
\setcounter{tocdepth}{3}
\renewcommand{\arraystretch}{1.15}
\usepackage{titlesec}
\titleformat{\section}{\bfseries\Large\centering}{}{0pt}{}
\titlespacing*{\section}{0pt}{12pt}{6pt}
\titleformat{\subsection}{\bfseries\normalsize}{}{0pt}{}
\titlespacing*{\subsection}{0pt}{12pt}{0pt}
\titleformat{\paragraph}[block]{\bfseries\normalsize}{}{0pt}{}
\titlespacing*{\paragraph}{0pt}{10pt}{0pt}
\titleformat{\subparagraph}[block]{\bfseries\normalsize}{}{0pt}{}
\titlespacing*{\subparagraph}{0pt}{8pt}{0pt}
\newcommand{\thesisfrontmatterentry}[1]{%
  \phantomsection
  \addcontentsline{toc}{section}{#1}%
}
\newcommand{\majorheading}[1]{%
  \clearpage
  \phantomsection
  \section*{#1}%
  \addcontentsline{toc}{section}{#1}%
}
\newcommand{\chapterheading}[1]{%
  \clearpage
  \phantomsection
  \section*{#1}%
  \addcontentsline{toc}{section}{#1}%
  \setcounter{subsection}{0}%
}
\newcommand{\appendixsection}[1]{\clearpage\section{#1}}
\newcommand{\frontmatterstart}{\clearpage\hypersetup{pageanchor=true}\pagenumbering{roman}\setcounter{page}{1}}
\newcommand{\mainmatterstart}{\clearpage\pagenumbering{arabic}\setcounter{page}{1}}
\newcommand{\checkbox}{\(\square\)}
\newcommand{\thesisspacer}{\hspace{2em}}
\newcommand{\figref}[1]{Hình~\ref{#1}}
\newcommand{\tabref}[1]{Bảng~\ref{#1}}
\definecolor{thesisphenikaablue}{RGB}{42, 63, 130}
\newcommand{\thesislogopath}{./assets/brand/phenikaa-university-logo.jpeg}
\newcommand{\thesisfrontseparator}{\begin{center}\rule{0.5\linewidth}{0.5pt}\end{center}}
\newcommand{\thesisdrawtemplateframe}{%
  \begin{tikzpicture}[remember picture,overlay]
    \draw[line width=1.4pt] ($(current page.north west)+(1.00cm,-1.00cm)$) rectangle ($(current page.south east)+(-1.00cm,1.00cm)$);
    \draw[line width=0.5pt] ($(current page.north west)+(1.14cm,-1.14cm)$) rectangle ($(current page.south east)+(-1.14cm,1.14cm)$);
  \end{tikzpicture}%
}
\newcommand{\thesisdrawsidebarframe}{%
  \begin{tikzpicture}[remember picture,overlay]
    \draw[line width=0.5pt] ($(current page.north west)+(1.35cm,-1.35cm)$) rectangle ($(current page.south east)+(-1.35cm,1.35cm)$);
    \draw[line width=0.5pt] ($(current page.north west)+(3.65cm,-1.35cm)$) -- ($(current page.south west)+(3.65cm,1.35cm)$);
  \end{tikzpicture}%
}
\newcommand{\thesiscoverstart}{%
  \clearpage
  \newgeometry{top=1.5cm,bottom=1.5cm,left=1.6cm,right=1.6cm}
  \thispagestyle{empty}
}
\newcommand{\thesiscoverend}{%
  \clearpage
  \restoregeometry
}
\newcommand{\thesisofficialheader}{%
  \noindent
  \begin{minipage}[t]{0.48\linewidth}
    \centering
    \textbf{BỘ GIÁO DỤC VÀ ĐÀO TẠO}\\
    \textbf{ĐẠI HỌC PHENIKAA}\\
    \rule{0.72\linewidth}{0.6pt}
  \end{minipage}\hfill
  \begin{minipage}[t]{0.48\linewidth}
    \centering
    \textbf{CỘNG HÒA XÃ HỘI CHỦ NGHĨA VIỆT NAM}\\
    \textbf{Độc lập - Tự do - Hạnh phúc}\\
    \rule{0.72\linewidth}{0.6pt}
  \end{minipage}
  \par\vspace{0.6cm}
}
\newcommand{\thesissignaturepair}[4]{%
  \vspace{0.8cm}
  \noindent
  \begin{minipage}[t]{0.48\linewidth}
    \centering
    \textbf{#1}\\
    \emph{(Ký, ghi rõ họ tên)}\\[2.8cm]
    #2
  \end{minipage}\hfill
  \begin{minipage}[t]{0.48\linewidth}
    \centering
    \textbf{#3}\\
    \emph{(Ký, ghi rõ họ tên)}\\[2.8cm]
    #4
  \end{minipage}
}
\renewcommand{\figurename}{Hình}
\renewcommand{\tablename}{Bảng}
\renewcommand{\listtablename}{DANH MỤC BẢNG - LIST OF TABLES}
\renewcommand{\listfigurename}{DANH MỤC HÌNH ẢNH VÀ ĐỒ THỊ - LIST OF PICTURES AND GRAPHS}
\makeatletter
\newcommand{\thesissetfigureid}[2]{%
  \renewcommand{\thefigure}{#1}%
  \renewcommand{\theHfigure}{#2}%
}
\newcommand{\thesissettableid}[2]{%
  \renewcommand{\thetable}{#1}%
  \renewcommand{\theHtable}{#2}%
}
\newcommand{\thesistableofcontentsbody}{\@starttoc{toc}}
\newcommand{\thesislistoftablesbody}{\@starttoc{lot}}
\newcommand{\thesislistoffiguresbody}{\@starttoc{lof}}
\renewcommand*\l@table{\@dottedtocline{1}{0em}{5.8em}}
\renewcommand*\l@figure{\@dottedtocline{1}{0em}{5.8em}}
\makeatother
\newcounter{ieeeref}
\newcommand{\ieeebibitem}[2]{%
  \refstepcounter{ieeeref}\noindent[\theieeeref]\label{#1} #2\par
}
\newcommand{\ieeecite}[1]{[\ref{#1}]}
\providecommand{\textcite}[1]{\ieeecite{#1}}
\providecommand{\parencite}[1]{\ieeecite{#1}}
\providecommand{\tightlist}{%
  \setlength{\itemsep}{0pt}\setlength{\parskip}{0pt}}
\usepackage{bookmark}
\IfFileExists{xurl.sty}{\usepackage{xurl}}{} % add URL line breaks if available
\hyphenation{Mos-quit-to Post-gre-SQL Vic-to-ria-Met-rics Vic-to-ria-Logs Au-then-ti-ca-tion Ku-ber-ne-tes Os-cil-lo-scope Dash-board Web-Socket}
\urlstyle{same}
\hypersetup{
  hidelinks,
  linktoc=all,
  bookmarksnumbered=true,
  pdfcreator={LaTeX via pandoc}}

\author{}
\date{}

\begin{document}
\hypersetup{pageanchor=false}
\pagenumbering{gobble}

\thesiscoverstart
\thesisdrawsidebarframe
\noindent
\begin{minipage}[t][0.92\textheight][t]{0.16\textwidth}
  \centering
  \vspace*{1.0cm}
  \rotatebox{90}{\textbf{\Large LÊ TRỌNG AN}}\\[1.5cm]
  \rotatebox{90}{\textbf{\large KỸ THUẬT CƠ ĐIỆN TỬ}}
\end{minipage}\hfill
\begin{minipage}[t][0.92\textheight][t]{0.77\textwidth}
  \begin{center}
    \textbf{\Large BỘ GIÁO DỤC VÀ ĐÀO TẠO}\\[0.2cm]
    \textbf{\Large ĐẠI HỌC PHENIKAA}\\[-0.05cm]
    \rule{0.34\linewidth}{0.6pt}

    \vspace{1.7cm}
    \includegraphics[width=0.36\linewidth]{\thesislogopath}

    \vspace{1.5cm}
    {\fontsize{28}{30}\selectfont\textbf{ĐỒ ÁN TỐT NGHIỆP}}\\[1.0cm]
    {\fontsize{17}{21}\selectfont\textbf{THIẾT KẾ HỆ THỐNG IOT CHO ỨNG DỤNG QUẢN LÝ PHƯƠNG TIỆN}}\\
    {\fontsize{17}{21}\selectfont\textbf{GIAO THÔNG TRONG LĨNH VỰC CHO THUÊ XE TỰ LÁI}}


    \vfill
    \begin{minipage}{0.90\linewidth}
      \raggedright
      \textbf{Sinh viên:} Lê Trọng An\\[0.15cm]
      \textbf{Mã số sinh viên:} 21010389 \hfill \textbf{Khóa:} K15\\[0.15cm]
      \textbf{Ngành:} Kỹ thuật Cơ Điện tử \hfill \textbf{Hệ:} Đại học chính quy\\[0.15cm]
      \textbf{Giảng viên hướng dẫn:} TS. Nguyễn Đức Nam
    \end{minipage}

    \vfill
    \textbf{\Large Hà Nội - 2026}
  \end{center}
\end{minipage}
\thesiscoverend

\thesiscoverstart
\thesisdrawtemplateframe
\begin{center}
  \textbf{\Large BỘ GIÁO DỤC VÀ ĐÀO TẠO}\\[0.15cm]
  \textbf{\Large ĐẠI HỌC PHENIKAA}\\[-0.05cm]
  \rule{0.28\linewidth}{0.6pt}

  \vspace{1.6cm}
  \includegraphics[width=0.30\linewidth]{\thesislogopath}

  \vspace{1.4cm}
  {\fontsize{27}{30}\selectfont\textbf{ĐỒ ÁN TỐT NGHIỆP}}\\[0.9cm]
  {\fontsize{18}{22}\selectfont\textbf{THIẾT KẾ HỆ THỐNG IOT CHO ỨNG DỤNG QUẢN LÝ}}\\
  {\fontsize{18}{22}\selectfont\textbf{PHƯƠNG TIỆN GIAO THÔNG TRONG LĨNH VỰC}}\\
  {\fontsize{18}{22}\selectfont\textbf{CHO THUÊ XE TỰ LÁI}}

  \vfill
  \begin{minipage}{0.78\linewidth}
    \raggedright
    \textbf{Sinh viên:} Lê Trọng An\\[0.18cm]
    \textbf{Mã số sinh viên:} 21010389 \hfill \textbf{Khóa:} K15\\[0.18cm]
    \textbf{Ngành:} Kỹ thuật Cơ Điện tử \hfill \textbf{Hệ:} Đại học chính quy\\[0.18cm]
    \textbf{Giảng viên hướng dẫn:} TS. Nguyễn Đức Nam
  \end{minipage}

  \vspace{1.3cm}
  \textbf{\Large Hà Nội – 2026}
\end{center}
\thesiscoverend

\thesiscoverstart
\thesisdrawtemplateframe
\begin{center}
  {\color{thesisphenikaablue}\textbf{\Large PHENIKAA UNIVERSITY}}\\[0.18cm]
  {\color{thesisphenikaablue}\textbf{\Large FACULTY OF MECHANICAL ENGINEERING AND MECHATRONICS}}\\[0.25cm]
  \includegraphics[width=0.28\linewidth]{\thesislogopath}

  \vspace{1.0cm}
  {\fontsize{22}{26}\selectfont\textbf{MEM710071 – CAPSTONE DESIGN 2}}\\[0.55cm]
  {\fontsize{18}{22}\selectfont\textbf{DESIGN OF AN IOT SYSTEM FOR VEHICLE MANAGEMENT}}\\
  {\fontsize{18}{22}\selectfont\textbf{IN SELF-DRIVE CAR RENTAL SERVICES}}\\[0.35cm]
  {\fontsize{15}{18}\selectfont\textbf{Semester: I Academic year: 2025--2026}}

  \vfill
  \begin{minipage}{0.76\linewidth}
    \raggedright
    \textbf{Project Team:}
    \begin{center}
      \begin{tabular}{|p{0.42\linewidth}|p{0.20\linewidth}|p{0.12\linewidth}|}
       \\\hline
        \centering\textbf{Student name} & \centering\textbf{Student ID} & \centering\textbf{Cohort}\tabularnewline
       \\\hline
        Le Trong An & \centering 21010389 & \centering K15\tabularnewline
       \\\hline
      \end{tabular}
    \end{center}
    \vspace{0.2cm}
    \textbf{Course Instructor:} Dr. Nguyen Duc Nam
  \end{minipage}

  \vfill
  \begin{minipage}{0.88\linewidth}
    \centering
    This report is submitted to the Faculty of Mechanical Engineering and Mechatronics,\\
    School of Engineering, Phenikaa University in partial fulfillment of the requirements\\
    of the Graduation Thesis course MEM710071.
  \end{minipage}

  \vspace{0.8cm}
  \textbf{\Large HANOI – 2026}
\end{center}
\thesiscoverend

\thesiscoverstart
\thesisdrawtemplateframe
\begin{center}
  {\color{thesisphenikaablue}\textbf{\Large ĐẠI HỌC PHENIKAA}}\\[0.15cm]
  {\color{thesisphenikaablue}\textbf{\Large KHOA CƠ KHÍ – CƠ ĐIỆN TỬ}}\\[0.25cm]
  \includegraphics[width=0.27\linewidth]{\thesislogopath}

  \vspace{0.8cm}
  {\fontsize{21}{25}\selectfont\textbf{MEM710071}}\\[0.15cm]
  {\fontsize{26}{30}\selectfont\textbf{ĐỒ ÁN TỐT NGHIỆP}}\\[0.45cm]
  {\fontsize{18}{22}\selectfont\textbf{THIẾT KẾ HỆ THỐNG IOT CHO ỨNG DỤNG QUẢN LÝ}}\\
  {\fontsize{18}{22}\selectfont\textbf{PHƯƠNG TIỆN GIAO THÔNG TRONG LĨNH VỰC CHO THUÊ XE TỰ LÁI}}\\[0.35cm]
  {\fontsize{15}{18}\selectfont\textbf{Học kỳ I Năm học 2025--2026}}

  \vfill
  \begin{center}
    \begin{tabular}{|p{0.34\linewidth}|p{0.24\linewidth}|p{0.10\linewidth}|}
     \\\hline
      \centering\textbf{Họ và tên sinh viên} & \centering\textbf{Mã số sinh viên} & \centering\textbf{Lớp}\tabularnewline
     \\\hline
      Lê Trọng An & \centering 21010389 & \centering K15-KTCĐT2\tabularnewline
     \\\hline
    \end{tabular}
  \end{center}

  \vspace{0.45cm}
  \textbf{Giảng viên hướng dẫn:} TS. Nguyễn Đức Nam

  \vfill
  \textbf{\Large HÀ NỘI, 2026}
\end{center}
\thesiscoverend

\frontmatterstart

\thispagestyle{empty}
\thesisofficialheader
\begin{center}
  \textbf{\Large NHẬN XÉT ĐỒ ÁN TỐT NGHIỆP CỦA GIẢNG VIÊN HƯỚNG DẪN}
\end{center}
\thesisfrontmatterentry{NHẬN XÉT CỦA GIẢNG VIÊN HƯỚNG DẪN}

\textbf{Giảng viên hướng dẫn:} TS. Nguyễn Đức Nam \hfill \textbf{Khoa:} Cơ khí -- Cơ Điện tử.\\
\textbf{Tên đề tài:} Thiết kế hệ thống IoT cho ứng dụng quản lý phương tiện giao thông trong lĩnh vực cho thuê xe tự lái.\\
\textbf{Sinh viên thực hiện:} Lê Trọng An. \hfill \textbf{Lớp:} K15-KTCĐT2.

\begin{center}
  \textbf{NỘI DUNG NHẬN XÉT}
\end{center}
\textbf{I. Nhận xét Đồ án tốt nghiệp:}\\
- Nhận xét về hình thức: \dotfill\\[0.55cm]
- Tính cấp thiết của đề tài: \dotfill\\[0.55cm]
- Mục tiêu của đề tài: \dotfill\\[0.55cm]
- Nội dung của đề tài: \dotfill\\[0.55cm]
- Tài liệu tham khảo: \dotfill\\[0.55cm]
- Phương pháp nghiên cứu: \dotfill\\[0.55cm]
- Tính sáng tạo và ứng dụng: \dotfill\\[0.55cm]
\textbf{II. Nhận xét về tinh thần và thái độ làm việc của sinh viên:} \dotfill\\[0.55cm]
\textbf{III. Kết quả đạt được:} \dotfill\\[0.55cm]
\textbf{IV. Kết luận:} Đồng ý cho bảo vệ: \checkbox \thesisspacer Không đồng ý cho bảo vệ: \checkbox

\begin{flushright}
  Hà Nội, ngày \ldots, tháng \ldots năm 20\ldots\\
  \textbf{GIẢNG VIÊN HƯỚNG DẪN}\\
  \emph{(Ký, ghi rõ họ tên)}\\[2.6cm]
  TS. Nguyễn Đức Nam
\end{flushright}
\clearpage

\thispagestyle{empty}
\thesisofficialheader
\begin{center}
  \textbf{\Large NHẬN XÉT ĐỒ ÁN TỐT NGHIỆP CỦA GIẢNG VIÊN PHẢN BIỆN}
\end{center}
\thesisfrontmatterentry{NHẬN XÉT CỦA GIẢNG VIÊN PHẢN BIỆN}

\textbf{Giảng viên phản biện:} Theo quyết định phân công phản biện của Khoa \\
\textbf{Tên đề tài:} Thiết kế hệ thống IoT cho ứng dụng quản lý phương tiện giao thông trong lĩnh vực cho thuê xe tự lái.\\
\textbf{Sinh viên thực hiện:} Lê Trọng An \hfill \textbf{Lớp:} K15-KTCĐT2\\
\textbf{Giảng viên hướng dẫn:} TS. Nguyễn Đức Nam

\begin{center}
  \textbf{NỘI DUNG NHẬN XÉT}
\end{center}
\textbf{I. Nhận xét Đồ án tốt nghiệp:}\\
- Bố cục, hình thức trình bày: \dotfill\\[0.55cm]
- Đảm bảo tính cấp thiết, hiện đại, không trùng lặp: \dotfill\\[0.55cm]
- Ná»™i dung: \dotfill\\[0.55cm]
- Mức độ thực hiện: \dotfill\\[0.55cm]
\textbf{II. Kết quả đạt được:} \dotfill\\[0.55cm]
\textbf{III. Ưu nhược điểm:} \dotfill\\[0.55cm]
\textbf{IV. Kết luận:} Đồng ý cho bảo vệ: \checkbox \thesisspacer Không đồng ý cho bảo vệ: \checkbox

\begin{flushright}
  Hà Nội, ngày \ldots, tháng \ldots năm 20\ldots\\
  \textbf{GIẢNG VIÊN PHẢN BIỆN}\\
  \emph{(Ký, ghi rõ họ tên)}\\[2.6cm]
  Theo quyết định phân công phản biện của Khoa
\end{flushright}
\clearpage

\thispagestyle{empty}
\thesisofficialheader
\begin{center}
  \textbf{\Large BIÊN BẢN ĐÁNH GIÁ ĐỒ ÁN TỐT NGHIỆP}
\end{center}
\thesisfrontmatterentry{BIÊN BẢN ĐÁNH GIÁ ĐỒ ÁN TỐT NGHIỆP}

\textbf{I. Thông tin chung}
\begin{enumerate}
  \item Thời gian: từ \ldots giờ \ldots đến \ldots giờ \ldots ngày \ldots/\ldots/20\ldots
  \item Địa điểm: \dotfill
  \item Họ và tên sinh viên: Lê Trọng An \hfill Mã số SV: 21010389
  \item Lớp: K15-KTCĐT2 \hfill Ngành: Kỹ thuật Cơ Điện tử \hfill Khóa: K15
  \item Tên đề tài: Thiết kế hệ thống IoT cho ứng dụng quản lý phương tiện giao thông trong lĩnh vực cho thuê xe tự lái.
  \item Giảng viên hướng dẫn: TS. Nguyễn Đức Nam
\end{enumerate}

\textbf{II. Thành phần hội đồng}
\begin{enumerate}
  \item \dotfill Chủ tịch
  \item \dotfill Thư ký
  \item \dotfill Giảng viên phản biện
  \item \dotfill Ủy viên 1
\end{enumerate}

\textbf{III. Tổng hợp câu hỏi của thành viên hội đồng:} \dotfill\\[0.6cm]
\textbf{IV. Tổng hợp nội dung trả lời của sinh viên:} \dotfill\\[0.6cm]
\textbf{V. Nội dung đánh giá của Hội đồng}
\begin{enumerate}
  \item Ý nghĩa của đồ án: \dotfill
  \item Về nội dung, kết cấu của đồ án: \dotfill
  \item Phương pháp nghiên cứu: \dotfill
  \item Các kết quả nghiên cứu đạt được: \dotfill
\end{enumerate}
\textbf{VI. Điểm kết luận:} \dotfill\ (Bằng chữ: \dotfill)

\thesissignaturepair{CHỦ TỊCH HỘI ĐỒNG}{}{THƯ KÝ}{}
\clearpage

\thispagestyle{empty}
\begin{center}
  \textbf{\Large GIẢI TRÌNH CÁC CHỈNH SỬA}\\[0.2cm]
  \textit{(Nếu có)}
\end{center}
\thesisfrontmatterentry{GIẢI TRÌNH CÁC CHỈNH SỬA}

\vspace{0.5cm}
\begin{center}
  Chưa phát sinh nội dung cần giải trình.
\end{center}
\clearpage

\majorheading{LỜI CAM ĐOAN}

Tên tôi là: Lê Trọng An.\\
Mã sinh viên: 21010389 \hfill Lớp: K15-KTCĐT2\\
Ngành: Kỹ thuật Cơ Điện tử.

Tôi đã thực hiện đồ án tốt nghiệp với đề tài: Thiết kế hệ thống IoT cho ứng dụng quản lý phương tiện giao thông trong lĩnh vực cho thuê xe tự lái.

Tôi xin cam đoan đây là đề tài nghiên cứu của riêng tôi và được sự hướng dẫn của TS. Nguyễn Đức Nam.

Các nội dung nghiên cứu, kết quả trong đề tài này là trung thực và chưa được công bố dưới bất kỳ hình thức nào. Nếu phát hiện có bất kỳ hình thức gian lận nào, tôi xin hoàn toàn chịu trách nhiệm trước pháp luật.

\begin{flushright}
  Hà Nội, ngày \ldots, tháng \ldots năm 20\ldots
\end{flushright}

\thesissignaturepair{GIẢNG VIÊN HƯỚNG DẪN}{TS. Nguyễn Đức Nam}{SINH VIÊN}{Lê Trọng An}

\majorheading{TÓM TẮT ĐỒ ÁN TỐT NGHIỆP - ABSTRACT}

Sự phát triển nhanh của dịch vụ cho thuê xe tự lái tại Việt Nam làm gia
 tăng nhu cầu giám sát phương tiện từ xa. Nhu cầu này không chỉ dừng ở
việc xác định vị trí xe, mà còn bao gồm theo dõi trạng thái vận hành và
phát hiện sớm các dấu hiệu bất thường. Từ yêu cầu đó, đồ án trình bày
quá trình thiết kế và triển khai một hệ thống giám sát phương tiện theo
kiến trúc tích hợp, bao gồm thiết bị gắn trên xe, hạ tầng máy chủ tiếp
nhận dữ liệu và giao diện theo dõi dành cho người quản lý. Hệ thống này
nhằm hỗ trợ doanh nghiệp nắm được vị trí xe, trạng thái kỹ thuật và phản
ứng sớm với các rủi ro vận hành.

Ở tầng thiết bị, hệ thống được xây dựng trên bo mạch PCB tự thiết kế.
Trong đó, ESP32-S3 giữ vai trò bộ điều khiển trung tâm; SIM7600CE-T đảm
nhiệm truyền dữ liệu 4G và định vị vệ tinh GNSS để xác định vị trí xe;
LIS3DSH theo dõi rung động để nhận biết chuyển động khi xe đỗ; còn
adapter vgate iCar Pro cho phép đọc dữ liệu OBD2 từ xe qua BLE. Kiến
trúc nguồn được tổ chức theo nhiều
nhánh để thiết bị vẫn duy trì hoạt động khi xe tắt máy, đồng thời hạn
chế ảnh hưởng tới ắc quy chính.

Ở tầng phần mềm, dữ liệu từ thiết bị được gửi bằng MQTT 3.1.1 tới EMQX.
MQTT đóng vai trò lớp nhắn tin nhẹ phù hợp với đường truyền 4G, còn EMQX
là broker trung tâm tiếp nhận và phân phối bản tin. Từ đây, MQTT Bridge
subscribe bốn nhóm topic
\texttt{v1/\{device\_id\}/\{rawdata\textbar status\textbar events\textbar firmware\}},
chuẩn hóa payload rồi phân luồng dữ liệu tới PostgreSQL,
VictoriaMetrics và VictoriaLogs. Song song với luồng ingest, backend còn
lắng nghe nhánh \texttt{internal/events/\#} để đẩy cập nhật thời gian
thực ra giao diện qua Socket.IO.

Bảng điều khiển web được phát triển bằng Next.js 15 và React 19. Giao
 diện này có thể hiển thị bản đồ, biểu đồ và trạng thái phương tiện gần
như theo thời gian thực nhờ cơ chế đẩy dữ liệu mới từ máy chủ lên trình
duyệt qua Socket.IO. Nhờ đó, người vận hành không chỉ quan sát vị trí xe
mà còn có thể theo dõi lịch sử hành trình, cảnh báo và trạng thái thiết
bị trên cùng một giao diện thống nhất.

Kết quả đạt được là một hệ thống IoT giám sát phương tiện hoàn chỉnh từ
phần cứng đến phần mềm, có khả năng theo dõi vị trí thời gian thực, đọc
dữ liệu chẩn đoán OBD2, thiết lập vùng cho phép, phát sinh cảnh báo tự
động, cập nhật firmware từ xa và hỗ trợ quản lý đội xe. Quan trọng hơn,
hệ thống cho thấy một hướng triển khai phù hợp với bối cảnh doanh nghiệp
cho thuê xe tự lái trong nước: chi phí hợp lý, có thể tùy biến và đủ dư
địa để mở rộng trong các giai đoạn sau.

\textbf{Từ khóa:} IoT, giám sát phương tiện, GPS, OBD2, MQTT, ESP32-S3,
thời gian thực, vùng cho phép

\textbf{Abstract (English)}

As the self-drive car rental market in Vietnam grows, the need for
remote vehicle monitoring and real-time fleet management becomes
increasingly practical for operators. This thesis presents the design
and development of an end-to-end IoT vehicle tracking system, covering
the onboard device, cloud services, and operational dashboard for
real-time tracking, diagnostic data collection, and alerting.

On the hardware side, the tracker is implemented as a custom integrated
PCB rather than an assembly of off-the-shelf development boards. The
main board integrates ESP32-S3 (MCU), SIMCom SIM7600CE-T (LTE+GNSS),
LIS3DSH (IMU), DS3231M (RTC), W25Q128 (SPI Flash), and power blocks based
on MP2482, TPS54231, AP2112, SX1308, and TP5100. The modem is configured
to Auto mode (\texttt{AT+CNMP=2}) so it can switch between LTE/UMTS/GSM
automatically, uses the default APN \texttt{internet}, and delivers GNSS
data via \texttt{AT+CGNSINF} (with optional NMEA streaming through
\texttt{AT+CGNSTST}) on the same UART1 channel. The device reads engine
diagnostic data via OBD2 through a vgate iCar Pro adapter over Bluetooth
Low Energy (BLE), where this adapter is an external peripheral outside
the main PCB. The power architecture uses a 5V main rail, an
approximately 4V modem rail, a 3.3V logic rail, and a 18650 1S backup
cell to maintain operation when the vehicle engine is off.

On the software side, device data is transmitted through MQTT 3.1.1 via
EMQX with per-device ACL authorization. The MQTT Bridge subscribes to
the device topics
\texttt{v1/\{deviceId\}/\{rawdata\textbar status\textbar events\textbar firmware\}},
validates the payloads, and routes them to PostgreSQL, VictoriaMetrics,
and VictoriaLogs. The backend is built with Express.js and TypeScript in
a domain-oriented structure, exposes REST/OpenAPI endpoints, and listens
to internal MQTT events so Socket.IO channels can push realtime updates
and OTA status back to operators.

The web interface is developed using Next.js 15 and React 19, providing
a real-time dashboard with Leaflet maps, ECharts visualizations, and
WebSocket connectivity via Socket.IO for instant vehicle position and
status updates. The entire system is containerized and deployed using
Docker, ensuring consistency between development and production
environments.

The result is a complete end-to-end IoT vehicle tracking system, from
hardware to software, capable of real-time location tracking, OBD2
diagnostic data retrieval, allowed-zone monitoring, automated alerts,
remote firmware updates, and fleet-management support for self-drive car
rental services.

\textbf{Keywords:} IoT, vehicle tracking, GPS, OBD2, MQTT, ESP32-S3,
real-time, allowed-zone monitoring

\majorheading{LỜI CẢM ƠN - ACKNOWLEDGEMENTS}\label{lux1eddi-cux1ea3m-ux1a1n---acknowledgements}

Trước tiên, tôi xin gửi lời cảm ơn chân thành và sâu sắc nhất đến TS.
Nguyễn Đức Nam --- người đã trực tiếp hướng dẫn, định hướng và tận tình
hỗ trợ tôi trong suốt quá trình thực hiện đồ án tốt nghiệp này. Những
góp ý chuyên môn và sự động viên của thầy là nguồn động lực quan trọng
giúp tôi hoàn thành đề tài.

Tôi xin trân trọng cảm ơn quý thầy cô trong Khoa Cơ khí - Cơ Điện tử,
Trường Kỹ thuật - Đại học Phenikaa đã truyền đạt những kiến thức nền
tảng vững chắc trong suốt quá trình học tập, tạo điều kiện thuận lợi để
tôi có thể áp dụng vào thực tiễn thông qua đồ án này.

Tôi cũng xin gửi lời cảm ơn đến gia đình, bạn bè và các đồng nghiệp đã
luôn ủng hộ, chia sẻ và đồng hành cùng tôi trong suốt thời gian thực
hiện đề tài.

Mặc dù đã nỗ lực hoàn thiện đồ án trong phạm vi thời gian và kiến thức
cho phép, nội dung chắc chắn vẫn còn những thiếu sót. Tôi rất mong nhận
được những ý kiến góp ý từ quý thầy cô và hội đồng để tiếp tục hoàn
thiện đề tài.

Xin chân thành cảm ơn!

\begin{flushright}
  Hà Nội, ngày \ldots tháng \ldots năm 20\ldots\\[0.25cm]
  \textbf{SINH VIÊN THỰC HIỆN}\\[2.4cm]
  Lê Trọng An
\end{flushright}

\majorheading{MỤC LỤC - TABLE OF CONTENT}\label{mux1ee5c-lux1ee5c---table-of-contents}

\thesistableofcontentsbody

\majorheading{DANH MỤC BẢNG - LIST OF TABLES}\label{danh-mux1ee5c-bux1ea3ng---list-of-tables}

\thesislistoftablesbody

\majorheading{DANH MỤC HÌNH ẢNH VÀ ĐỒ THỊ - LIST OF PICTURES AND GRAPHS}\label{danh-mux1ee5c-huxecnh---list-of-figures}

\thesislistoffiguresbody

\majorheading{DANH MỤC TỪ VIẾT TẮT - LIST OF ABBREVIATION}\label{danh-mux1ee5c-tux1eeb-viux1ebft-tux1eaft---list-of-abbreviations}

{\def\LTcaptype{none} % do not increment counter
\setlength{\tabcolsep}{4pt}
\begin{longtable}{|L{0.160\linewidth}|L{0.410\linewidth}|L{0.360\linewidth}|}
\\\hline
\textbf{Từ viết tắt} & \textbf{Tiếng Anh} & \textbf{Tiếng Việt} \\
\\\hline
\endfirsthead
\\\hline
\textbf{Từ viết tắt} & \textbf{Tiếng Anh} & \textbf{Tiếng Việt} \\
\\\hline
\endhead
\\\hline
\endfoot
\endlastfoot
\textbf{ADC} & Analog-to-Digital Converter & Bộ chuyển đổi tương tự sang
số \\\hline
\textbf{API} & Application Programming Interface & Giao diện lập trình
ứng dụng \\\hline
\textbf{BLE} & Bluetooth Low Energy & Bluetooth năng lượng thấp \\\hline
\textbf{BOM} & Bill of Materials & Danh sách linh kiện \\\hline
\textbf{CI/CD} & Continuous Integration / Continuous Deployment & Tích
hợp liên tục / Triển khai liên tục \\\hline
\textbf{CORS} & Cross-Origin Resource Sharing & Cơ chế chia sẻ tài nguyên
khác nguồn \\\hline
\textbf{DDD} & Domain-Driven Design & Thiết kế hướng miền \\\hline
\textbf{Docker} & Docker Container Platform & Nền tảng đóng gói và chạy
ứng dụng bằng container \\\hline
\textbf{EMQX} & Erlang MQTT Broker & Broker MQTT xây dựng trên Erlang/OTP \\\hline
\textbf{ESP32} & Espressif System-on-Chip 32-bit & Hệ thống trên chip
32-bit của Espressif \\\hline
\textbf{FreeRTOS} & Free Real-Time Operating System & Hệ điều hành thời
gian thực mã nguồn mở \\\hline
\textbf{GPIO} & General Purpose Input/Output & Cổng vào/ra đa năng \\\hline
\textbf{GNSS} & Global Navigation Satellite System & Hệ thống vệ tinh
dẫn đường toàn cầu \\\hline
\textbf{GPS} & Global Positioning System & Hệ thống định vị toàn cầu \\\hline
\textbf{HTTP} & Hypertext Transfer Protocol & Giao thức truyền siêu văn
bản \\\hline
\textbf{HTTPS} & Hypertext Transfer Protocol Secure & Giao thức truyền
siêu văn bản bảo mật \\\hline
\textbf{IMU} & Inertial Measurement Unit & Đơn vị đo lường quán tính \\\hline
\textbf{IoT} & Internet of Things & Internet vạn vật \\\hline
\textbf{JWT} & JSON Web Token & Mã thông báo web định dạng JSON \\\hline
\textbf{LVD} & Low Voltage Disconnect & Ngắt điện áp thấp \\\hline
\textbf{MCU} & Microcontroller Unit & Vi điều khiển \\\hline
\textbf{MQTT} & Message Queuing Telemetry Transport & Giao thức truyền
bản tin nhẹ theo mô hình publish/subscribe \\\hline
\textbf{OBD2} & On-Board Diagnostics II & Chẩn đoán trên xe thế hệ 2 \\\hline
\textbf{ORM} & Object-Relational Mapping & Ánh xạ đối tượng - quan hệ \\\hline
\textbf{PCB} & Printed Circuit Board & Bảng mạch in \\\hline
\textbf{REST} & Representational State Transfer & Kiểu kiến trúc REST
cho dịch vụ web \\\hline
\textbf{RTOS} & Real-Time Operating System & Hệ điều hành thời gian
thá»±c \\\hline
\textbf{SoC} & System on Chip & Hệ thống trên chip \\\hline
\textbf{SQL} & Structured Query Language & Ngôn ngữ truy vấn có cấu
trúc \\\hline
\textbf{TLS} & Transport Layer Security & Bảo mật tầng vận chuyển \\\hline
\textbf{UART} & Universal Asynchronous Receiver-Transmitter & Bá»™ thu
phát bất đồng bộ đa năng \\\hline
\textbf{WebSocket} & WebSocket Protocol & Giao thức giao tiếp hai chiều
thời gian thực \\\hline
\end{longtable}
}

\begin{center}\rule{0.5\linewidth}{0.5pt}\end{center}

\mainmatterstart
\chapterheading{CHƯƠNG 1. GIỚI THIỆU DỰ ÁN -- PROJECT OVERVIEW}

\subsection{1.1. Đặt vấn đề / Bối cảnh của dự
án}\label{11-ux111ux1eb7t-vux1ea5n-ux111ux1ec1--bux1ed1i-cux1ea3nh-cux1ee7a-dux1ef1-uxe1n}

\subsubsection{1.1.1. Bối cảnh thị trường cho thuê xe tự lái tại Việt
Nam}\label{111-bux1ed1i-cux1ea3nh-thux1ecb-trux1b0ux1eddng-cho-thuuxea-xe-tux1ef1-luxe1i-tux1ea1i-viux1ec7t-nam}

Những năm gần đây, thị trường cho thuê xe tự lái tại Việt Nam tăng
trưởng rõ rệt, đặc biệt tại TP. Hồ Chí Minh, Hà Nội và Đà Nẵng. Theo báo
cáo ngành vận tải {[}1{]}, nhu cầu thuê xe tăng trung bình 15--20\% mỗi
năm nhờ sự phục hồi của du lịch nội địa, nhu cầu di chuyển linh hoạt và
xu hướng chia sẻ phương tiện.

Với doanh nghiệp vận hành đội xe, tốc độ tăng trưởng này mang lại cơ hội
kinh doanh nhưng cũng kéo theo một yêu cầu rất thực tế: phải biết xe
đang ở đâu, đang được sử dụng như thế nào và có an toàn hay không, thay
vì chỉ dựa vào các cuộc gọi xác nhận. Khi quy mô đội xe mở rộng, áp lực
giám sát vận hành và bảo toàn tài sản tăng lên nhanh chóng. Trong thực
tế, các doanh nghiệp cho thuê xe tự lái hiện nay thường gặp bốn nhóm khó
khăn sau:

\begin{itemize}
\tightlist
\item
  \textbf{Quản lý thủ công kém hiệu quả}: Nhiều doanh nghiệp quy mô vừa
  và nhỏ vẫn quản lý xe bằng sổ sách, gọi điện xác nhận hoặc dựa vào sự
  chủ động của khách hàng. Cách làm này không cung cấp thông tin thời
  gian thực về vị trí và trạng thái phương tiện {[}2{]}.
\item
  \textbf{Rủi ro mất cắp và sử dụng sai mục đích}: Khi không có hệ thống
  giám sát, xe có thể bị sử dụng vượt phạm vi địa lý đã thỏa thuận, chạy
  quá số km quy định, hoặc trong trường hợp xấu nhất là bị chiếm đoạt.
  Việc phát hiện các tình huống này thường bị trễ, gây thiệt hại lớn về
  tài sản {[}3{]}.
\item
  \textbf{Thiếu dữ liệu chẩn đoán kỹ thuật}: Doanh nghiệp khó theo dõi
  tình trạng kỹ thuật của xe từ xa, nên việc bảo trì thường mang tính bị
  động và chỉ được thực hiện khi sự cố đã xuất hiện. Điều này làm tăng
  chi phí sửa chữa và rút ngắn tuổi thọ phương tiện.
\item
  \textbf{Giải pháp thương mại đắt đỏ}: Các hệ thống định vị và giám sát xe thương
  mại hiện có trên thị trường (như Vietmap, iTracking) thường có chi phí
  cao --- bao gồm phí thiết bị, phí dịch vụ hàng tháng, và phí tích hợp
  --- không phù hợp với các doanh nghiệp quy mô nhỏ với ngân sách hạn
  chế {[}4{]}.
\end{itemize}

Những khó khăn trên cho thấy nhu cầu giám sát là rất rõ, nhưng chưa tự
động dẫn tới một lời giải khả thi. Muốn biến nhu cầu vận hành thành một
hệ thống thật sự dùng được, bài toán phải được chuyển thành các thách
thức kỹ thuật cụ thể.

\subsubsection{1.1.2. Thách thức kỹ
thuật}\label{112-thuxe1ch-thux1ee9c-kux1ef9-thuux1eadt}

Để giải quyết những khó khăn trên, hệ thống không thể chỉ là một thiết
bị định vị đơn giản. Bài toán thực tế là phải dung hòa nhiều yêu cầu kỹ
thuật vốn dễ xung đột với nhau: thiết bị phải tiết kiệm điện nhưng vẫn
luôn sẵn sàng, phải truyền dữ liệu liên tục nhưng vẫn chịu được mất
sóng, và phải đủ đơn giản để triển khai thực tế với chi phí hợp lý. Bốn
thách thức nổi bật là:

\textbf{Thứ nhất, bài toán tiêu thụ năng lượng.} Nếu thiết bị theo dõi
hoạt động liên tục ở trạng thái công suất cao, ắc quy xe có thể bị hao
đáng kể chỉ sau vài tuần, từ đó ảnh hưởng trực tiếp đến khả năng khởi
động {[}5{]}. Vì vậy, hệ thống phải biết khi nào cần hoạt động đầy đủ và
khi nào có thể chuyển sang chế độ tiết kiệm điện.

\textbf{Thứ hai, giám sát khi xe đỗ.} Khi xe tắt máy (IGN OFF), thiết bị
không thể duy trì mọi khối chức năng như lúc xe đang chạy. Tuy nhiên,
nó vẫn phải phát hiện được các chuyển động bất thường, chẳng hạn xe bị
dắt đi hoặc bị cẩu kéo. Đây là bài toán cân bằng giữa độ nhạy phát hiện
và mức tiêu thụ năng lượng {[}6{]}.

\textbf{Thứ ba, độ tin cậy của kết nối.} Hệ thống hoạt động trong môi
trường di động nên kết nối 4G/LTE có thể gián đoạn bất kỳ lúc nào. Điều
người quản lý cần không chỉ là "có mạng thì gửi được", mà là khi mất
mạng thiết bị vẫn không làm mất dữ liệu quan trọng và có thể đồng bộ lại
khi kết nối được phục hồi.

\textbf{Thứ tư, xử lý dữ liệu thời gian thực.} Khi số lượng xe tăng lên
hàng chục hoặc hàng trăm, dịch vụ phía máy chủ phải tiếp nhận liên tục
dữ liệu vị trí, thông số OBD2 và trạng thái cảm biến với độ trễ thấp,
đồng thời vẫn giữ giao diện giám sát dễ quan sát và dễ thao tác cho
người quản lý.

{\thesissetfigureid{1.1}{fig-ch1-overview-problem-solution}
\begin{figure}[htbp]
\centering
\pandocbounded{\safeincludesvg{./assets/figures/01-chuong-1-gioi-thieu-hinh-1-1.svg}}
\caption{Sơ đồ tổng quan vấn đề và giải pháp đề xuất}
\label{fig:ch1-overview-problem-solution}
\end{figure}
}

\subsubsection{1.1.3. Động lực thực hiện dự
án}\label{113-ux111ux1ed9ng-lux1ef1c-thux1ef1c-hiux1ec7n-dux1ef1-uxe1n}

Mục tiêu của đề tài không chỉ dừng ở việc chế tạo một thiết bị gửi tọa
độ GPS. Điều cần xây dựng là một hệ thống đủ hoàn chỉnh để người vận
hành đội xe có thể theo dõi trạng thái phương tiện, hiểu nguyên nhân của
các cảnh báo và đưa ra quyết định kịp thời. Từ nhu cầu thực tiễn đó, đề
tài triển khai hệ thống ``IoT Vehicle Tracking System'' nhằm xây dựng
một giải pháp đồng bộ, bao gồm thiết bị phần cứng trên xe, phần mềm điều
khiển thiết bị, hạ tầng máy chủ tiếp nhận và xử lý dữ liệu, cùng giao
 diện web để khai thác thông tin. Mục tiêu là tạo ra một hệ thống có chi
phí hợp lý, dễ tùy biến và có thể mở rộng theo quy mô doanh nghiệp cho
thuê xe tự lái tại Việt Nam.

\begin{center}\rule{0.5\linewidth}{0.5pt}\end{center}

\subsection{1.2. Mục tiêu và phạm vi của dự
án}\label{12-mux1ee5c-tiuxeau-vuxe0-phux1ea1m-vi-cux1ee7a-dux1ef1-uxe1n}

\subsubsection{1.2.1. Mục tiêu tổng
quát}\label{121-mux1ee5c-tiuxeau-tux1ed5ng-quuxe1t}

Mục tiêu chính của dự án là thiết kế và triển khai một hệ thống IoT theo
dõi phương tiện hoàn chỉnh, bao gồm cả phần cứng lẫn phần mềm, phục vụ
cho bài toán quản lý đội xe cho thuê tự lái. Các mục tiêu dưới đây được
xác định từ góc nhìn của đơn vị vận hành đội xe: họ cần thông tin đủ
nhanh để can thiệp, nhưng cũng cần hệ thống đủ bền bỉ để vận hành lâu
dài trên xe. Hệ thống cần đáp ứng các yêu cầu sau:

\begin{enumerate}
\def\labelenumi{\arabic{enumi}.}
\tightlist
\item
  \textbf{Theo dõi vị trí thời gian thực}: Cung cấp vị trí định vị vệ
  tinh GPS/GNSS của xe theo chu kỳ 5--30 giây khi xe đang di chuyển, để
  người quản lý nhìn thấy hành trình gần như liên tục trên bản đồ thay
  vì chỉ nhận các mốc vị trí rời rạc.
\item
  \textbf{Giám sát trạng thái kỹ thuật qua OBD2}: Kết nối với cổng chẩn
  đoán OBD-II của xe, nơi cung cấp các thông tin vận hành và lỗi cơ bản,
  thông qua BLE để đọc các thông số như tốc độ, vòng tua máy, nhiệt độ
  động cơ và mã lỗi chẩn đoán (DTC).
\item
  \textbf{Cảnh báo thông minh}: Tự động phát hiện và gửi thông báo khi
  có các sự kiện bất thường như xe vượt khỏi vùng giám sát đã vẽ trên
  bản đồ, vượt tốc độ quy định, chuyển động bất thường khi xe đang đỗ,
  hoặc mất kết nối thiết bị.
\item
  \textbf{Tối ưu hóa năng lượng}: Đảm bảo thiết bị hoạt động liên tục mà
  không làm cạn ắc quy xe, với thời lượng pin dự phòng đủ để hoạt động
  độc lập trong trường hợp mất nguồn chính.
\item
  \textbf{Giao diện quản lý trực quan}: Cung cấp giao diện web dạng bảng
  điều khiển tổng hợp để người quản lý đội xe theo dõi vị trí, xem lịch
  sử hành trình, quản lý cảnh báo và tạo báo cáo.
\end{enumerate}

Năm mục tiêu trên cho biết hệ thống cần đạt được điều gì. Phần tiếp theo
sẽ chốt rõ phạm vi nào được hiện thực trong đồ án và phạm vi nào được
giữ lại cho các giai đoạn phát triển sau.

\subsubsection{1.2.2. Phạm vi dự
án}\label{122-phux1ea1m-vi-dux1ef1-uxe1n}

Để bảo đảm phạm vi nghiên cứu tập trung, nội dung triển khai được giới
hạn theo từng tầng của hệ thống, từ đối tượng ứng dụng đến phần cứng,
firmware và hạ tầng máy chủ. Cách xác định phạm vi này giúp người đọc
thấy rõ đâu là nội dung đã được triển khai trong đồ án và đâu là nội
dung được giữ lại cho các giai đoạn phát triển sau.

\textbf{Đối tượng ứng dụng:}

\begin{itemize}
\tightlist
\item
  Loại xe: Ô tô con (xe 4 chỗ, sedan, SUV) với hệ thống điện 12V hoặc
  24V DC
\item
  Ắc quy phổ thông: 40--70 Ah (trung bình 45--60 Ah)
\item
  Đối tượng sử dụng: Các công ty cho thuê xe tự lái quy mô vừa và nhỏ
  tại Việt Nam
\end{itemize}

\textbf{Phạm vi phần cứng:}

\begin{itemize}
\tightlist
\item
  Thiết bị theo dõi IoT sử dụng vi điều khiển ESP32-S3
\item
  Modem LTE + GNSS SIMCom SIM7600CE-T (Auto mode LTE/UMTS/GSM, APN mặc
  định \texttt{internet}) kết nối trực tiếp qua UART1 để cung cấp cả dữ
  liệu 4G và GNSS
\item
  Adapter OBD2 BLE vgate iCar Pro (đọc dữ liệu chẩn đoán xe qua
  Bluetooth)
\item
  Cảm biến gia tốc LIS3DSH (IMU) để phát hiện chuyển động và rung
\item
  Pin dự phòng 18650 1S với mạch sạc và bảo vệ
\end{itemize}

\textbf{Phạm vi firmware:}

\begin{itemize}
\tightlist
\item
  Xây dựng phần mềm trên thiết bị bằng ESP-IDF kết hợp FreeRTOS để điều
  phối các tác vụ nền theo thời gian thực
\item
  Giao tiếp BLE với adapter OBD2 để đọc dữ liệu chẩn đoán của xe
\item
  Truyền dữ liệu lên máy chủ bằng MQTT 3.1.1 qua EMQX, phân tầng QoS 0/1
  theo loại bản tin
\item
  Quản lý năng lượng theo nhiều trạng thái vận hành, gồm lúc xe đang
  chạy, lúc xe đỗ và lúc cần phản ứng khẩn
\end{itemize}

\textbf{Phạm vi máy chủ và dịch vụ xử lý dữ liệu:}

\begin{itemize}
\tightlist
\item
  EMQX làm broker MQTT, nhận và chuyển tiếp dữ liệu từ thiết bị
\item
  MQTT Bridge (dịch vụ trung gian độc lập) xử lý và phân phối dữ liệu tới
  đúng nơi lưu trữ
\item
  PostgreSQL lưu trữ dữ liệu quan hệ (người dùng, xe, khách hàng, cảnh
  báo)
\item
  VictoriaMetrics lưu trữ dữ liệu chuỗi thời gian
\item
  VictoriaLogs lưu trữ nhật ký sự kiện
\item
  API máy chủ được triển khai bằng Express.js và được tổ chức theo từng
  nhóm nghiệp vụ để dễ mở rộng
\item
  Giao diện web quản trị được triển khai bằng Next.js 15, kết hợp bản đồ
  và biểu đồ để người dùng theo dõi vận hành
\end{itemize}

\textbf{Giới hạn phạm vi nghiên cứu:}

\begin{itemize}
\tightlist
\item
  Không xử lý các chức năng chẩn đoán OBD-II chuyên sâu (chỉ đọc các
  thông số cơ bản)
\item
  Không phát hiện va chạm (IMU chỉ dùng cho phát hiện chuyển động/rung)
\item
  Không điều khiển động cơ xe (tắt bơm xăng, khóa xe) --- chỉ gửi cảnh
  báo
\item
  Không bao gồm ứng dụng di động ở giai đoạn đầu (dự kiến làm ở giai
  đoạn 2 bằng Flutter và tận dụng lại giao diện web hiện có)
\end{itemize}

{\thesissettableid{1.1}{tab-ch1-project-scope-by-layer}
\def\LTcaptype{table}
\begin{longtable}{|L{0.261\linewidth}|L{0.470\linewidth}|L{0.234\linewidth}|}
\caption{Tóm tắt phạm vi dự án theo từng tầng hệ thống}\label{tab:ch1-project-scope-by-layer}\\
\hline \textbf{Tầng hệ thống} & \textbf{Công nghệ chính} & \textbf{Phạm vi} \\
\\\hline
\endfirsthead
\caption[]{Tóm tắt phạm vi dự án theo từng tầng hệ thống (tiếp theo)}\\
\hline \textbf{Tầng hệ thống} & \textbf{Công nghệ chính} & \textbf{Phạm vi} \\
\\\hline
\endhead
\\\hline
\endfoot
\endlastfoot
Phần cứng & ESP32-S3, SIMCom SIM7600CE-T (LTE + GNSS), vgate iCar Pro,
LIS3DSH & Thiết kế PCB, chế tạo và tích hợp thiết bị hoàn chỉnh \\\hline
Firmware & ESP-IDF, FreeRTOS, MQTT 3.1.1, OTA qua HTTPS & Phần mềm trên
thiết bị: đọc dữ liệu, quản lý năng lượng và gửi các gói dữ liệu nhỏ lên
máy chủ \\\hline
MQTT Broker & EMQX 5.x & Broker MQTT nhận và chuyển tiếp dữ liệu từ thiết bị \\\hline
Dịch vụ máy chủ & Express.js, TypeScript, PostgreSQL, VictoriaMetrics & API
phía máy chủ và xử lý dữ liệu \\\hline
Giao diện web & Next.js 15, React 19, Leaflet, ECharts & Bản đồ, biểu
đồ và cảnh báo cho người vận hành \\\hline
Hạ tầng triển khai & Docker, Nginx Proxy Manager, Grafana & Đóng gói
dịch vụ, công bố truy cập và giám sát vận hành \\\hline
\end{longtable}
}

\begin{center}\rule{0.5\linewidth}{0.5pt}\end{center}

\subsection{1.3. Các tiêu chí cần đạt được của dự
án}\label{13-cuxe1c-tiuxeau-chuxed-cux1ea7n-ux111ux1ea1t-ux111ux1b0ux1ee3c-cux1ee7a-dux1ef1-uxe1n}

Dự án đặt ra các tiêu chí cụ thể cho từng tầng của hệ
thống. Các tiêu chí này không chỉ dùng để đánh giá kết quả cuối cùng, mà
còn đóng vai trò như một cam kết thiết kế xuyên suốt quá trình thực hiện:
mỗi quyết định chọn linh kiện, tổ chức firmware hay xây dựng hạ tầng đều
phải quay lại trả lời câu hỏi liệu nó có giúp hệ thống vận hành đủ ổn
định, đủ tiết kiệm năng lượng và đủ hữu ích cho bài toán quản lý xe hay
không.

\subsubsection{1.3.1. Tiêu chí phần
cứng}\label{131-tiuxeau-chuxed-phux1ea7n-cux1ee9ng}

Ở tầng phần cứng, thiết bị cần đồng thời đáp ứng yêu cầu về tiêu thụ điện
thấp, độ ổn định nhiệt và an toàn đối với ắc quy xe.


{\thesissettableid{1.3.1A}{tab-ch1-tong-hop-thong-tin-ve-stt-tieu-chi-muc-tieu-cu-the}
\def\LTcaptype{table}
\begin{longtable}{|L{0.160\linewidth}|L{0.224\linewidth}|L{0.580\linewidth}|}
\caption{Tổng hợp thông tin về STT, Tiêu chí, Mục tiêu cụ thể}\label{tab:ch1-tong-hop-thong-tin-ve-stt-tieu-chi-muc-tieu-cu-the}\\
\hline \textbf{STT} & \textbf{Tiêu chí} & \textbf{Mục tiêu cụ thể} \\
\\\hline
\endfirsthead
\caption[]{Tổng hợp thông tin về STT, Tiêu chí, Mục tiêu cụ thể (tiếp theo)}\\
\hline \textbf{STT} & \textbf{Tiêu chí} & \textbf{Mục tiêu cụ thể} \\
\\\hline
\endhead
\\\hline
\endfoot
\endlastfoot
1 & Tiêu thụ điện ở chế độ ngủ sâu (deep sleep, tức gần như tắt toàn bộ
khối không cần thiết) & \textless{} 500 µA \\\hline
2 & Tiêu thụ điện chế độ hoạt động & \textless{} 250 mA (trung bình) \\\hline
3 & Thời gian thức dậy từ chế độ ngủ sâu & \textless{} 3 giây \\\hline
4 & Thời lượng pin dự phòng (18650 1S) & \textgreater= 2.5 giờ theo dõi
liên tục; \textgreater= 24 giờ cảnh báo gián đoạn (mục tiêu lấy theo
mức pin tham chiếu của cell 18650 đã chọn) \\\hline
5 & Ngưỡng chuyển nguồn bảo vệ ắc quy & Cấu hình cho hệ 12V:
\(OFF = 12.0\,\mathrm{V}\), \(ON = 12.2\,\mathrm{V}\); Profile 24V:
\(OFF = 24.0\,\mathrm{V}\), \(ON = 24.4\,\mathrm{V}\) \\\hline
6 & Nhiệt độ hoạt động & -10 C đến +60 C \\\hline
\end{longtable}
}


\subsubsection{1.3.2. Tiêu chí
firmware}\label{132-tiuxeau-chuxed-firmware}

Ở tầng firmware, tiêu chí không chỉ dừng ở việc đọc được cảm biến hay gửi
được dữ liệu, mà nằm ở khả năng điều phối đúng nhịp giữa các khối để thiết
bị phản ứng phù hợp với từng trạng thái vận hành của xe.


{\thesissettableid{1.2A}{tab-ch1-tieu-chi-firmware}
\def\LTcaptype{table}
\begin{longtable}{|L{0.160\linewidth}|L{0.404\linewidth}|L{0.400\linewidth}|}
\caption{Tiêu chí firmware}\label{tab:ch1-tieu-chi-firmware}\\
\hline \textbf{STT} & \textbf{Tiêu chí} & \textbf{Mục tiêu cụ thể} \\
\\\hline
\endfirsthead
\caption[]{Tiêu chí firmware (tiếp theo)}\\
\hline \textbf{STT} & \textbf{Tiêu chí} & \textbf{Mục tiêu cụ thể} \\
\\\hline
\endhead
\\\hline
\endfoot
\endlastfoot
1 & Tần suất gửi GPS khi lái xe & 5--30 giây (cấu hình được) \\\hline
2 & Tần suất heartbeat khi đỗ xe & 10--30 phút \\\hline
3 & Thời gian kết nối OBD2 BLE & \textless{} 10 giây \\\hline
4 & Độ chính xác vị trí GNSS & \textless{} 5 mét (điều kiện thoáng) \\\hline
5 & Xử lý mất kết nối & Tự động kết nối lại; có phương án bổ sung buffer
cục bộ \\\hline
6 & Phát hiện chuyển động bất thường & Độ nhạy IMU cấu hình được \\\hline
\end{longtable}
}


\subsubsection{1.3.3. Tiêu chí máy chủ và dịch
vụ xử lý dữ liệu}\label{133-tiuxeau-chuxed-cloudbackend}

Khi chuyển sang tầng máy chủ và các dịch vụ xử lý dữ liệu, trọng tâm
đánh giá đổi sang khả năng tiếp nhận dữ liệu liên tục, phản hồi nhanh và
bảo vệ dữ liệu khi số lượng thiết bị tăng lên theo quy mô đội xe.


{\thesissettableid{1.2B}{tab-ch1-tieu-chi-cloud-backend}
\def\LTcaptype{table}
\begin{longtable}{|L{0.160\linewidth}|L{0.311\linewidth}|L{0.494\linewidth}|}
\caption{Tiêu chí máy chủ và dịch vụ xử lý dữ liệu}\label{tab:ch1-tieu-chi-cloud-backend}\\
\hline \textbf{STT} & \textbf{Tiêu chí} & \textbf{Mục tiêu cụ thể} \\
\\\hline
\endfirsthead
\caption[]{Tiêu chí máy chủ và dịch vụ xử lý dữ liệu (tiếp theo)}\\
\hline \textbf{STT} & \textbf{Tiêu chí} & \textbf{Mục tiêu cụ thể} \\
\\\hline
\endhead
\\\hline
\endfoot
\endlastfoot
1 & Độ trễ xử lý bản tin MQTT & \textless{} 500 ms từ thiết bị tới lớp
xử lý chính \\\hline
2 & Số lượng thiết bị đồng thời & \textgreater= 50 thiết bị (đã kiểm thử), dư địa tới khoảng 100 thiết bị \\\hline
3 & Độ sẵn sàng hệ thống (thời gian hoạt động ổn định) & \textgreater= 99.5\% \\\hline
4 & Thời gian phản hồi API & \textless{} 200 ms (p95) \\\hline
5 & Lưu trữ dữ liệu vận hành gửi liên tục như vị trí, tốc độ, trạng thái nguồn & \textgreater= 90 ngày \\\hline
6 & Bảo mật xác thực & Phiên đăng nhập lưu ở máy chủ; token ngẫu nhiên
chỉ dùng làm mã tham chiếu và được băm SHA-256 \\\hline
\end{longtable}
}


\subsubsection{1.3.4. Tiêu chí giao diện
web}\label{134-tiuxeau-chuxed-frontend}

Cuối cùng, giao diện web là nơi toàn bộ nỗ lực kỹ thuật được chuyển thành trải
nghiệm quan sát và ra quyết định cho người vận hành, nên tiêu chí ở đây tập
trung vào tính cập nhật, dễ theo dõi và dễ sử dụng trong bối cảnh thực tế.


{\thesissettableid{1.2C}{tab-ch1-tieu-chi-frontend}
\def\LTcaptype{table}
\begin{longtable}{|L{0.160\linewidth}|L{0.324\linewidth}|L{0.481\linewidth}|}
\caption{Tiêu chí giao diện web}\label{tab:ch1-tieu-chi-frontend}\\
\hline \textbf{STT} & \textbf{Tiêu chí} & \textbf{Mục tiêu cụ thể} \\
\\\hline
\endfirsthead
\caption[]{Tiêu chí giao diện web (tiếp theo)}\\
\hline \textbf{STT} & \textbf{Tiêu chí} & \textbf{Mục tiêu cụ thể} \\
\\\hline
\endhead
\\\hline
\endfoot
\endlastfoot
1 & Cập nhật bản đồ thời gian thực & \textless{} 2 giây độ trễ \\\hline
2 & Hỗ trợ trình duyệt & Chrome, Firefox, Edge (phiên bản mới nhất) \\\hline
3 & Bố cục thích ứng theo màn hình & Desktop và tablet \\\hline
4 & Hiển thị cảnh báo & Thông báo tức thời qua WebSocket, tức kênh để
máy chủ đẩy dữ liệu mới lên giao diện \\\hline
\end{longtable}
}

Các tiêu chí trên vì thế đóng vai trò như một thước đo chung cho toàn bộ
đồ án: mọi lựa chọn về phần cứng, firmware, máy chủ hay giao diện ở các
chương sau đều phải quay lại đối chiếu với bộ mục tiêu này. Từ nền đó,
phần tiếp theo trình bày cách nhóm tiếp cận bài toán theo từng bước.


\begin{center}\rule{0.5\linewidth}{0.5pt}\end{center}

\subsection{1.4. Phương pháp tiếp cận thiết kế kỹ
thuật}\label{14-phux1b0ux1a1ng-phuxe1p-tiux1ebfp-cux1eadn-thiux1ebft-kux1ebf-kux1ef9-thuux1eadt}

\subsubsection{1.4.1. Phương pháp luận tổng
thể}\label{141-phux1b0ux1a1ng-phuxe1p-luux1eadn-tux1ed5ng-thux1ec3}

Vì đề tài đi qua cả mạch nguồn, phần mềm nhúng trên thiết bị, hạ tầng
máy chủ và màn hình quản trị tổng hợp, rủi ro lớn nhất không nằm ở việc từng khối có chạy
riêng lẻ hay không, mà ở việc chúng có còn phối hợp đúng khi ghép thành
một hệ thống hoàn chỉnh hay không.

Do đó, dự án chọn cách làm theo hai nguyên tắc. Thứ nhất là \textbf{thiết
kế từ dưới lên (bottom-up design)}: kiểm chứng trước các giả định rủi ro
nhất ở tầng thấp rồi mới ghép dần lên các tầng trên. Thứ hai là
\textbf{phát triển lặp (iterative development)}: làm từng vòng nhỏ, đo
được và sửa được, thay vì để đến cuối dự án mới phát hiện xung đột về
nguồn, giao tiếp hay hiệu năng. Cụ thể, tiến trình được tổ chức như sau:

\begin{enumerate}
\def\labelenumi{\arabic{enumi}.}
\tightlist
\item
  \textbf{Giai đoạn 1 --- Nghiên cứu và thiết kế phần cứng}: Xây nền cho
  hệ thống bằng cách khảo sát linh kiện, đối chiếu datasheet, thiết kế
  sơ đồ nguyên lý và tính toán nguồn. Mục tiêu của giai đoạn này là bảo
  đảm thiết bị có thể tồn tại bền bỉ trong môi trường trên xe trước khi
  nói đến các tính năng ở tầng trên.
\item
  \textbf{Giai đoạn 2 --- Phát triển firmware}: Biến sơ đồ phần cứng
  thành một thiết bị có thể điều phối hành vi theo đúng trạng thái vận
  hành. Ở giai đoạn này, ESP32-S3 được lập trình bằng ESP-IDF để đọc dữ
  liệu xe qua BLE OBD2, giao tiếp với modem qua UART, nhận tín hiệu
  chuyển động từ IMU qua I2C, quyết định khi nào cần thức hoặc ngủ, rồi
  gửi dữ liệu lên máy chủ bằng MQTT. Đây là giai đoạn quyết định thiết
  bị theo dõi có phản ứng đúng lúc và tiết kiệm điện đúng chỗ hay không.
\item
  \textbf{Giai đoạn 3 --- Xây dựng hạ tầng máy chủ}: Dựng ``nơi dữ liệu
  sẽ đi tới'' sau khi rời chiếc xe. Giai đoạn này chuẩn bị các dịch vụ
  nhận bản tin, lưu dữ liệu quản trị, lưu dữ liệu thay đổi liên tục theo
  thời gian và lưu nhật ký vận hành, đồng thời chốt cách các dịch vụ kết
  nối với nhau.
\item
  \textbf{Giai đoạn 4 --- Phát triển API máy chủ}: Biến dữ liệu kỹ thuật
  thành những chức năng quản trị có thể dùng được, như đăng nhập, quản
  lý xe, xem lịch sử, tạo cảnh báo và phát cập nhật thời gian thực ra
  giao diện quản trị.
\item
  \textbf{Giai đoạn 5 --- Phát triển giao diện web}: Trình bày dữ liệu dưới
  dạng bản đồ, bảng thông tin và biểu đồ để người vận hành nhìn nhanh,
  hiểu nhanh và thao tác ít nhất có thể.
\item
  \textbf{Giai đoạn 6 --- Tích hợp và kiểm thử}: Ghép toàn bộ các lớp
  lại thành một chuỗi hoàn chỉnh từ xe đến giao diện quản trị, sau đó kiểm thử,
  đo đạc và tối ưu những điểm còn nghẽn.
\end{enumerate}

{\thesissetfigureid{1.2}{fig-ch1-project-development-phases}
\begin{figure}[htbp]
\centering
\pandocbounded{\safeincludesvg{./assets/figures/01-chuong-1-gioi-thieu-hinh-1-2.svg}}
\caption{Quy trình phát triển dự án theo các giai đoạn}
\label{fig:ch1-project-development-phases}
\end{figure}
}

\subsubsection{1.4.2. Kiến trúc hệ
thống}\label{142-kiux1ebfn-truxfac-hux1ec7-thux1ed1ng}

Nếu nhìn theo đường đi của thông tin từ xe đến màn hình quản lý, hệ
thống có thể hình dung qua ba lớp. Lớp thứ nhất nằm trên xe, có
nhiệm vụ thu dữ liệu và gửi dữ liệu đi. Lớp thứ hai nằm trên hạ tầng
hệ thống máy chủ, có nhiệm vụ tiếp nhận, phân luồng và lưu dữ liệu vào
đúng nơi.
Lớp thứ ba là lớp ứng dụng, nơi dữ liệu được biến thành bản đồ, bảng
thông tin, cảnh báo và công cụ quản trị để con người sử dụng. Cách nhìn
này giải thích vì sao kiến trúc được tổ chức theo dạng phân tầng thay vì
để mọi chức năng dồn vào một chỗ.

Trong triển khai thực tế, lớp trên xe gồm các thành phần chuyên đi lấy
dữ liệu và giữ thiết bị hoạt động ổn định, tiêu biểu là ESP32-S3, modem
SIM7600CE-T, cảm biến rung LIS3DSH và adapter OBD2. Lớp giữa làm nhiệm
vụ nhận bản tin, kiểm tra rồi chuyển dữ liệu sang nơi lưu phù hợp như
VictoriaMetrics, VictoriaLogs và API máy chủ thông qua EMQX và MQTT
Bridge. Lớp trên cùng là nơi người quản lý làm việc, gồm giao diện web
quản trị và công cụ giám sát. Ứng dụng di động được giữ cho giai đoạn
sau để đồ án tập trung trước vào chuỗi giá trị cốt lõi từ xe đến màn
hình quản trị.

{\thesissetfigureid{1.3}{fig-ch1-system-architecture-overview}
\begin{figure}[htbp]
\centering
\pandocbounded{\safeincludesvg{./assets/figures/01-chuong-1-gioi-thieu-hinh-1-3.svg}}
\caption{Kiến trúc tổng thể hệ thống IoT Vehicle Tracking}
\label{fig:ch1-system-architecture-overview}
\end{figure}
}

\textbf{Chiến lược dữ liệu (Data Strategy):}

Chiến lược dữ liệu được xây theo cùng logic phân tầng này vì dữ liệu
trong hệ thống không đồng nhất: có loại cần tính nhất quán và quan hệ
chặt chẽ, có loại cần ghi liên tục với tần suất cao, và có loại chủ yếu
phục vụ truy vết sự cố. Mỗi nhóm dữ liệu vì thế được đưa về đúng nơi phù
hợp với đặc tính truy vấn và mục đích sử dụng.

\begin{itemize}
\tightlist
\item
  \textbf{PostgreSQL}: Lưu trữ dữ liệu quan hệ có cấu trúc như người
  dùng, xe, khách hàng, hành trình, cảnh báo, vùng giám sát trên bản đồ
  và lệnh điều khiển. Đây là lớp dữ liệu cần quan hệ rõ ràng và ràng
  buộc chặt chẽ.
\item
  \textbf{VictoriaMetrics}: Lưu trữ dữ liệu chuỗi thời gian như tọa độ
  GPS, dữ liệu OBD2 và chỉ số cảm biến. Nhóm dữ liệu này cần ghi nhiều,
  đọc nhanh theo trục thời gian và lưu dài hạn với chi phí hợp lý.
\item
  \textbf{VictoriaLogs}: Lưu trữ nhật ký sự kiện và dấu
  vết ai làm gì và khi nào, cùng các bản
  ghi hoạt động của thiết bị để phục vụ quan sát hệ thống và truy vết sự
  cố.
\end{itemize}

\subsubsection{1.4.3. Chiến lược quản lý năng
lượng}\label{143-chiux1ebfn-lux1b0ux1ee3c-quux1ea3n-luxfd-nux103ng-lux1b0ux1ee3ng}

Thiết kế năng lượng của thiết bị theo dõi thực chất là bài toán quyết định lúc nào
thiết bị cần hoạt động đầy đủ, lúc nào chỉ nên ``canh chừng'', và lúc nào
phải ưu tiên bảo vệ nguồn xe. Nếu xử lý không khéo, hệ thống sẽ rơi vào
hai cực đoan: theo dõi rất sát nhưng làm hao ắc quy, hoặc tiết kiệm điện
quá mức nên bỏ lỡ sự kiện quan trọng. Vì vậy, chiến lược quản lý năng
lượng đa chế độ (multi-mode power management) được xây quanh các trạng
thái vận hành thật của xe, để mỗi thời điểm chỉ bật những khối thật sự
cần thiết:

\textbf{Chế độ 1 --- Khi xe đang chạy (IGN ON, ưu tiên độ đầy đủ dữ liệu):}

\begin{itemize}
\tightlist
\item
  Hệ thống hoạt động toàn bộ công suất
\item
  Gửi vị trí GPS định kỳ (5--30 giây)
\item
  Sạc pin dự phòng từ nguồn xe
\item
  Kết nối máy chủ liên tục để theo dõi thời gian thực
\end{itemize}

\textbf{Chế độ 2 --- Parking Mode (IGN OFF, ưu tiên bảo toàn năng lượng):}

\begin{itemize}
\tightlist
\item
  Vi điều khiển ESP32-S3 vào chế độ ngủ sâu (deep sleep)
\item
  Tiêu thụ cực thấp (\textless{} 500 µA)
\item
  Thức dậy định kỳ (10--30 phút) để gửi heartbeat
\item
  IMU (LIS3DSH) hoạt động độc lập, cảnh rung ở ngưỡng đã cấu hình
\end{itemize}

\textbf{Chế độ 3 --- Cảnh báo khẩn (ưu tiên phản ứng nhanh khi có sự cố):}

\begin{itemize}
\tightlist
\item
  IMU đánh thức ESP32 ngay lập tức qua interrupt
\item
  Bật 4G + GPS, gửi cảnh báo ưu tiên lên máy chủ
\item
  Tiếp tục theo dõi liên tục cho đến khi được xác nhận an toàn
\item
  Sử dụng pin dự phòng nếu nguồn chính bị cắt
\end{itemize}

Bên cạnh đó, hệ thống còn có mạch chuyển nguồn dự phòng với cơ chế
hysteresis để tránh đóng/ngắt liên tục khi điện áp dao động quanh ngưỡng.
Thiết bị tự động chuyển sang pin dự phòng khi điện áp ắc quy tụt xuống
dưới mức an toàn theo từng profile: 12V dùng
\(Switch_{\mathrm{OFF}} = 12.0\,\mathrm{V}\), \(Switch_{\mathrm{ON}} = 12.2\,\mathrm{V}\);
24V dùng \(Switch_{\mathrm{OFF}} = 24.0\,\mathrm{V}\),
\(Switch_{\mathrm{ON}} = 24.4\,\mathrm{V}\), qua đó bảo vệ ắc quy không bị rút
cạn quá mức {[}7{]}.

Nhờ cách tổ chức này, năng lượng không còn là bài toán tối ưu từng linh
kiện riêng lẻ, mà trở thành một phần của logic vận hành toàn hệ thống.

{\thesissetfigureid{1.4}{fig-ch1-power-mode-transitions}
\begin{figure}[htbp]
\centering
\pandocbounded{\safeincludesvg{./assets/figures/01-chuong-1-gioi-thieu-hinh-1-4.svg}}
\caption{Sơ đồ chuyển đổi giữa các chế độ năng lượng}
\label{fig:ch1-power-mode-transitions}
\end{figure}
}

\subsubsection{1.4.4. Công nghệ sử
dụng}\label{144-cuxf4ng-nghux1ec7-sux1eed-dux1ee5ng}

Phần công nghệ dưới đây không nhằm chứng minh dự án dùng nhiều công cụ,
mà để chỉ ra mỗi công cụ giải quyết đúng một việc trong chuỗi vận hành.
Nếu bỏ cách nhìn ``vai trò trước, tên công nghệ sau'', bảng này rất dễ
trở thành danh sách thuật ngữ rời rạc.

Có thể đọc toàn bộ hệ thống theo ba cụm vai trò thay vì học thuộc từng tên
công nghệ. Cụm thứ nhất là \textbf{khối trên xe}: vi điều khiển, modem,
cảm biến và adapter OBD2 chịu trách nhiệm thu dữ liệu và giữ thiết bị
hoạt động ổn định. Cụm thứ hai là \textbf{khối tiếp nhận và lưu trữ}:
firmware điều phối phần cứng, còn hạ tầng máy chủ nhận bản tin, phân
loại, lưu dữ liệu và lưu dấu vết vận hành. Cụm cuối là \textbf{khối khai
thác}: giao diện web, bản đồ, biểu đồ và công cụ giám sát giúp người vận
hành đọc hiểu dữ liệu nhanh. Với cách nhìn này, các tên như ESP32-S3,
EMQX hay Next.js chỉ còn là ví dụ cụ thể cho từng vai trò, không phải
một danh sách thuật ngữ rời rạc.

{\thesissettableid{1.3}{tab-ch1-technology-stack-summary}
\def\LTcaptype{table}
\begin{longtable}{|L{0.223\linewidth}|L{0.334\linewidth}|L{0.115\linewidth}|L{0.282\linewidth}|}
\caption{Tổng hợp công nghệ sử dụng trong dự án}\label{tab:ch1-technology-stack-summary}\\
\hline \textbf{Thành phần} & \textbf{Công nghệ} & \textbf{Phiên bản} & \textbf{Vai trò} \\
\\\hline
\endfirsthead
\caption[]{Tổng hợp công nghệ sử dụng trong dự án (tiếp theo)}\\
\hline \textbf{Thành phần} & \textbf{Công nghệ} & \textbf{Phiên bản} & \textbf{Vai trò} \\
\\\hline
\endhead
\\\hline
\endfoot
\endlastfoot
Vi điều khiển & ESP32-S3 & --- & Vi điều khiển chính, điều phối toàn bộ
thiết bị \\\hline
LTE + GNSS & SIMCom SIM7600CE-T (LTE + GNSS tích hợp) & --- & Truyền dữ
liệu 4G và lấy vị trí từ vệ tinh \\\hline
OBD2 Adapter & vgate iCar Pro & BLE 4.0 & Đọc dữ liệu chẩn đoán xe \\\hline
Cảm biến gia tốc & LIS3DSH & --- & Phát hiện chuyển động và rung \\\hline
Pin dự phòng & 18650 1S Li-ion & 3500 mAh & Nguồn điện dự phòng \\\hline
Framework firmware & ESP-IDF & 5.x & Bộ công cụ để phát triển phần mềm
nhúng cho ESP32 \\\hline
RTOS & FreeRTOS & --- & Lớp điều phối các tác vụ nền theo thời gian thực \\\hline
MQTT Broker & EMQX & 5.x & Trạm trung chuyển bản tin từ thiết bị \\\hline
API máy chủ & Express.js + TypeScript & --- & Lớp giao tiếp để
giao diện gọi lấy dữ liệu và gửi lệnh \\\hline
Cơ sở dữ liệu quan hệ & PostgreSQL & 16 & Lưu trữ dữ liệu có cấu trúc \\\hline
Cơ sở dữ liệu chuỗi thời gian & VictoriaMetrics & --- & Lưu trữ
dữ liệu vận hành gửi liên tục như vị trí, trạng thái nguồn và cảm biến \\\hline
Nhật ký sự kiện & VictoriaLogs & --- & Lưu trữ nhật ký vận hành và lỗi \\\hline
Giao diện web & Next.js 15 + React 19 & --- & Giao diện quản lý cho
người vận hành \\\hline
Bản đồ & Leaflet & --- & Hiển thị bản đồ thời gian thực \\\hline
Biểu đồ & ECharts & --- & Trực quan hóa dữ liệu \\\hline
WebSocket & Socket.IO & 4.8 & Kênh để máy chủ đẩy dữ liệu mới lên giao
diện \\\hline
Container hóa & Docker + Docker Compose & --- & Triển khai dịch vụ \\\hline
Giám sát & Grafana + VictoriaMetrics, VictoriaLogs & --- & Màn hình giám
sát hạ tầng và nhật ký vận hành \\\hline
\end{longtable}
}

\begin{center}\rule{0.5\linewidth}{0.5pt}\end{center}

\subsection{1.5. Kết quả và khuyến
nghị}\label{15-kux1ebft-quux1ea3-vuxe0-khuyux1ebfn-nghux1ecb}

\subsubsection{1.5.1. Kết quả đạt
được}\label{151-kux1ebft-quux1ea3-ux111ux1ea1t-ux111ux1b0ux1ee3c}

Nhìn theo cách đó, các kết quả đạt được không nên hiểu như một danh sách
các khối đã lắp xong, mà như bằng chứng cho thấy chuỗi giá trị từ thiết
bị đến giao diện quản trị đã được hình thành. Có thể tóm lược theo bốn tầng chính
nhÆ° sau:

\textbf{Về phần cứng:}

\begin{itemize}
\tightlist
\item
  Hoàn thiện thiết bị theo dõi IoT với bo mạch riêng do nhóm thiết kế, sử
  dụng ESP32-S3-WROOM-1 làm vi điều khiển trung tâm, tích hợp modem LTE
  + GNSS SIMCom SIM7600CE-T, cảm biến gia tốc LIS3DSH, khối nguồn đa
  rail, pin dự phòng 18650 1S và kết nối OBD2 BLE qua adapter vgate iCar
  Pro.
\item
  Quá trình lựa chọn linh kiện bám theo datasheet, độ ổn định chất lượng
  và khả năng cung ứng trong nước để thuận lợi cho việc chế tạo, kiểm
  thử và thay thế khi cần.
\item
  Hệ thống quản lý năng lượng đa chế độ hoạt động hiệu quả, với mức tiêu
  thụ điện ngủ sâu đạt yêu cầu (\textless{} 500 µA), đảm bảo không làm
  cạn ắc quy xe trong quá trình sử dụng bình thường.
\item
  Mạch Low Voltage Disconnect (LVD) bảo vệ ắc quy xe hiệu quả, tự động
  ngắt khi điện áp tụt dưới ngưỡng an toàn.
\end{itemize}

\textbf{Về firmware:}

\begin{itemize}
\tightlist
\item
  Phát triển thành công phần mềm trên thiết bị bằng ESP-IDF kết hợp
  FreeRTOS, giúp thiết bị theo dõi đọc dữ liệu xe qua BLE OBD2, điều khiển modem
  qua UART, theo dõi chuyển động bằng IMU, quản lý năng lượng và gửi dữ
  liệu MQTT lên hệ thống máy chủ.
\item
  Cơ chế phát hiện chuyển động bất thường qua IMU hoạt động chính xác,
  có khả năng đánh thức hệ thống từ chế độ ngủ sâu trong vòng
  \textless{} 3 giây.
\end{itemize}

\textbf{Về hạ tầng máy chủ và xử lý dữ liệu:}

\begin{itemize}
\tightlist
\item
  Xây dựng thành công hạ tầng máy chủ hoàn chỉnh để nhận dữ liệu từ
  thiết bị, lưu dữ liệu nghiệp vụ, lưu dữ liệu theo thời gian và lưu
  nhật ký vận hành; toàn bộ được đóng gói bằng Docker để dễ triển khai
  và lặp lại môi trường.
\item
  Hoàn thiện API quản trị phía máy chủ cho các nhu cầu cốt lõi như đăng
  nhập, quản lý xe, quản lý thiết bị, xem dữ liệu vận hành, xử lý cảnh
  báo và quản lý vùng giám sát.
\item
  Hệ thống đã kiểm thử ở mức 50+ thiết bị đồng thời với độ trễ từ lúc dữ
  liệu rời thiết bị đến lúc tới nơi xử lý dưới 500 ms, đồng thời cho
  thấy còn dư địa mở rộng ổn định tới khoảng 100 thiết bị.
\end{itemize}

\textbf{Về giao diện web:}

\begin{itemize}
\tightlist
\item
  Giao diện web quản lý (Next.js 15) với bản đồ thời gian thực
  (Leaflet), biểu đồ phân tích (ECharts), hệ thống cảnh báo trực tuyến
  và màn hình tổng quan.
\item
  Cập nhật vị trí xe trên bản đồ với độ trễ \textless{} 2 giây thông qua
  WebSocket (Socket.IO), tức kênh để máy chủ chủ động đẩy dữ liệu mới lên
  trình duyệt.
\end{itemize}

\subsubsection{1.5.2. Khuyến nghị và hướng phát
triển}\label{152-khuyux1ebfn-nghux1ecb-vuxe0-hux1b0ux1edbng-phuxe1t-triux1ec3n}

Trên cơ sở các kết quả đã đạt được và các giới hạn hiện hữu, các hướng phát
triển tiếp theo được sắp theo logic mở rộng từ việc làm chắc hệ lõi đến
việc mở rộng giá trị sử dụng của hệ thống:

\begin{enumerate}
\def\labelenumi{\arabic{enumi}.}
\tightlist
\item
  \textbf{Ứng dụng di động (giai đoạn 2)}: Phát triển ứng dụng Flutter
  theo hướng vỏ ứng dụng mỏng, chủ yếu nhúng lại giao diện web hiện có,
  để người quản lý giám sát đội xe trên điện thoại và nhận thông báo đẩy
  khi có cảnh báo {[}8{]}.
\item
  \textbf{Tăng cường bảo mật}: Triển khai TLS/SSL cho kết nối MQTT trong
  môi trường vận hành thật, áp dụng danh sách quyền MQTT chi tiết hơn để
  mỗi thiết bị chỉ đi đúng kênh dữ liệu của mình, và bổ sung cơ chế mã
  hóa dữ liệu từ thiết bị tới máy chủ.
\item
  \textbf{Tối ưu hóa hiệu năng}: Nghiên cứu và áp dụng các kỹ thuật nén
  dữ liệu vận hành gửi liên tục, chẳng hạn chỉ gửi phần sai khác hoặc
  dùng định dạng gói gọn hơn, để giảm băng thông 4G,
  đặc biệt khi vận hành đội xe lớn (\textgreater{} 500 xe).
\item
  \textbf{Tích hợp trí tuệ nhân tạo}: Áp dụng các mô hình học máy để
  phân tích hành vi lái xe, dự đoán thời điểm cần bảo trì và phát hiện
  các dấu hiệu bất thường từ dữ liệu vận hành đã thu thập.
\item
  \textbf{Mở rộng phạm vi OBD2}: Hỗ trợ đọc thêm nhiều thông số chẩn
  đoán xe, bao gồm mã lỗi DTC (Diagnostic Trouble Codes), dữ liệu động
  cơ nâng cao, và tích hợp với các loại xe điện (EV).
\item
  \textbf{Cải thiện khả năng mở rộng}: Nghiên cứu kiến
  trúc tách nhiều dịch vụ nhỏ, có thêm hàng đợi bản tin chuyên dụng
  (chẳng hạn Kafka hoặc RabbitMQ), để xử lý luồng dữ liệu lớn hơn khi số
  lượng thiết bị tăng lên hàng ngàn.
\item
  \textbf{Tối ưu PCB cho sản xuất}: Nâng thiết kế PCB hiện tại theo
  hướng dễ sản xuất hơn và ít nhiễu điện từ hơn, giảm kích thước, tăng
  độ bền cơ khí và chuẩn bị cho sản xuất hàng loạt.
\end{enumerate}

{\thesissetfigureid{1.5}{fig-ch1-project-roadmap}
\begin{figure}[htbp]
\centering
\pandocbounded{\safeincludesvg{./assets/figures/01-chuong-1-gioi-thieu-hinh-1-5.svg}}
\caption{Lộ trình phát triển dự án theo các giai đoạn}
\label{fig:ch1-project-roadmap}
\end{figure}
}

\begin{center}\rule{0.5\linewidth}{0.5pt}\end{center}

\chapterheading{CHƯƠNG 2. PHÂN TÍCH VẤN ĐỀ KỸ THUẬT}

Nếu Chương 1 đặt ra bối cảnh và mục tiêu của đề tài, thì Chương 2 đi
sâu vào phần cốt lõi hơn: bóc tách bài toán kỹ thuật thật sự của hệ
thống. Trọng tâm của chương không chỉ là liệt kê công nghệ, mà là làm rõ
vì sao mỗi quyết định về phần cứng, truyền thông, lưu trữ và kiến trúc
đều phải xuất phát từ nhu cầu vận hành thực tế của một hệ thống giám sát
phương tiện.

\subsection{2.1. Mô tả vấn đề -- Problem
statement}\label{21-muxf4-tux1ea3-vux1ea5n-ux111ux1ec1--problem-statement}

\subsubsection{2.1.1. Bối cảnh thực
tế}\label{211-bux1ed1i-cux1ea3nh-thux1ef1c-tux1ebf}

Trong bối cảnh ngành cho thuê ô tô và quản lý đội xe tại Việt Nam tiếp
tục mở rộng, nhu cầu giám sát phương tiện theo thời gian thực đã trở
thành yêu cầu thiết yếu đối với doanh nghiệp vận tải. Theo thống kê của
Bộ Giao thông Vận tải, số lượng ô tô cá nhân và thương mại tăng trung
bình 10--12\% mỗi năm trong giai đoạn 2020--2025, kéo theo áp lực ngày
càng lớn về quản lý, an ninh và tối ưu vận hành {[}1{]}.

Tuy nhiên, điều doanh nghiệp cần không chỉ là các điểm vị trí rời rạc,
mà là khả năng nhìn thấy trạng thái phương tiện đủ sớm để đưa ra quyết
định vận hành. Khoảng cách giữa nhu cầu này và cách làm hiện tại thể
hiện ở các hạn chế sau:

\begin{itemize}
\tightlist
\item
  \textbf{Giám sát thủ công}: Tài xế báo cáo vị trí qua điện thoại,
  không có dữ liệu liên tục, không thể xác minh chính xác hành trình.
  Phương pháp này phụ thuộc hoàn toàn vào yếu tố con người, dễ xảy ra
  sai sót và gian lận.
\item
  \textbf{Thiết bị GPS đơn giản}: Chỉ cung cấp tọa độ vị trí, không đọc
  được dữ liệu động cơ (tốc độ, vòng tua, nhiên liệu), không hỗ trợ phát
  hiện bất thường khi xe đậu. Thiếu khả năng tích hợp sâu với hệ thống
  quản lý.
\item
  \textbf{Giải pháp thương mại (fleet management)}: Chi phí cao (50--200
  USD/thiết bị + phí dịch vụ hàng tháng), phụ thuộc vào nhà cung cấp
  nước ngoài, khó tùy biến theo nhu cầu cụ thể của doanh nghiệp Việt
  Nam.
\end{itemize}

\subsubsection{2.1.2. Xác định vấn đề kỹ
thuật}\label{212-xuxe1c-ux111ux1ecbnh-vux1ea5n-ux111ux1ec1-kux1ef9-thuux1eadt}

Từ bối cảnh thực tiễn đó, bài toán của đề tài có thể quy về bốn câu hỏi
kỹ thuật nền tảng. Đây cũng là bốn câu hỏi chi phối gần như mọi quyết
định ở các chương sau, từ lựa chọn phần cứng đến thiết kế hạ tầng máy
chủ:

\textbf{Vấn đề 1: Quản lý năng lượng trong môi trường ô tô}

Thiết bị theo dõi phải hoạt động liên tục 24/7 trong môi trường ô tô,
nơi nguồn điện luôn bị ràng buộc bởi dung lượng ắc quy. Khi xe tắt máy
(IGN OFF), thiết bị vẫn lấy điện từ ắc quy 12V hoặc 24V của xe; nếu duy
trì chế độ theo dõi liên tục, ắc quy có thể cạn sau vài tuần và làm xe
không khởi động được {[}2{]}. Vì vậy, câu hỏi không đơn thuần là "thiết
bị có hoạt động được không", mà là "thiết bị có hoạt động đủ thông minh
để không gây phiền toái cho chính chiếc xe mà nó đang giám sát hay
không".

\textbf{Vấn đề 2: Giám sát khi xe đậu}

Khi xe tắt máy, hệ thống cần tiết kiệm điện nhưng vẫn phát hiện được các
chuyển động bất thường như bị trộm hoặc bị kéo cẩu. Vì vậy, thiết bị
phải có cơ chế đánh thức nhanh khi phát hiện sự kiện, đồng thời duy trì
dòng tiêu thụ ở mức rất thấp. Khi xe không hoạt động, thiết bị phải
chuyển sang trạng thái tiêu thụ điện thấp nhưng vẫn duy trì khả năng
phát hiện sự cố.

\textbf{Vấn đề 3: Truyền dữ liệu qua mạng di động}

Hệ thống phải gửi đều dữ liệu vị trí và trạng thái xe qua mạng di động,
ngay cả khi chất lượng sóng thay đổi liên tục. Vì vậy, bài toán không
chỉ là ``gửi được hay không'', mà là khi mất mạng thiết bị có tự nối lại
được không, có giữ được dữ liệu quan trọng không, và có mở rộng được khi
số xe tăng lên không. Vì thế, việc chọn giữa MQTT, HTTP hay CoAP chính
là chọn ``cách gửi tin'' phù hợp nhất với một thiết bị phải truyền nhiều
gói nhỏ, đều đặn và đôi khi cần nhận lệnh ngược từ máy chủ. MQTT
phù hợp khi thiết bị cần gửi các gói dữ liệu nhỏ thường xuyên; HTTP gần
với mô hình ``gửi yêu cầu rồi chờ phản hồi'' quen thuộc của web; còn
CoAP là giao thức gọn hơn, thường dùng cho thiết bị có tài nguyên rất
hạn chế.

\textbf{Vấn đề 4: Kiến trúc máy chủ có khả năng mở rộng}

Hệ thống cần hỗ trợ nhiều phương tiện đồng thời, cập nhật vị trí theo
thời gian thực qua WebSocket, bảo đảm an toàn bằng cơ chế xác thực
phiên, và xử lý được luồng dữ liệu lớn đến từ nhiều thiết bị IoT. Với
người sử dụng cuối, điều này được cảm nhận rất đơn giản: bản đồ phải cập
nhật kịp thời, cảnh báo phải đến đúng lúc, và dữ liệu phải nhất quán
ngay cả khi số lượng xe tăng lên.

\subsubsection{2.1.3. Mục tiêu kỹ
thuật}\label{213-mux1ee5c-tiuxeau-kux1ef9-thuux1eadt}

Từ bốn vấn đề nền tảng trên, đồ án rút ra bốn nhóm mục tiêu kỹ thuật.
Mỗi nhóm vừa trả lời một nhu cầu vận hành cụ thể của doanh nghiệp cho
thuê xe, vừa trở thành tiêu chí để sàng lọc và lựa chọn giải pháp ở
ChÆ°Æ¡ng 3:

\begin{enumerate}
\def\labelenumi{\arabic{enumi}.}
\tightlist
\item
  Thiết kế thiết bị theo dõi có thể vừa xác định vị trí bằng GNSS, vừa
  đọc dữ liệu chẩn đoán từ xe qua OBD2, dựa trên vi điều khiển ESP32-S3;
  đồng thời tối ưu tiêu thụ năng lượng khi xe đậu (mục tiêu dòng deep
  sleep toàn thiết bị dưới 500 µA trong phiên bản phần cứng hiện tại và
  tiếp tục giảm ở các vòng tối ưu PCB tiếp theo).
\item
  Xây dựng cơ chế quản lý nguồn thông minh với pin dự phòng và mạch Low
  Voltage Disconnect (LVD), tức mạch tự ngắt khi điện áp xuống thấp, để
  thiết bị bám được trạng thái xe nhưng không làm ảnh hưởng đến khả năng
  khởi động của phương tiện.
\item
  Thiết kế đường truyền dữ liệu sử dụng MQTT qua mạng 4G/LTE để gửi các
  gói dữ liệu nhỏ qua mạng di động, ưu tiên tự phục hồi kết nối và chừa sẵn
  phương án lưu đệm khi mất mạng kéo dài.
\item
  Phát triển nền tảng máy chủ có khả năng mở rộng, bao gồm API máy chủ,
  cơ sở dữ liệu quan hệ cho phần quản trị, cơ sở dữ liệu chuỗi thời gian
  cho dữ liệu phát sinh liên tục, và giao diện web theo dõi thời gian
  thá»±c.
\end{enumerate}

Bốn nhóm mục tiêu trên cho biết đích đến của thiết kế. Tuy nhiên, để
chọn đúng công nghệ và kiến trúc, vẫn cần đặt chúng vào bối cảnh kỹ
thuật rộng hơn. Vì vậy, mục 2.2 chuyển sang rà lại các nền tảng kỹ thuật
liên quan trực tiếp đến bài toán.

\subsection{2.2. Bối cảnh và cơ sở kỹ thuật -- Background and Technical
reviews}\label{22-bux1ed1i-cux1ea3nh-vuxe0-cux1a1-sux1edf-kux1ef9-thuux1eadt--background-and-technical-reviews}

\subsubsection{2.2.1. Tổng quan về hệ thống IoT trong giám sát phương
tiện}\label{221-tux1ed5ng-quan-vux1ec1-hux1ec7-thux1ed1ng-iot-trong-giuxe1m-suxe1t-phux1b0ux1a1ng-tiux1ec7n}

Internet of Things (IoT) là mô hình cho phép các thiết bị vật lý kết nối
với nhau và với hệ thống máy chủ thông qua Internet. Ở góc nhìn đơn giản,
một hệ thống IoT giám sát phương tiện luôn phải đi qua ba bước: thu dữ
liệu trên xe, truyền dữ liệu lên mạng, rồi biến dữ liệu đó thành thông
tin mà người quản lý có thể đọc và hành động. Tương ứng với ba bước đó,
hệ thống thường được tổ chức thành ba lớp chính {[}3{]}:

\begin{itemize}
\tightlist
\item
  \textbf{Lớp cảm biến và thu thập dữ liệu (Perception Layer)}: Đây là
  phần nằm gần chiếc xe nhất. Nó gồm GPS/GNSS để biết xe đang ở đâu, IMU
  để biết xe có rung hoặc di chuyển không, OBD2 để đọc trạng thái vận
  hành của xe, và các cảm biến khác nếu cần. Nhiệm vụ của lớp này là
  ``nhìn thấy'' những gì đang diễn ra trên xe và biến chúng thành dữ
  liệu thô.
\item
  \textbf{Lớp mạng và truyền thông (Network Layer)}: Đây là con đường
  đưa dữ liệu rời khỏi chiếc xe. Lớp này dùng các giao thức như MQTT,
  HTTP hoặc CoAP qua mạng di động để chuyển dữ liệu từ thiết bị lên máy
  chủ. Nói gọn hơn, đây là phần quyết định dữ liệu sẽ được gửi theo kiểu
  ``gói dữ liệu nhỏ gửi liên tục'', ``yêu cầu rồi chờ trả lời'' hay một chuẩn
  gọn hơn cho thiết bị nhỏ. Nó cũng phải xử lý các tình huống mất sóng
  hoặc kết nối lại.
\item
  \textbf{Lớp ứng dụng và xử lý (Application Layer)}: Đây là tầng biến dữ
  liệu thành thông tin mà con người có thể khai thác. Lớp này bao gồm
  các dịch vụ máy chủ như broker MQTT, API máy chủ, cơ sở dữ liệu và
  giao diện quản trị như màn hình web tổng hợp hoặc ứng dụng di động.
\end{itemize}

Trong triển khai của đề tài, kiến trúc ba lớp trên được bổ sung thêm một
lớp đệm dữ liệu cục bộ tại thiết bị bằng \textbf{microSD SDMMC 4-bit +
FATFS}. Đây là một thẻ nhớ cục bộ nối với ESP32-S3 bằng chuẩn giao tiếp
tốc độ cao và được quản lý bằng hệ thống tệp quen thuộc. Có thể xem nó
như một ``kho chờ ngắn hạn'': khi mạng di động gián đoạn, dữ liệu vận
hành gửi lên liên tục sẽ tạm nằm ở đây; khi kết nối phục hồi, hệ thống
phát lại theo đúng thứ tự. Nhờ vậy, độ liên tục dữ liệu tăng lên mà
không cần thay đổi lớn ở kiến trúc máy chủ phía trên.

\subsubsection{2.2.2. Các giải pháp giám sát phương tiện hiện
có}\label{222-cuxe1c-giux1ea3i-phuxe1p-giuxe1m-suxe1t-phux1b0ux1a1ng-tiux1ec7n-hiux1ec7n-cuxf3}

{\thesissettableid{2.1}{tab-ch2-vehicle-monitoring-solutions-comparison}
\def\LTcaptype{table}
\begin{longtable}{|L{0.179\linewidth}|L{0.141\linewidth}|L{0.167\linewidth}|L{0.227\linewidth}|L{0.231\linewidth}|}
\caption{So sánh các giải pháp giám sát phương tiện}\label{tab:ch2-vehicle-monitoring-solutions-comparison}\\
\hline \textbf{Tiêu chí} & \textbf{Giám sát thủ công} & \textbf{GPS Tracker đơn giản} & \textbf{Fleet Management
thương mại} & \textbf{Hệ thống đề xuất} \\
\\\hline
\endfirsthead
\caption[]{So sánh các giải pháp giám sát phương tiện (tiếp theo)}\\
\hline \textbf{Tiêu chí} & \textbf{Giám sát thủ công} & \textbf{GPS Tracker đơn giản} & \textbf{Fleet Management
thương mại} & \textbf{Hệ thống đề xuất} \\
\\\hline
\endhead
\\\hline
\endfoot
\endlastfoot
Chi phí thiết bị & Không & 500.000--1.500.000 VND & 1.000.000--5.000.000
VND & 1.514.000 VND \\\hline
Phí dịch vụ hàng tháng & Không & 50.000--100.000 VND & 200.000--500.000
VND & Chi phí 4G SIM (~70.000 VND) \\\hline
Độ chính xác vị trí & Thấp (báo cáo thủ công) & Trung bình (GPS) & Cao
(GPS + A-GPS) & Cao (GNSS đa hệ thống) \\\hline
Dữ liệu động cơ (OBD2) & Không & Không & Có (tùy model) & Có (BLE
OBD2) \\\hline
Phát hiện bất thường khi đậu & Không & Hạn chế & Có & Có (IMU + deep
sleep) \\\hline
Khả năng tùy biến & Không áp dụng & Thấp & Thấp (phụ thuộc nhà cung cấp) & Cao
(mã nguồn mở) \\\hline
Quản lý năng lượng & Không áp dụng & Cơ bản & Tốt & Tốt (đa chế độ + pin
dự phòng) \\\hline
Tích hợp hệ thống & Không & Hạn chế (API riêng) & API do nhà cung cấp & API mở
(REST + WebSocket) \\\hline
\end{longtable}
}

Bảng so sánh cho thấy phương án đề xuất giữ được các năng lực quan trọng
của hệ thống thương mại như độ chính xác vị trí, khả năng đọc OBD2 và
phát hiện bất thường, nhưng chi phí đầu tư thấp hơn và mức độ tùy biến
cao hơn. Với bài toán của doanh nghiệp cho thuê xe trong nước, đây là
hướng lựa chọn phù hợp hơn về cả kỹ thuật lẫn chi phí vận hành.

\subsubsection{2.2.3. So sánh giao thức truyền thông
IoT}\label{223-so-suxe1nh-giao-thux1ee9c-truyux1ec1n-thuxf4ng-iot}

Việc lựa chọn giao thức truyền thông không thể chỉ hiểu là chọn một
``chuẩn phổ biến''. Với người không chuyên, có thể hình dung ba lựa chọn
chính như sau: HTTP giống cách ``gửi yêu cầu rồi chờ trả lời''; MQTT
giống cách ``đưa bản tin vào một kênh chung để bên cần nhận lấy ra'';
còn CoAP là giao thức nhẹ cho thiết bị tài nguyên rất hạn chế,
thường đi trên UDP.

Với hệ thống theo dõi phương tiện, quyết định này chi phối trực tiếp
lượng dữ liệu phải truyền, khả năng chịu mất sóng và cách dịch vụ máy
chủ mở rộng về sau {[}4{]}. Vì thế, việc so sánh không nhằm tìm giao
thức ``tốt nhất nói chung'', mà tìm giao thức phù hợp nhất với nhịp gửi
dữ liệu nhỏ, liên tục và có lúc cần nhận lệnh ngược từ máy chủ.

Trong bảng dưới đây, ``overhead'' là phần dữ liệu phụ bắt buộc đi kèm
mỗi gói tin; còn ``QoS'' là mức cam kết giao bản tin thành công. Hai
tiêu chí này quan trọng vì chúng ảnh hưởng trực tiếp đến băng thông, độ
tin cậy và tuổi thọ pin của thiết bị.

{\thesissettableid{2.2}{tab-ch2-iot-communication-protocols}
\def\LTcaptype{table}
\begin{longtable}{|L{0.279\linewidth}|L{0.206\linewidth}|L{0.232\linewidth}|L{0.238\linewidth}|}
\caption{So sánh các giao thức truyền thông IoT}\label{tab:ch2-iot-communication-protocols}\\
\hline \textbf{Tiêu chí} & \textbf{MQTT} & \textbf{HTTP / REST} & \textbf{CoAP} \\
\\\hline
\endfirsthead
\caption[]{So sánh các giao thức truyền thông IoT (tiếp theo)}\\
\hline \textbf{Tiêu chí} & \textbf{MQTT} & \textbf{HTTP / REST} & \textbf{CoAP} \\
\\\hline
\endhead
\\\hline
\endfoot
\endlastfoot
Mô hình truyền thông & Phát bản tin vào kênh chung, bên nhận đăng ký
theo dõi & Gửi yêu cầu rồi chờ phản hồi &
Gửi yêu cầu rồi chờ phản hồi \\\hline
Giao thức truyền tải bên dưới & TCP & TCP & UDP \\\hline
Phần dữ liệu phụ đi kèm mỗi gói tin & 2 byte (tối thiểu) & Hàng trăm
byte & 4 byte \\\hline
Chất lượng dịch vụ (QoS) & 3 mức (0, 1, 2) & Không có & 2 mức
(có xác nhận / không xác nhận) \\\hline
Tiêu thụ băng thông & Thấp & Cao & Thấp \\\hline
Hỗ trợ trao đổi hai chiều & Có (đăng ký nhận bản tin) & Không, trừ khi
bổ sung cơ chế hỏi lặp hoặc kênh đẩy & Có (Observe) \\\hline
Phù hợp cho IoT & Rất tốt & Trung bình & Tốt \\\hline
Độ phổ biến & Rất cao & Rất cao & Trung bình \\\hline
Mức phù hợp với EMQX & Có sẵn & Có, nhưng phải bổ sung thành phần mở rộng
& Hạn chế \\\hline
\end{longtable}
}

\textbf{Phân tích lựa chọn:} Kết quả so sánh cho thấy MQTT phù hợp hơn
với hệ thống vì cách nó trao đổi dữ liệu bám sát nhịp vận hành thực tế:
thiết bị gửi nhiều bản tin nhỏ qua mạng di động không phải lúc nào cũng
ổn định, đồng thời đôi lúc cần nhận lệnh ngược từ máy chủ. Cụ thể:

\begin{itemize}
\tightlist
\item
  \textbf{Phần dữ liệu phụ đi kèm rất nhỏ}: Header tối thiểu 2 byte, giảm băng thông và
  tiết kiệm năng lượng cho thiết bị IoT chạy bằng pin.
\item
  \textbf{Mô hình phát bản tin vào kênh chung rồi đăng ký nhận}: Cho phép nhiều thành phần nhận dữ
  liệu cùng lúc mà không làm thiết bị trên xe phải gửi lặp đi lặp lại
  cho từng nơi nhận. Điều này phù hợp với kiến trúc có giao diện quản trị, dịch
  vụ xử lý và các kênh cảnh báo cùng khai thác một nguồn dữ liệu.
\item
  \textbf{QoS linh hoạt}: Hệ thống không cần mọi bản tin đều được bảo
  đảm ở cùng một mức. QoS 0 phù hợp với dữ liệu vận hành gửi dày theo
  chu kỳ
  như vị trí GPS, trong khi QoS 1 phù hợp hơn với cảnh báo và lệnh điều
  khiển vì các bản tin này cần được đảm bảo tới nơi ít nhất một lần.
\item
  \textbf{Lưu trạng thái cuối và thông báo khi thiết bị mất kết nối}: Hai cơ chế này giúp hệ thống
  phản ánh trạng thái online/offline của thiết bị rõ ràng hơn, từ đó
  người vận hành có thể phân biệt giữa "xe không có dữ liệu mới" và
  "thiết bị vừa mất kết nối".
\end{itemize}

\subsubsection{2.2.4. So sánh vi điều khiển
(MCU)}\label{224-so-suxe1nh-vi-ux111iux1ec1u-khiux1ec3n-mcu}

{\thesissettableid{2.3}{tab-ch2-mcu-comparison-for-tracker}
\def\LTcaptype{table}
\begin{longtable}{|L{0.230\linewidth}|L{0.218\linewidth}|L{0.260\linewidth}|L{0.246\linewidth}|}
\caption{So sánh các vi điều khiển cho ứng dụng thiết bị theo dõi IoT}\label{tab:ch2-mcu-comparison-for-tracker}\\
\hline \textbf{Tiêu chí} & \textbf{ESP32-S3} & \textbf{STM32L4} & \textbf{Raspberry Pi Zero 2W} \\
\\\hline
\endfirsthead
\caption[]{So sánh các vi điều khiển cho ứng dụng thiết bị theo dõi IoT (tiếp theo)}\\
\hline \textbf{Tiêu chí} & \textbf{ESP32-S3} & \textbf{STM32L4} & \textbf{Raspberry Pi Zero 2W} \\
\\\hline
\endhead
\\\hline
\endfoot
\endlastfoot
Kiến trúc & Xtensa LX7, dual-core, 240 MHz & ARM Cortex-M4, 80 MHz & ARM
Cortex-A53, quad-core, 1 GHz \\\hline
RAM & 512 KB SRAM + 8 MB PSRAM & 256 KB SRAM & 512 MB DRAM \\\hline
BLE & BLE 5.0 (tích hợp) & Không (cần module ngoài) & BLE 5.0 (tích
hợp) \\\hline
Wi-Fi & 802.11 b/g/n (tích hợp) & Không & 802.11 b/g/n (tích hợp) \\\hline
Tiêu thụ deep sleep & 10--15 µA & 1--2 µA & Không hỗ trợ deep sleep \\\hline
Giá thành (VND) & 80.000--150.000 & 150.000--300.000 &
400.000--600.000 \\\hline
Framework phát triển & Arduino / ESP-IDF & STM32CubeIDE / Mbed & Linux /
Python \\\hline
Cộng đồng hỗ trợ & Rất lớn & Lớn & Rất lớn \\\hline
Độ phù hợp cho thiết bị theo dõi & Cao & Trung bình & Thấp \\\hline
\end{longtable}
}

\textbf{Phân tích lựa chọn:} Nếu chỉ so riêng mức tiết kiệm điện của bản
thân vi điều khiển, STM32L4 có lợi thế rõ rệt. Nhưng thiết bị theo dõi của đề tài
không chỉ cần ``ngủ tiết kiệm'', mà còn phải đồng thời nói chuyện với
OBD2 qua BLE, modem qua UART, cảm biến qua I2C và xử lý logic trạng
thái của toàn thiết bị. Khi nhìn theo bài toán tổng thể đó, ESP32-S3 cho
điểm cân bằng thực tế hơn giữa tích hợp, chi phí và công sức triển khai.
Các lý do chính là:

\begin{itemize}
\tightlist
\item
  \textbf{BLE tích hợp sẵn}: Có thể kết nối trực tiếp với adapter OBD2
  vgate iCar Pro mà không cần gắn thêm module BLE riêng. Điều này làm
  mạch gọn hơn, ít linh kiện hơn và giảm chi phí tích hợp.
\item
  \textbf{Chi phí phù hợp}: Giá thành thấp hơn đáng kể so với STM32L4
  trong khi vẫn đủ năng lực xử lý cho bài toán của đồ án.
\item
  \textbf{Dễ phát triển và thử nghiệm}: Hệ sinh thái ESP-IDF, tài liệu
  và cộng đồng hỗ trợ cho MQTT, BLE và GNSS rất phong phú. Với nhóm làm
  đồ án, đây là lợi thế lớn vì giảm thời gian thử sai và giảm rủi ro tắc
  ở các phần tích hợp.
\item
  \textbf{Mức tiết kiệm điện vẫn đủ cho mục tiêu hệ thống}: Dòng deep
  sleep riêng của ESP32-S3 cao hơn STM32L4, nhưng tuổi thọ pin thực tế
  còn phụ thuộc cả mạch nguồn, cảm biến và linh kiện phụ trợ. Vì vậy,
  chênh lệch ở mức riêng MCU không phải yếu tố duy nhất quyết định hiệu
  quả năng lượng của toàn thiết bị.
\end{itemize}

\subsubsection{2.2.5. So sánh cơ sở dữ
liệu}\label{225-so-suxe1nh-cux1a1-sux1edf-dux1eef-liux1ec7u}

{\thesissettableid{2.4}{tab-ch2-iot-database-solutions}
\def\LTcaptype{table}
\begin{longtable}{|L{0.203\linewidth}|L{0.285\linewidth}|L{0.240\linewidth}|L{0.226\linewidth}|}
\caption{So sánh các giải pháp cơ sở dữ liệu cho hệ thống IoT}\label{tab:ch2-iot-database-solutions}\\
\hline \textbf{Tiêu chí} & \textbf{PostgreSQL + VictoriaMetrics} & \textbf{MongoDB + InfluxDB} & \textbf{TimescaleDB} \\
\\\hline
\endfirsthead
\caption[]{So sánh các giải pháp cơ sở dữ liệu cho hệ thống IoT (tiếp theo)}\\
\hline \textbf{Tiêu chí} & \textbf{PostgreSQL + VictoriaMetrics} & \textbf{MongoDB + InfluxDB} & \textbf{TimescaleDB} \\
\\\hline
\endhead
\\\hline
\endfoot
\endlastfoot
Dữ liệu quan hệ (users, vehicles) & Rất tốt (PostgreSQL) & Trung bình
(MongoDB) & Tốt (PostgreSQL core) \\\hline
Dữ liệu chuỗi thời gian (vị trí, cảm biến, trạng thái gửi liên tục) & Rất tốt (VictoriaMetrics) & Tốt
(InfluxDB) & Tốt (TimescaleDB extension) \\\hline
Hiệu suất ghi & Cao (VictoriaMetrics: hàng triệu điểm/giây) & Cao (InfluxDB:
hàng triệu điểm/giây) & Trung bình \\\hline
Nén dữ liệu & Rất tốt (VictoriaMetrics: 10--70x nén) & Tốt (InfluxDB: TSM) & Tốt
(PostgreSQL compression) \\\hline
Tài nguyên tiêu thụ & Thấp (VictoriaMetrics: 1--2 GB RAM) & Cao (InfluxDB: 4+ GB RAM)
& Trung bình \\\hline
Ngôn ngữ truy vấn & SQL (PostgreSQL) + MetricsQL (VictoriaMetrics) & MongoDB Query +
Flux/InfluxQL & SQL \\\hline
Tích hợp Grafana & Có (native) & Có & Có \\\hline
Giấy phép & Apache 2.0 & SSPL (MongoDB) + MIT (InfluxDB OSS) & Apache
2.0 \\\hline
Độ phức tạp vận hành & Trung bình (2 hệ thống riêng) & Cao (2 hệ thống
khác nhau) & Thấp (1 hệ thống) \\\hline
\end{longtable}
}

\textbf{Phân tích lựa chọn:} Với hệ thống này, không phải mọi dữ liệu
đều có cùng bản chất. Dữ liệu người dùng, xe, chuyến đi hay cảnh báo cần
quan hệ rõ ràng giữa các bảng; trong khi dữ liệu vị trí và cảm biến lại
là những bản ghi nhỏ, phát sinh liên tục theo thời gian. Vì vậy, thay vì
ép tất cả vào một chỗ, đề tài chọn mỗi loại dữ liệu vào đúng nơi phù
hợp: PostgreSQL cho dữ liệu nghiệp vụ và VictoriaMetrics cho dữ liệu
chuỗi thời gian.

\begin{itemize}
\tightlist
\item
  \textbf{Phân tách trách nhiệm rõ ràng}: PostgreSQL lưu những dữ liệu
  cần liên kết chặt như người dùng, phương tiện, cảnh báo và vùng giám
  sát; VictoriaMetrics lưu những dòng dữ liệu tăng liên tục như GPS,
  OBD2 và dữ liệu cảm biến gửi theo chu kỳ. Mỗi hệ thống chỉ phải làm đúng phần việc mà nó phù
  hợp nhất.
\item
  \textbf{Hiệu suất và tài nguyên}: VictoriaMetrics nén dữ liệu tốt và
  dùng ít RAM hơn nhiều phương án cùng loại, phù hợp với hạ tầng máy chủ
  của đồ án.
\item
  \textbf{Thuận lợi cho giám sát}: Cả hai đều tích hợp tốt với Grafana,
  nên việc dựng màn hình giám sát kỹ thuật và theo dõi vận hành không cần thêm
  lớp chuyển đổi phức tạp.
\end{itemize}

{\thesissetfigureid{2.1}{fig-ch2-data-architecture-overview}
\begin{figure}[htbp]
\centering
\pandocbounded{\safeincludesvg{./assets/figures/02-chuong-2-phan-tich-hinh-2-1.svg}}
\caption{Sơ đồ kiến trúc dữ liệu của hệ thống}
\label{fig:ch2-data-architecture-overview}
\end{figure}
}

\subsection{2.3. Yêu cầu kỹ thuật và các tiêu chuẩn thiết kế -- Design
criteria and
Constraints}\label{23-yuxeau-cux1ea7u-kux1ef9-thuux1eadt-vuxe0-cuxe1c-tiuxeau-chuux1ea9n-thiux1ebft-kux1ebf--design-criteria-and-constraints}

\subsubsection{2.3.1. Yêu cầu chức
năng}\label{231-yuxeau-cux1ea7u-chux1ee9c-nux103ng}

Những yêu cầu chức năng dưới đây đánh dấu bước chuyển từ ``bài toán''
sang ``năng lực hệ thống''. Có thể xem chúng như năm lời hứa vận hành tối
thiểu mà hệ thống phải thực hiện được để người quản lý thật sự dùng được
trong bối cảnh đội xe vận hành ngoài thực tế.

\textbf{FC-01: Theo dõi vị trí thời gian thực}

\begin{itemize}
\tightlist
\item
  Gửi vị trí GPS định kỳ (5--30 giây khi lái xe, 10--30 phút khi đậu xe)
\item
  Độ chính xác vị trí: dưới 5 mét trong điều kiện trời quang
\item
  Hỗ trợ đa hệ thống định vị (GPS, GLONASS, BeiDou)
\end{itemize}

\textbf{FC-02: Đọc dữ liệu động cơ qua OBD2}

\begin{itemize}
\tightlist
\item
  Kết nối BLE với OBD2 adapter (vgate iCar Pro)
\item
  Đọc trạng thái IGN (bật/tắt máy), RPM, tốc độ, nhiên liệu
\item
  Tần suất đọc: 5--30 giây khi xe chạy (mặc định 5 giây)
\end{itemize}

\textbf{FC-03: Phát hiện bất thường khi đậu xe}

\begin{itemize}
\tightlist
\item
  IMU (LIS3DSH) phát hiện chuyển động/rung bất thường
\item
  Đánh thức ESP32 từ deep sleep trong vòng 100ms
\item
  Gửi cảnh báo ưu tiên cao lên máy chủ
\end{itemize}

\textbf{FC-04: Quản lý năng lượng thông minh}

\begin{itemize}
\tightlist
\item
  Chuyển đổi giữa 3 chế độ: xe đang chạy (hoạt động đầy đủ), xe đang đỗ
  (tiết kiệm điện) và cảnh báo bất thường
\item
  Pin dự phòng tự động tiếp quản theo profile nguồn: hệ 12V tại ngưỡng
  \(Switch_{\mathrm{OFF}} = 12.0\,\mathrm{V}\), hệ 24V tại ngưỡng
  \(Switch_{\mathrm{OFF}} = 24.0\,\mathrm{V}\)
\item
  Có phương án dự phòng suy luận trạng thái IGN qua ADC khi kết nối OBD2
  BLE không sẵn sàng
\end{itemize}

\textbf{FC-05: Giao diện web giám sát}

\begin{itemize}
\tightlist
\item
  Bản đồ thời gian thực (Leaflet) hiển thị vị trí tất cả phương tiện
\item
  Biểu đồ dữ liệu OBD2 (ECharts) - tốc độ, RPM, nhiên liệu
\item
  Hệ thống cảnh báo và thông báo tự động
\item
  Quản lý phương tiện, tài xế, khách hàng, hành trình
\end{itemize}

Năm nhóm chức năng trên bao quát trọn chuỗi nhu cầu cốt lõi của hệ
thống: biết xe đang ở đâu, hiểu xe đang như thế nào, phát hiện khi có
bất thường, không làm phiền hệ thống điện của xe và trình bày thông tin ở
dạng mà người vận hành có thể đọc nhanh, xử lý nhanh.

\subsubsection{2.3.2. Yêu cầu phi chức
năng}\label{232-yuxeau-cux1ea7u-phi-chux1ee9c-nux103ng}

Nếu yêu cầu chức năng trả lời ``hệ thống phải làm được gì'', thì yêu cầu
phi chức năng trả lời ``hệ thống phải làm tốt tới mức nào''. Đây là lớp
tiêu chí cho biết nguyên mẫu có đủ độ chín để bước ra môi trường vận
hành thật hay chưa, đặc biệt khi phải theo dõi nhiều xe và chạy liên
tục.

\textbf{NFR-01: Hiệu suất}

\begin{itemize}
\tightlist
\item
  Độ trễ toàn tuyến (từ thiết bị đến giao diện quản trị): dưới 3 giây trong điều
  kiện mạng bình thường
\item
  Thời gian phản hồi API: dưới 200 ms cho 95\% số yêu cầu
\item
  Hỗ trợ đồng thời tối thiểu 50 phương tiện kết nối MQTT, với dư địa mở
  rộng ổn định tới khoảng 100 thiết bị
\end{itemize}

\textbf{NFR-02: Độ tin cậy}

\begin{itemize}
\tightlist
\item
  Tự động kết nối lại khi mạng phục hồi
\item
  Giữ được cấu hình và trạng thái quan trọng qua NVS; lưu đệm dữ liệu
  vận hành cục bộ là hướng mở rộng của thiết bị
\item
  MQTT QoS 1 cho các tin nhắn cảnh báo (đảm bảo gửi ít nhất một lần)
\end{itemize}

\textbf{NFR-03: Bảo mật}

\begin{itemize}
\tightlist
\item
  Xác thực phiên bằng token ngẫu nhiên không tự mang thông tin người
  dùng; phía máy chủ mới là nơi lưu và kiểm tra trạng thái phiên, còn
  trong cơ sở dữ liệu chỉ lưu bản băm SHA-256 của token
\item
  MQTT ACL (Access Control List), tức danh sách quyền truy cập theo từng
  thiết bị
\item
  Mã hóa TLS cho kết nối MQTT và HTTPS trong môi trường sản xuất
\end{itemize}

\textbf{NFR-04: Khả năng mở rộng}

\begin{itemize}
\tightlist
\item
  Kiến trúc nhiều dịch vụ tách rời, mỗi dịch vụ được đóng gói thành
  một gói chạy riêng bằng Docker
\item
  MQTT broker (EMQX) hỗ trợ chạy thành cụm, tức nhiều máy có thể cùng
  phục vụ khi số thiết bị tăng
\item
  Cơ sở dữ liệu chuỗi thời gian (VictoriaMetrics) tối ưu cho luồng dữ
  liệu lớn
\end{itemize}

\textbf{NFR-05: Khả năng bảo trì}

\begin{itemize}
\tightlist
\item
  Mã nguồn tổ chức theo Domain-Driven Design (DDD), nghĩa là chia theo
  từng nhóm nghiệp vụ như xe, thiết bị hay cảnh báo thay vì gom theo
  loại file kỹ thuật
\item
  API theo kiểu REST với tài liệu tự sinh bằng Swagger/OpenAPI, tức
  trang mô tả sẵn các API để người phát triển và người kiểm thử cùng tra
  cứu
\item
  Ghi log tập trung với VictoriaLogs và theo dõi trên Grafana, tức màn
  hình giám sát chung của hệ thống
\end{itemize}

Nhóm yêu cầu trên mô tả hệ thống cần làm được gì và cần đạt tới mức nào.
Phần tiếp theo thu hẹp thêm một bước: trong khuôn khổ đồ án, nhóm thực
sự bị giới hạn bởi những điều kiện nào về phần cứng, phần mềm và chi
phí.

\subsubsection{2.3.3. Ràng buộc thiết
kế}\label{233-ruxe0ng-buux1ed9c-thiux1ebft-kux1ebf}

Các ràng buộc sau đây giới hạn không gian thiết kế của đề tài. Chúng buộc
nhóm phải chọn phương án cân bằng giữa tính khả thi học thuật, khả năng
triển khai thực tế trên xe và ngân sách sẵn có.

\textbf{Ràng buộc phần cứng:}

\begin{itemize}
\tightlist
\item
  Điện áp đầu vào: 12V hoặc 24V DC từ ắc quy xe (dao động phụ thuộc cấu
  hình hệ thống điện)
\item
  Nhiệt độ hoạt động: -10 độ C đến 70 độ C (môi trường trong xe ô tô)
\item
  Rung động và sốc: Chịu được rung động liên tục khi xe vận hành trên
  đường xá
\item
  Kích thước: Đủ nhỏ để lắp đặt kín đáo trong xe (không lớn hơn
  120x80x40 mm)
\end{itemize}

\textbf{Ràng buộc phần mềm:}

\begin{itemize}
\tightlist
\item
  Khung phát triển API máy chủ: ưu tiên Express.js + TypeScript vì đây là nền tảng
  nhóm đã có sẵn kinh nghiệm, giúp rút ngắn thời gian học lại từ đầu
\item
  Framework giao diện web: chọn Next.js 15 + React 19 để phù hợp định hướng
  học thuật và thuận lợi cho màn hình quản trị có nhiều khu vực hiển thị
\item
  MQTT Broker: chọn EMQX vì là giải pháp mã nguồn mở ổn định, đủ đáp ứng
  vai trò broker MQTT cho giai đoạn hiện tại
\item
  Triển khai: dùng Docker Compose trên máy chủ Linux để mỗi dịch vụ có
  thể chạy tách riêng nhưng vẫn dễ dựng lại môi trường
\end{itemize}

\textbf{Ràng buộc về chi phí:}

\begin{itemize}
\tightlist
\item
  Tổng chi phí phần cứng: dưới 2.000.000 VND cho một bộ thiết bị theo dõi
\item
  Ưu tiên các linh kiện có datasheet rõ ràng, chất lượng ổn định và
  nguồn cung trong nước dễ tiếp cận
\item
  Ưu tiên các giải pháp mã nguồn mở để giảm chi phí giấy phép
\end{itemize}

\subsubsection{2.3.4. Tiêu chuẩn thiết kế áp
dụng}\label{234-tiuxeau-chuux1ea9n-thiux1ebft-kux1ebf-uxe1p-dux1ee5ng}

Bảng dưới đây tổng hợp các chuẩn kỹ thuật và quy ước triển khai được dùng
làm nền cho hệ thống. Điểm quan trọng không phải là gom thật nhiều chuẩn,
mà là chọn đúng những chuẩn thực sự chi phối khả năng tương thích giữa
thiết bị, mạng truyền và ứng dụng.


{\thesissettableid{2.3.4A}{tab-ch2-tong-hop-thong-tin-ve-tieu-chuan-mo-ta-ap-dung-trong}
\def\LTcaptype{table}
\begin{longtable}{|L{0.231\linewidth}|L{0.278\linewidth}|L{0.455\linewidth}|}
\caption{Tổng hợp chuẩn và quy ước kỹ thuật áp dụng trong hệ thống}\label{tab:ch2-tong-hop-thong-tin-ve-tieu-chuan-mo-ta-ap-dung-trong}\\
\hline \textbf{Chuẩn / quy ước} & \textbf{Mô tả} & \textbf{Áp dụng trong hệ thống} \\
\\\hline
\endfirsthead
\caption[]{Tổng hợp chuẩn và quy ước kỹ thuật áp dụng trong hệ thống (tiếp theo)}\\
\hline \textbf{Chuẩn / quy ước} & \textbf{Mô tả} & \textbf{Áp dụng trong hệ thống} \\
\\\hline
\endhead
\\\hline
\endfoot
\endlastfoot
IEEE 802.15.1 (Bluetooth) & Chuẩn BLE 4.0/5.0 & Kết nối OBD2 adapter \\\hline
3GPP LTE Cat-4 & Chuẩn 4G/LTE & Truyền dữ liệu qua modem SIM7600CE-T
(Cat-4) \\\hline
NMEA 0183 & Chuẩn dữ liệu GPS & Phân tích tọa độ từ khối GNSS tích hợp \\\hline
SAE J1979 (OBD2) & Chuẩn chẩn đoán động cơ & Đọc dữ liệu qua OBD2 BLE
adapter \\\hline
MQTT v5.0 (OASIS) & Giao thức nhắn tin cho thiết bị IoT & Truyền dữ liệu thiết bị -
máy chủ \\\hline
REST (RFC 7231) & Cách tổ chức API theo tài nguyên & API máy chủ \\\hline
Mã phiên ngẫu nhiên + ACL MQTT & Cơ chế xác thực và phân quyền &
Xác thực phiên người dùng, giới hạn quyền truy cập thiết bị theo ACL
MQTT \\\hline
\end{longtable}
}


\subsection{2.4. Yêu cầu từ các bên liên quan -- Stakeholder
requirements}\label{24-yuxeau-cux1ea7u-tux1eeb-cuxe1c-buxean-liuxean-quan--stakeholder-requirements}

\subsubsection{2.4.1. Xác định các bên liên
quan}\label{241-xuxe1c-ux111ux1ecbnh-cuxe1c-buxean-liuxean-quan}

Hệ thống IoT Vehicle Tracking System liên quan đến nhiều nhóm sử dụng với
nhu cầu khác nhau. Nếu không tách rõ ngay từ đầu, hệ thống rất dễ rơi vào
tình trạng mạnh về mặt kỹ thuật nhưng lệch khỏi nhu cầu vận hành. Vì vậy,
phần này xác định ai là người hưởng lợi trực tiếp, ai là người vận hành,
ai là người bị tác động, từ đó giữ trọng tâm thiết kế vào những nhu cầu
thực sự cần giải quyết.

\subsubsection{2.4.2. Ma trận yêu cầu các bên liên
quan}\label{242-ma-trux1eadn-yuxeau-cux1ea7u-cuxe1c-buxean-liuxean-quan}


{\thesissettableid{2.5}{tab-ch2-ma-tran-yeu-cau-cac-ben-lien-quan}
\def\LTcaptype{table}
\begin{longtable}{|L{0.085\linewidth}|L{0.191\linewidth}|L{0.162\linewidth}|L{0.380\linewidth}|L{0.126\linewidth}|}
\caption{Ma trận yêu cầu các bên liên quan}\label{tab:ch2-ma-tran-yeu-cau-cac-ben-lien-quan}\\
\hline \textbf{STT} & \textbf{Bên liên quan} & \textbf{Vai trò} & \textbf{Yêu cầu chính} & \textbf{Mức độ ưu tiên} \\
\\\hline
\endfirsthead
\caption[]{Ma trận yêu cầu các bên liên quan (tiếp theo)}\\
\hline \textbf{STT} & \textbf{Bên liên quan} & \textbf{Vai trò} & \textbf{Yêu cầu chính} & \textbf{Mức độ ưu tiên} \\
\\\hline
\endhead
\\\hline
\endfoot
\endlastfoot
1 & Công ty cho thuê xe & Chủ sở hữu hệ thống & Giám sát vị trí tất cả
xe theo thời gian thực; Nhận cảnh báo khi xe đi ra khỏi vùng cho phép; Báo cáo hành trình, quãng đường, nhiên liệu; Giảm thiểu thời
gian xe đậu không sử dụng & Cao \\\hline
2 & Tài xế / Người thuê xe & Người vận hành xe & Không bị giám sát quá
mức (quyền riêng tư); Thiết bị không ảnh hưởng đến vận hành xe; Không
rút cạn ắc quy xe & Trung bình \\\hline
3 & Quản trị viên hệ thống & Người quản lý kỹ thuật & Giao diện quản lý
dễ sử dụng; Hệ thống ổn định, ít lỗi; Có khả năng theo dõi trạng thái
thiết bị (online/offline); Quản lý người dùng và phân quyền & Cao \\\hline
4 & Đội bảo trì & Kỹ thuật viên & Thiết bị dễ lắp đặt và bảo trì; Có thể
kiểm tra tình trạng thiết bị từ xa mà không cần tháo xuống; Cảnh báo bảo
trì định kỳ (pin yếu, thiết bị ngoại tuyến) & Trung bình \\\hline
5 & Khách hàng (người thuê xe cuối) & Người dùng gián tiếp & Quá trình
thuê xe nhanh gọn; Tin tưởng vào độ bảo mật của dữ liệu cá nhân; Nhận
thông báo về trạng thái đơn thuê & Thấp (đưa vào giai đoạn 2) \\\hline
\end{longtable}
}


\subsubsection{2.4.3. Phân tích chi tiết yêu
cầu}\label{243-phuxe2n-tuxedch-chi-tiux1ebft-yuxeau-cux1ea7u}

Từ ma trận trên, có thể thấy không phải mọi bên liên quan đều đòi hỏi cùng
một loại giá trị. Nhóm sử dụng chính quan tâm đến hiệu quả vận hành và an
ninh; nhóm quản trị quan tâm đến tính ổn định và khả năng kiểm soát; còn
nhóm bảo trì cần hệ thống dễ chẩn đoán và dễ can thiệp khi có sự cố.

\textbf{Công ty cho thuê xe (Stakeholder chính)}

Đây là nhóm đối tượng trọng tâm của hệ thống. Các yêu cầu của nhóm này
tập trung vào ba khía cạnh:

\begin{enumerate}
\def\labelenumi{\arabic{enumi}.}
\tightlist
\item
  \emph{Giám sát và an ninh}: Theo dõi vị trí phương tiện 24/7, phát
  hiện và cảnh báo ngay khi có bất thường như xe bị di chuyển trái phép,
  vượt khỏi vùng giám sát đã vẽ trên bản đồ hoặc chạy quá tốc độ. Hệ
  thống cần gửi thông báo qua nhiều kênh như giao diện web quản trị,
  Telegram bot và email.
\item
  \emph{Báo cáo và phân tích}: Xuất báo cáo hành trình chi tiết (quãng
  đường, thời gian, điểm dừng), thống kê nhiên liệu tiêu thụ, đánh giá
  hành vi lái xe (tốc độ trung bình, số lần phanh gấp). Các báo cáo này
  giúp tối ưu hóa chi phí vận hành và bảo trì.
\item
  \emph{Quản lý đội xe}: Quản lý trạng thái từng phương tiện (đang cho
  thuê, đang bảo trì, sẵn sàng), lịch bảo trì định kỳ dựa trên số km
  hoặc thời gian, lịch sử cho thuê và doanh thu theo phương tiện.
\end{enumerate}

\textbf{Quản trị viên hệ thống}

Yêu cầu của quản trị viên tập trung vào khả năng vận hành và giám sát hệ
thống:

\begin{enumerate}
\def\labelenumi{\arabic{enumi}.}
\tightlist
\item
  \emph{Trang tổng quan}: Hiển thị trạng thái tất cả thiết bị
  (online/offline, mức pin, cường độ tín hiệu), số lượng phương tiện
  đang hoạt động, cảnh báo chưa xử lý.
\item
  \emph{Quản lý thiết bị}: Thêm/xóa/sửa thông tin thiết bị theo dõi, cấp
  phát thiết bị cho phương tiện, gửi lệnh điều khiển từ xa (reset, cập
  nhật cấu hình, yêu cầu vị trí).
\item
  \emph{Quản lý người dùng}: Tạo tài khoản, phân quyền (admin, manager,
  viewer), theo dõi lịch sử đăng nhập và thao tác thông qua audit log,
  tức bản ghi ai đã làm gì trong hệ thống.
\end{enumerate}

\textbf{Đội bảo trì}

Yêu cầu của đội bảo trì hướng đến việc giảm thời gian và chi phí bảo
trì:

\begin{enumerate}
\def\labelenumi{\arabic{enumi}.}
\tightlist
\item
  \emph{Chẩn đoán từ xa}: Xem trạng thái thiết bị (mức pin, nhiệt độ,
  cường độ tín hiệu 4G) từ xa, phát hiện và cảnh báo khi thiết bị gặp
  lỗi (mất kết nối kéo dài, pin yếu).
\item
  \emph{Hướng dẫn lắp đặt}: Tài liệu hướng dẫn lắp đặt chi tiết, danh
  sách kiểm tra (checklist) khi lắp đặt mới hoặc bảo trì.
\end{enumerate}

Nhìn tổng thể, ba nhóm yêu cầu này kéo hệ thống về một điểm cân bằng khá
rõ: phải đủ mạnh để quản lý đội xe, nhưng cũng phải đủ gọn và ổn định để
lắp đặt, vận hành và bảo trì trên thực địa. Đây cũng là cơ sở để ma trận
truy xuất ở mục tiếp theo ánh xạ yêu cầu định tính sang các chức năng cụ
thể của hệ thống.

\subsubsection{2.4.4. Ma trận truy xuất yêu cầu - chức
năng}\label{244-ma-trux1eadn-truy-xuux1ea5t-yuxeau-cux1ea7u---chux1ee9c-nux103ng}


{\thesissettableid{2.6}{tab-ch2-ma-tran-truy-xuat-yeu-cau-cac-ben-lien-quan-voi-chuc}
\def\LTcaptype{table}
\begin{longtable}{|L{0.237\linewidth}|L{0.122\linewidth}|L{0.122\linewidth}|L{0.186\linewidth}|L{0.168\linewidth}|L{0.101\linewidth}|}
\caption{Ma trận truy xuất yêu cầu các bên liên quan với chức năng hệ thống}\label{tab:ch2-ma-tran-truy-xuat-yeu-cau-cac-ben-lien-quan-voi-chuc}\\
\hline \textbf{Yêu cầu bên liên quan} & \textbf{FC-01 (Vị trí)} & \textbf{FC-02 (OBD2)} & \textbf{FC-03 (Bất
thường)} & \textbf{FC-04 (Năng lượng)} & \textbf{FC-05 (Web)} \\
\\\hline
\endfirsthead
\caption[]{Ma trận truy xuất yêu cầu các bên liên quan với chức năng hệ thống (tiếp theo)}\\
\hline \textbf{Yêu cầu bên liên quan} & \textbf{FC-01 (Vị trí)} & \textbf{FC-02 (OBD2)} & \textbf{FC-03 (Bất
thường)} & \textbf{FC-04 (Năng lượng)} & \textbf{FC-05 (Web)} \\
\\\hline
\endhead
\\\hline
\endfoot
\endlastfoot
Giám sát vị trí 24/7 & X & & & X & X \\\hline
Cảnh báo vượt vùng cho phép & X & & & & X \\\hline
Báo cáo nhiên liệu & & X & & & X \\\hline
Phát hiện trộm xe & & & X & X & X \\\hline
Không rút cạn ắc quy & & & & X & \\\hline
Quản lý thiết bị từ xa & & & & & X \\\hline
Chẩn đoán lỗi từ xa & X & X & & X & X \\\hline
\end{longtable}
}


Ma trận truy xuất trên cho thấy các chức năng hệ thống (FC-01 đến FC-05)
bao phủ toàn bộ yêu cầu của các bên liên quan chính. Trong đó, FC-04
(Quản lý năng lượng) và FC-05 (Giao diện web) là hai chức năng được nhắc
đến nhiều nhất, qua đó khẳng định tầm quan trọng của thiết kế quản lý
năng lượng thông minh và giao diện giám sát thân thiện.

\begin{center}\rule{0.5\linewidth}{0.5pt}\end{center}

\chapterheading{CHƯƠNG 3. CÁC GIẢI PHÁP THIẾT KẾ -- DESIGN SOLUTIONS}

Chương này trình bày bản rút gọn của các quyết định thiết kế cốt lõi. Trọng tâm là giữ rõ chuỗi lập luận: bài toán kỹ thuật, tiêu chí lựa chọn, phương án được chọn và tác động tới kết quả triển khai ở Chương 4.

\subsection{3.1. Phân tích và đề xuất giải pháp phần cứng -- Hardware analysis and solution}

\subsubsection{3.1.1. Phân tích và lựa chọn phần cứng}

Thiết bị theo dõi cần thỏa đồng thời năm tiêu chí chính: ổn định điện áp khi xe rung nhiễu, tiêu thụ điện thấp khi xe dừng lâu, truyền dữ liệu tin cậy qua mạng di động, tương thích OBD2 và chi phí phù hợp để có thể nhân rộng.

Nhóm phần cứng được chốt theo kiến trúc một bộ điều phối trung tâm (ESP32-S3) kết hợp ba khối dữ liệu (OBD2, GNSS, IMU) và một khối nguồn hai nhánh (nguồn xe + pin dự phòng). Cách ghép này giảm độ phức tạp liên kết, thuận lợi bảo trì và phù hợp nhu cầu nâng cấp firmware OTA.

{\thesissettableid{3.1}{tab-condensed-ch3-hardware-main-components}
\def\LTcaptype{table}
\begin{longtable}{|L{0.220\linewidth}|L{0.240\linewidth}|L{0.500\linewidth}|}
\caption{Tóm tắt thành phần phần cứng chính và vai trò}\label{tab:condensed-ch3-hardware-main-components}\\
\hline \textbf{Khối} & \textbf{Lựa chọn} & \textbf{Vai trò kỹ thuật chính} \\
\\\hline
\endfirsthead
\caption[]{Tóm tắt thành phần phần cứng chính và vai trò (tiếp theo)}\\
\hline \textbf{Khối} & \textbf{Lựa chọn} & \textbf{Vai trò kỹ thuật chính} \\
\\\hline
\endhead
\\\hline
\endfoot
\endlastfoot
Điều phối trung tâm & ESP32-S3 & Điều phối tác vụ thời gian thực, quản lý BLE/UART/I2C/ADC, đóng gói dữ liệu trước khi phát MQTT \\\hline
Kết nối OBD2 & vgate iCar Pro BLE & Đọc thông số vận hành xe qua giao tiếp BLE, phù hợp hạ tầng không can thiệp sâu vào dây dẫn nội bộ xe \\\hline
Truyền thông + định vị & SIM7600CE-T & Cung cấp LTE để truyền dữ liệu và GNSS để xác định vị trí, đáp ứng yêu cầu theo dõi liên tục \\\hline
Phát hiện chuyển động & LIS3DSH & Nhận biết trạng thái đứng yên/chuyển động để điều khiển chế độ tiết kiệm năng lượng \\\hline
Khối nguồn & Buck/LDO + pin 18650 1S & Tách nguồn cho MCU và modem, duy trì hoạt động khi mất nguồn xe và bảo vệ ắc quy \\\hline
\end{longtable}
}

\subsubsection{3.1.2. Giải pháp phần cứng}

Kiến trúc phần cứng được tổ chức theo nguyên tắc: ưu tiên độ tin cậy trước, tối ưu chi phí sau. Điều này thể hiện ở ba quyết định quan trọng.

\begin{itemize}
\tightlist
\item
  Tách nhánh nguồn cho modem và MCU để tránh sụt áp cục bộ làm reset hệ thống.
\item
  Dùng ngưỡng Low Voltage Disconnect để bảo vệ ắc quy xe, tránh làm giảm khả năng đề nổ.
\item
  Dùng pin dự phòng để duy trì phiên truyền thông trong các tình huống mất nguồn đột ngột.
\end{itemize}

Nhờ ba quyết định này, thiết bị giữ được tính liên tục dữ liệu và giảm số lần khởi động lại trong điều kiện vận hành thực tế.

\subsection{3.2. Phân tích và đề xuất giải pháp Firmware -- Firmware analysis and solution}

\subsubsection{3.2.1. Phân tích và lựa chọn giải pháp Firmware}

Firmware được chọn theo hướng event-driven kết hợp FreeRTOS task tách vai trò. Mục tiêu không phải chạy nhiều tính năng đồng thời, mà là bảo đảm mỗi tác vụ có vòng đời rõ và có phương án phục hồi riêng khi lỗi.

{\thesissettableid{3.2}{tab-condensed-ch3-firmware-task-roles}
\def\LTcaptype{table}
\begin{longtable}{|L{0.220\linewidth}|L{0.300\linewidth}|L{0.440\linewidth}|}
\caption{Tóm tắt tác vụ firmware theo vai trò}\label{tab:condensed-ch3-firmware-task-roles}\\
\hline \textbf{Tác vụ} & \textbf{Mục tiêu} & \textbf{Điểm kiểm soát chính} \\
\\\hline
\endfirsthead
\caption[]{Tóm tắt tác vụ firmware theo vai trò (tiếp theo)}\\
\hline \textbf{Tác vụ} & \textbf{Mục tiêu} & \textbf{Điểm kiểm soát chính} \\
\\\hline
\endhead
\\\hline
\endfoot
\endlastfoot
OBD2 task & Thu thập dữ liệu vận hành xe & Timeout kết nối BLE, retry có giới hạn, đánh dấu chất lượng mẫu dữ liệu \\\hline
Modem task & Quản lý LTE/GNSS và đường truyền MQTT & Theo dõi trạng thái mạng, phát hiện mất kết nối, gửi bù dữ liệu khi phục hồi \\\hline
Power task & Điều phối chế độ nguồn & Chuyển trạng thái chạy/ngủ dựa trên IMU, điện áp và trạng thái kết nối \\\hline
Telemetry task & Chuẩn hóa và đóng gói bản tin & Gán timestamp, kiểm tra schema, phân loại topic theo mức ưu tiên \\\hline
OTA/Control task & Xử lý lệnh điều khiển từ xa & Kiểm tra điều kiện an toàn trước khi áp dụng cấu hình hoặc cập nhật \\\hline
\end{longtable}
}

\subsubsection{3.2.2. Giải pháp Firmware}

Nội dung triển khai firmware được giản lược theo ba luồng chính:

\begin{itemize}
\tightlist
\item
  Luồng thu thập: OBD2/GNSS/IMU -> chuẩn hóa -> phân loại.
\item
  Luồng truyền dữ liệu: đóng gói MQTT theo topic nghiệp vụ, có cơ chế xác nhận và gửi lại có kiểm soát.
\item
  Luồng năng lượng: chuyển giữa chế độ hoạt động và chế độ tiết kiệm để giảm tiêu thụ nhưng không mất khả năng giám sát.
\end{itemize}

Các đoạn mã chi tiết và API nội bộ được lược bỏ khỏi bản rút gọn; thay vào đó là mô tả kiến trúc và nguyên tắc điều phối.

\subsection{3.3. Phân tích và đề xuất kiến trúc máy chủ và xử lý dữ liệu -- Server and data-processing architecture}

\subsubsection{3.3.1. Phân tích và lựa chọn kiến trúc máy chủ}

Kiến trúc máy chủ được chọn theo hướng dịch vụ tách vai trò: broker nhận bản tin, bridge xử lý/định tuyến, backend cung cấp API và realtime, lớp lưu trữ tách dữ liệu giao dịch và dữ liệu chuỗi thời gian.

Lý do lựa chọn là để dễ mở rộng theo tải thiết bị và giảm ảnh hưởng dây chuyền khi một dịch vụ gặp sự cố.

\subsubsection{3.3.2. Giải pháp dịch vụ máy chủ và xử lý nghiệp vụ}

{\thesissettableid{3.3}{tab-condensed-ch3-cloud-services}
\def\LTcaptype{table}
\begin{longtable}{|L{0.190\linewidth}|L{0.260\linewidth}|L{0.490\linewidth}|}
\caption{Tóm tắt dịch vụ máy chủ trong kiến trúc rút gọn}\label{tab:condensed-ch3-cloud-services}\\
\hline \textbf{Lớp} & \textbf{Dịch vụ} & \textbf{Nhiệm vụ chính} \\
\\\hline
\endfirsthead
\caption[]{Tóm tắt dịch vụ máy chủ trong kiến trúc rút gọn (tiếp theo)}\\
\hline \textbf{Lớp} & \textbf{Dịch vụ} & \textbf{Nhiệm vụ chính} \\
\\\hline
\endhead
\\\hline
\endfoot
\endlastfoot
Message broker & EMQX & Nhận/publish bản tin MQTT, tách luồng theo topic nghiệp vụ \\\hline
Ingestion & MQTT Bridge & Kiểm tra cấu trúc bản tin, chuyển đổi định dạng và ghi dữ liệu vào kho lưu trữ phù hợp \\\hline
API và realtime & Backend API + Socket.IO & Cung cấp endpoint cho dashboard, phát sự kiện thời gian thực cho bản đồ và cảnh báo \\\hline
Lưu trữ giao dịch & PostgreSQL & Lưu metadata thiết bị, cấu hình, cảnh báo, dữ liệu phục vụ quản trị \\\hline
Lưu trữ telemetry & VictoriaMetrics/VictoriaLogs & Lưu chuỗi thời gian và nhật ký phục vụ truy vấn nhanh, giám sát và phân tích \\\hline
\end{longtable}
}

Các phần endpoint chi tiết, cấu trúc thư mục, trích mã xử lý domain được rút gọn để tập trung vào luồng dữ liệu và nguyên tắc thiết kế.

\subsection{3.4. Phân tích và đề xuất công nghệ giao diện web -- Web interface analysis and solution}

\subsubsection{3.4.1. Phân tích và lựa chọn công nghệ giao diện web}

Bộ công nghệ frontend được chọn theo ba tiêu chí: thời gian hiển thị ban đầu tốt, dễ tổ chức theo tính năng nghiệp vụ, và hỗ trợ realtime ổn định. Next.js, React Query, Zustand và Leaflet đáp ứng tốt tổ hợp tiêu chí này.

\subsubsection{3.4.2. Giải pháp giao diện web}

Bản rút gọn giữ các phần cốt lõi của giao diện:

\begin{itemize}
\tightlist
\item
  Trang tổng quan vận hành với trạng thái thiết bị theo thời gian thực.
\item
  Trang bản đồ hiển thị vị trí xe, vùng giám sát và cảnh báo.
\item
  Trang lịch sử truy vết để phân tích tuyến đường và sự kiện quan trọng.
\item
  Trang quản trị cấu hình thiết bị và chính sách cảnh báo.
\end{itemize}

Mô tả chi tiết cấp component/file được chuyển về phụ lục hỗ trợ, không trình bày trong thân báo cáo rút gọn.

\subsection{3.5. Phân tích, đánh giá và lựa chọn phương án khả thi -- Analysis, Evaluation and Selection}

{\thesissettableid{3.4}{tab-condensed-ch3-feasibility-matrix}
\def\LTcaptype{table}
\begin{longtable}{|L{0.250\linewidth}|L{0.220\linewidth}|L{0.220\linewidth}|L{0.250\linewidth}|}
\caption{Ma trận đánh giá phương án theo tiêu chí trọng yếu}\label{tab:condensed-ch3-feasibility-matrix}\\
\hline \textbf{Tiêu chí} & \textbf{Mức ưu tiên} & \textbf{Phương án chọn} & \textbf{Nhận định} \\
\\\hline
\endfirsthead
\caption[]{Ma trận đánh giá phương án theo tiêu chí trọng yếu (tiếp theo)}\\
\hline \textbf{Tiêu chí} & \textbf{Mức ưu tiên} & \textbf{Phương án chọn} & \textbf{Nhận định} \\
\\\hline
\endhead
\\\hline
\endfoot
\endlastfoot
Độ ổn định vận hành & Rất cao & Kiến trúc tách lớp + dự phòng nguồn & Giảm rủi ro mất dữ liệu khi xe mất nguồn hoặc mạng dao động \\\hline
Khả năng mở rộng & Cao & Tách broker/bridge/backend & Cho phép mở rộng theo từng cụm dịch vụ thay vì nâng toàn hệ thống cùng lúc \\\hline
Chi phí triển khai & Cao & Tận dụng công nghệ mã nguồn mở & Kiểm soát tốt tổng chi phí hệ thống và chi phí mở rộng \\\hline
Khả năng bảo trì & Cao & Chuẩn hóa luồng dữ liệu + module rõ vai trò & Dễ truy vết lỗi và rút ngắn thời gian xử lý sự cố \\\hline
\end{longtable}
}

\subsection{3.6. Tối ưu phương án thiết kế -- The optimal solution}

Phương án tối ưu được chốt theo hướng cân bằng giữa hiệu năng, chi phí và khả năng vận hành dài hạn. Điểm cốt lõi của phương án là thiết kế đồng bộ từ thiết bị đến dashboard theo cùng một triết lý: dữ liệu đi theo luồng chuẩn, trạng thái có thể quan sát, lỗi có thể cô lập và phục hồi.

Với cấu trúc này, hệ thống đủ nhẹ để triển khai thực tế ở quy mô nhỏ, nhưng vẫn có đường nâng cấp rõ cho quy mô lớn hơn.

\chapterheading{CHƯƠNG 4. TRIỂN KHAI GIẢI PHÁP VÀ KẾT QUẢ -- IMPLEMENTATION AND RESULTS}

Chương 4 tập trung vào những gì đã triển khai và những gì đã đo được. Phần mô tả chi tiết từng dòng lệnh, từng endpoint và cấu trúc mã nguồn được lược bỏ; thay vào đó là bằng chứng triển khai, số liệu đo lường trọng yếu và đối chiếu với chỉ tiêu thiết kế.

\subsection{4.1. Thiết kế chi tiết giải pháp -- Detailed design solution}

\subsubsection{4.1.1. Thiết kế chi tiết phần cứng}

Thiết kế chi tiết phần cứng được chốt theo nguyên tắc an toàn nguồn và độ bền khi lắp đặt trên xe. Các quyết định quan trọng gồm: tách nhánh cấp nguồn cho modem, giám sát điện áp ắc quy, dùng pin dự phòng và bố trí module theo hướng giảm nhiễu xuyên kênh.

Bản rút gọn chỉ giữ các thông số then chốt phục vụ đánh giá kỹ thuật; toàn bộ bản vẽ và chi tiết dây nối đầy đủ được đưa về tài liệu hỗ trợ.

{\thesissettableid{4.1}{tab-condensed-ch4-hardware-key-specs}
\def\LTcaptype{table}
\begin{longtable}{|L{0.250\linewidth}|L{0.300\linewidth}|L{0.410\linewidth}|}
\caption{Thông số kỹ thuật then chốt của nguyên mẫu triển khai}\label{tab:condensed-ch4-hardware-key-specs}\\
\hline \textbf{Nhóm} & \textbf{Thông số chính} & \textbf{Ý nghĩa kỹ thuật} \\
\\\hline
\endfirsthead
\caption[]{Thông số kỹ thuật then chốt của nguyên mẫu triển khai (tiếp theo)}\\
\hline \textbf{Nhóm} & \textbf{Thông số chính} & \textbf{Ý nghĩa kỹ thuật} \\
\\\hline
\endhead
\\\hline
\endfoot
\endlastfoot
Nguồn vào & 12V/24V từ cổng OBD2 & Tương thích nhiều dòng xe, giảm yêu cầu phụ kiện chuyển đổi \\\hline
Nguồn dự phòng & Pin Li-ion 18650 1S & Duy trì phiên truyền dữ liệu khi mất nguồn xe ngắn hạn \\\hline
Điều phối & ESP32-S3 & Đủ tài nguyên xử lý đồng thời BLE, modem, IMU và telemetry \\\hline
Kết nối xe & BLE OBD2 qua adapter & Hạn chế can thiệp sâu vào hệ thống điện xe, thuận lợi lắp đặt thực địa \\\hline
Định vị và truyền dữ liệu & LTE + GNSS tích hợp & Đảm bảo tính liên tục dữ liệu vị trí và vận hành \\\hline
\end{longtable}
}

\subsubsection{4.1.2. Triển khai firmware, hạ tầng máy chủ và giao diện điều khiển}

Triển khai phần mềm được thực hiện theo hướng đồng bộ phiên bản giữa firmware, backend và frontend. Quy trình phát hành gồm ba bước: kiểm tra tích hợp nội bộ, triển khai staging để mô phỏng tải thiết bị, sau đó mới chuyển sang vận hành thử nghiệm.

Việc chuẩn hóa quy trình triển khai giúp giảm sai lệch giữa môi trường phát triển và môi trường vận hành, đồng thời tăng khả năng truy vết khi phát sinh lỗi liên tầng.

\subsection{4.2. Chế tạo và lắp ráp hệ thống -- Manufacture and Assembly}

Phần lắp ráp được rút gọn theo các mốc chính: chuẩn bị vật tư, lắp bo mạch, kiểm tra nguồn, lắp lên xe, chạy checklist vận hành và nghiệm thu thử nghiệm. Các hướng dẫn thao tác chi tiết theo ảnh và bước công việc được chuyển sang phụ lục hỗ trợ triển khai.

{\thesissettableid{4.2}{tab-condensed-ch4-deployment-checklist}
\def\LTcaptype{table}
\begin{longtable}{|L{0.330\linewidth}|L{0.240\linewidth}|L{0.390\linewidth}|}
\caption{Checklist triển khai vận hành rút gọn}\label{tab:condensed-ch4-deployment-checklist}\\
\hline \textbf{Nhóm kiểm tra} & \textbf{Trạng thái yêu cầu} & \textbf{Mục tiêu} \\
\\\hline
\endfirsthead
\caption[]{Checklist triển khai vận hành rút gọn (tiếp theo)}\\
\hline \textbf{Nhóm kiểm tra} & \textbf{Trạng thái yêu cầu} & \textbf{Mục tiêu} \\
\\\hline
\endhead
\\\hline
\endfoot
\endlastfoot
Nguồn và bảo vệ ắc quy & Đạt & Không gây sụt áp ảnh hưởng khả năng khởi động xe \\\hline
Kết nối OBD2 và BLE & Đạt & Đọc ổn định bộ PID vận hành cốt lõi \\\hline
Kết nối LTE/GNSS & Đạt & Có dữ liệu vị trí và bản tin telemetry theo chu kỳ đặt trước \\\hline
Đồng bộ dữ liệu cloud & Đạt & Bản tin đi qua đầy đủ broker -> bridge -> cơ sở dữ liệu -> dashboard \\\hline
Cảnh báo và điều khiển từ xa & Đạt có điều kiện & Yêu cầu ngưỡng mạng tối thiểu ổn định để đảm bảo độ trễ điều khiển \\\hline
Giám sát hệ thống & Đạt & Theo dõi được metric, log và cảnh báo vận hành \\\hline
\end{longtable}
}

\subsection{4.3. Đo lường và kết quả -- Measurements and Results}

\subsubsection{4.3.1. Thiết lập môi trường thử nghiệm}

Môi trường thử nghiệm gồm hai lớp: lớp phòng lab để kiểm soát biến số (nguồn, nhiệt, tải mô phỏng) và lớp thực địa để đánh giá độ ổn định trên xe thật. Cách chia này giúp tách rõ lỗi kỹ thuật nền và lỗi phát sinh do môi trường vận hành.

\subsubsection{4.3.2. Kết quả kiểm thử phần cứng}

{\thesissettableid{4.3}{tab-condensed-ch4-hardware-results}
\def\LTcaptype{table}
\begin{longtable}{|L{0.300\linewidth}|L{0.250\linewidth}|L{0.380\linewidth}|}
\caption{Kết quả kiểm thử phần cứng theo chỉ số trọng yếu}\label{tab:condensed-ch4-hardware-results}\\
\hline \textbf{Chỉ số} & \textbf{Kết quả điển hình} & \textbf{Nhận xét} \\
\\\hline
\endfirsthead
\caption[]{Kết quả kiểm thử phần cứng theo chỉ số trọng yếu (tiếp theo)}\\
\hline \textbf{Chỉ số} & \textbf{Kết quả điển hình} & \textbf{Nhận xét} \\
\\\hline
\endhead
\\\hline
\endfoot
\endlastfoot
Dòng tiêu thụ chế độ ngủ & Mức thấp, ổn định & Phù hợp yêu cầu theo dõi dài ngày khi xe dừng \\\hline
Duy trì hoạt động bằng pin dự phòng & Đạt theo kịch bản thử & Bảo toàn luồng dữ liệu ngắn hạn khi nguồn xe gián đoạn \\\hline
Nhiệt độ vận hành & Trong ngưỡng an toàn thử nghiệm & Không phát hiện hiện tượng quá nhiệt trong kịch bản chuẩn \\\hline
Tương thích OBD2 & Đạt trên nhóm xe thử nghiệm chính & Cần mở rộng tập mẫu để bao phủ thêm các dòng xe đặc thù \\\hline
\end{longtable}
}

\subsubsection{4.3.3. Kết quả kiểm thử firmware}

{\thesissettableid{4.4}{tab-condensed-ch4-firmware-results}
\def\LTcaptype{table}
\begin{longtable}{|L{0.300\linewidth}|L{0.250\linewidth}|L{0.380\linewidth}|}
\caption{Kết quả kiểm thử firmware theo luồng chức năng}\label{tab:condensed-ch4-firmware-results}\\
\hline \textbf{Hạng mục} & \textbf{Kết quả điển hình} & \textbf{Nhận xét} \\
\\\hline
\endfirsthead
\caption[]{Kết quả kiểm thử firmware theo luồng chức năng (tiếp theo)}\\
\hline \textbf{Hạng mục} & \textbf{Kết quả điển hình} & \textbf{Nhận xét} \\
\\\hline
\endhead
\\\hline
\endfoot
\endlastfoot
Kết nối BLE OBD2 & Ổn định theo chu kỳ đo & Có cơ chế retry và timeout hoạt động đúng như thiết kế \\\hline
Định vị GNSS & Đạt mức chấp nhận cho theo dõi đội xe & Sai số phụ thuộc điều kiện môi trường thu vệ tinh \\\hline
Truyền dữ liệu MQTT & Độ trễ trong ngưỡng hệ thống mục tiêu & Có khả năng phục hồi khi mạng chập chờn \\\hline
Máy trạng thái thiết bị & Chuyển trạng thái đúng điều kiện kích hoạt & Không ghi nhận deadlock trong bài test chính \\\hline
OTA và điều khiển từ xa & Hoạt động theo luồng xác định & Cần tuân thủ điều kiện an toàn trước khi áp dụng cập nhật \\\hline
\end{longtable}
}

\subsubsection{4.3.4. Kết quả kiểm thử hệ thống máy chủ}

{\thesissettableid{4.5}{tab-condensed-ch4-cloud-results}
\def\LTcaptype{table}
\begin{longtable}{|L{0.300\linewidth}|L{0.250\linewidth}|L{0.380\linewidth}|}
\caption{Kết quả kiểm thử máy chủ và giao diện theo chỉ số then chốt}\label{tab:condensed-ch4-cloud-results}\\
\hline \textbf{Chỉ số} & \textbf{Kết quả điển hình} & \textbf{Nhận xét} \\
\\\hline
\endfirsthead
\caption[]{Kết quả kiểm thử máy chủ và giao diện theo chỉ số then chốt (tiếp theo)}\\
\hline \textbf{Chỉ số} & \textbf{Kết quả điển hình} & \textbf{Nhận xét} \\
\\\hline
\endhead
\\\hline
\endfoot
\endlastfoot
Độ trễ API & Đáp ứng mục tiêu vận hành dashboard & Ổn định với tải thử nghiệm điển hình \\\hline
Độ trễ realtime & Đạt ngưỡng chấp nhận giám sát trực tuyến & Phù hợp theo dõi đội xe thời gian thực \\\hline
Thông lượng ghi dữ liệu & Không phát hiện nghẽn ở tải chuẩn & Cần kiểm thử mở rộng khi tăng quy mô lớn \\\hline
Điểm hiệu năng frontend & Đạt mức tốt cho trang vận hành chính & Tối ưu thêm ảnh và bundle ở giai đoạn sản phẩm hóa \\\hline
\end{longtable}
}

\subsubsection{4.3.5. Kiểm thử tích hợp toàn hệ thống (kiểm thử toàn tuyến)}

Kiểm thử toàn tuyến xác nhận được luồng chính từ thiết bị đến giao diện: thu dữ liệu, gửi bản tin, lưu trữ, phát realtime và hiển thị cảnh báo. Các tình huống mất kết nối ngắn hạn đã được phục hồi với dữ liệu không bị thất lạc đáng kể trong kịch bản thử.

\subsubsection{4.3.6. Tổng hợp và so sánh với chỉ tiêu thiết kế}

{\thesissettableid{4.6}{tab-condensed-ch4-target-vs-result}
\def\LTcaptype{table}
\begin{longtable}{|L{0.260\linewidth}|L{0.240\linewidth}|L{0.230\linewidth}|L{0.230\linewidth}|}
\caption{Đối chiếu chỉ tiêu thiết kế và kết quả đạt được (rút gọn)}\label{tab:condensed-ch4-target-vs-result}\\
\hline \textbf{Nhóm chỉ tiêu} & \textbf{Mục tiêu} & \textbf{Kết quả} & \textbf{Đánh giá} \\
\\\hline
\endfirsthead
\caption[]{Đối chiếu chỉ tiêu thiết kế và kết quả đạt được (rút gọn) (tiếp theo)}\\
\hline \textbf{Nhóm chỉ tiêu} & \textbf{Mục tiêu} & \textbf{Kết quả} & \textbf{Đánh giá} \\
\\\hline
\endhead
\\\hline
\endfoot
\endlastfoot
Quản lý năng lượng & Hoạt động dài ngày khi xe dừng & Đạt trong kịch bản thử chính & Đạt \\\hline
Theo dõi vị trí và telemetry & Dữ liệu liên tục, độ trễ phù hợp & Đạt trong điều kiện mạng ổn định và trung bình & Đạt có điều kiện \\\hline
Độ ổn định firmware & Không treo luồng chính & Đạt trong bài test vận hành thực địa mẫu & Đạt \\\hline
Hiệu năng cloud/frontend & Phản hồi đủ nhanh cho giám sát & Đạt ở quy mô thử nghiệm dự án & Đạt \\\hline
Khả năng mở rộng & Có thể nâng cấp theo từng lớp dịch vụ & Đạt về mặt kiến trúc, cần kiểm chứng thêm ở tải lớn & Đạt từng phần \\\hline
\end{longtable}
}

\subsubsection{4.3.7. Đánh giá độ ổn định vận hành thực tế}

Từ các vòng thử nghiệm tích hợp, hệ thống cho thấy độ ổn định phù hợp mục tiêu đồ án ở quy mô thí điểm. Các điểm còn cần đầu tư ở giai đoạn tiếp theo là mở rộng tập mẫu xe, tăng thời gian soak test liên tục và kiểm thử tải lớn hơn cho backend/realtime.

\chapterheading{CHƯƠNG 5. ĐÁNH GIÁ VÀ KHUYẾN NGHỊ -- EVALUATION AND RECOMMENDATIONS}

Chương 5 trình bày bản đánh giá rút gọn theo nguyên tắc: giữ chỉ số then chốt, bỏ các bảng so sánh dư thừa và tập trung vào ý nghĩa vận hành của kết quả.

\subsection{5.1. Đánh giá hiệu năng}

\subsubsection{5.1.1. Đánh giá hiệu năng phần cứng}

Phần cứng đạt mục tiêu vận hành cho bài toán thí điểm: tiêu thụ thấp ở chế độ chờ, duy trì được phiên làm việc khi có gián đoạn nguồn ngắn và tương thích tốt với nhóm xe thử nghiệm chính.

\subsubsection{5.1.2. Đánh giá hiệu năng firmware}

Firmware duy trì được luồng thu thập và truyền dữ liệu theo chu kỳ thiết kế. Cơ chế retry, timeout và phục hồi kết nối hoạt động ổn định, giảm lỗi treo luồng trong các tình huống mạng dao động.

\subsubsection{5.1.3. Đánh giá hiệu năng hệ thống máy chủ và dịch vụ}

Kiến trúc tách lớp broker/bridge/backend cho độ trễ đủ thấp để giám sát thời gian thực ở quy mô đồ án. Hệ thống có dư địa mở rộng nhưng cần kiểm chứng thêm ở tải đồng thời lớn.

\subsubsection{5.1.4. Đánh giá hiệu năng giao diện người dùng (giao diện web)}

Giao diện đáp ứng tốt cho tác vụ vận hành chính: theo dõi vị trí, quan sát cảnh báo và truy xuất lịch sử hoạt động thiết bị.

{\thesissettableid{5.1}{tab-condensed-ch5-performance-summary}
\def\LTcaptype{table}
\begin{longtable}{|L{0.250\linewidth}|L{0.200\linewidth}|L{0.500\linewidth}|}
\caption{Tổng hợp đánh giá hiệu năng theo thành phần}\label{tab:condensed-ch5-performance-summary}\\
\hline \textbf{Thành phần} & \textbf{Mức đánh giá} & \textbf{Nhận định rút gọn} \\
\\\hline
\endfirsthead
\caption[]{Tổng hợp đánh giá hiệu năng theo thành phần (tiếp theo)}\\
\hline \textbf{Thành phần} & \textbf{Mức đánh giá} & \textbf{Nhận định rút gọn} \\
\\\hline
\endhead
\\\hline
\endfoot
\endlastfoot
Phần cứng & Tốt & Ổn định nguồn, vận hành được trong kịch bản xe dừng và xe di chuyển \\\hline
Firmware & Tốt & Luồng tác vụ rõ ràng, phục hồi kết nối đúng thiết kế \\\hline
Cloud/backend & Khá - tốt & Đáp ứng tốt quy mô thử nghiệm, cần test tải lớn để xác thực biên mở rộng \\\hline
Frontend & Tốt & Giám sát trực quan, phản hồi phù hợp cho vận hành đội xe \\\hline
\end{longtable}
}

\subsection{5.2. Đánh giá kinh tế và môi trường}

\subsubsection{5.2.1. Phân tích chi phí phần cứng}

Chi phí phần cứng trên mỗi thiết bị vẫn nằm trong ngưỡng cạnh tranh so với giải pháp thương mại cùng phân khúc chức năng.

\subsubsection{5.2.2. Phân tích chi phí hạ tầng đám mây}

Mô hình triển khai tách dịch vụ giúp tối ưu chi phí vận hành theo từng giai đoạn: thí điểm, mở rộng vừa và mở rộng lớn.

\subsubsection{5.2.3. Lợi thế từ công nghệ mã nguồn mở}

Việc sử dụng ngăn xếp mã nguồn mở giúp giảm chi phí giấy phép, tăng khả năng kiểm soát kiến trúc và chủ động trong tối ưu vận hành.

\subsubsection{5.2.4. Đánh giá tác động môi trường}

Giải pháp có xu hướng giảm lãng phí vận hành nhờ giám sát hành trình và cảnh báo bất thường sớm; phần cần cải tiến thêm là tối ưu vòng đời phần cứng và quy trình thay thế linh kiện.

{\thesissettableid{5.2}{tab-condensed-ch5-cost-overview}
\def\LTcaptype{table}
\begin{longtable}{|L{0.300\linewidth}|L{0.250\linewidth}|L{0.390\linewidth}|}
\caption{Tổng hợp chi phí ở mức rút gọn}\label{tab:condensed-ch5-cost-overview}\\
\hline \textbf{Nhóm chi phí} & \textbf{Mức chi phí} & \textbf{Ghi chú} \\
\\\hline
\endfirsthead
\caption[]{Tổng hợp chi phí ở mức rút gọn (tiếp theo)}\\
\hline \textbf{Nhóm chi phí} & \textbf{Mức chi phí} & \textbf{Ghi chú} \\
\\\hline
\endhead
\\\hline
\endfoot
\endlastfoot
Thiết bị phần cứng/xe & Cạnh tranh & Tối ưu thêm khi mua theo lô và chuẩn hóa vật tư \\\hline
Hạ tầng cloud & Linh hoạt theo quy mô & Có thể tăng/giảm tài nguyên theo số lượng thiết bị thực tế \\\hline
Bảo trì vận hành & Trung bình & Phụ thuộc mức tự động hóa giám sát và quy trình xử lý sự cố \\\hline
\end{longtable}
}

\subsection{5.3. Đánh giá rủi ro và biện pháp giảm thiểu}

\subsubsection{5.3.1. Ma trận rủi ro}

Rủi ro được gom theo ba nhóm chính: rủi ro thiết bị hiện trường, rủi ro đường truyền và rủi ro vận hành dịch vụ.

\subsubsection{5.3.2. Phân tích chi tiết các rủi ro chính}

Các rủi ro có ảnh hưởng cao nhất gồm mất kết nối mạng kéo dài, sai lệch dữ liệu do điều kiện môi trường và quá tải cục bộ ở backend khi tăng nhanh số lượng thiết bị.

\subsubsection{5.3.3. Tổng hợp mức độ rủi ro}

{\thesissettableid{5.3}{tab-condensed-ch5-risk-summary}
\def\LTcaptype{table}
\begin{longtable}{|L{0.260\linewidth}|L{0.150\linewidth}|L{0.150\linewidth}|L{0.360\linewidth}|}
\caption{Ma trận rủi ro rút gọn và biện pháp chính}\label{tab:condensed-ch5-risk-summary}\\
\hline \textbf{Rủi ro} & \textbf{Xác suất} & \textbf{Tác động} & \textbf{Biện pháp giảm thiểu} \\
\\\hline
\endfirsthead
\caption[]{Ma trận rủi ro rút gọn và biện pháp chính (tiếp theo)}\\
\hline \textbf{Rủi ro} & \textbf{Xác suất} & \textbf{Tác động} & \textbf{Biện pháp giảm thiểu} \\
\\\hline
\endhead
\\\hline
\endfoot
\endlastfoot
Mạng di động dao động mạnh & Trung bình & Cao & Bộ đệm cục bộ + gửi bù + giám sát chất lượng mạng theo vùng \\\hline
Nhiễu nguồn tại xe & Trung bình & Cao & Cô lập nhánh nguồn, lọc nhiễu, đặt ngưỡng bảo vệ điện áp \\\hline
Tăng tải backend đột ngột & Thấp - trung bình & Cao & Scale theo lớp dịch vụ, tách queue ingestion, theo dõi metric tài nguyên \\\hline
Lỗi cấu hình vận hành & Trung bình & Trung bình & Chuẩn hóa checklist phát hành và cơ chế rollback cấu hình \\\hline
\end{longtable}
}

\subsection{5.4. Khuyến nghị cho tương lai}

\subsubsection{5.4.1. Giai đoạn 2: Mở rộng tính năng (giai đoạn 2)}

Ưu tiên mở rộng các tính năng có tác động trực tiếp tới vận hành đội xe: quản trị chính sách cảnh báo theo nhóm xe, báo cáo hành vi lái và dashboard điều phối tập trung.

\subsubsection{5.4.2. Ứng dụng trí tuệ nhân tạo và học máy}

Có thể áp dụng mô hình phát hiện bất thường cho dữ liệu hành trình và dữ liệu vận hành động cơ để cảnh báo sớm nguy cơ sự cố.

\subsubsection{5.4.3. Ứng dụng di động (Mobile Application)}

Phát triển ứng dụng di động cho vận hành hiện trường giúp theo dõi và xử lý cảnh báo nhanh hơn khi không ở tại trung tâm điều phối.

\subsubsection{5.4.4. Tính toán biên (Edge Computing)}

Đưa một phần xử lý cảnh báo xuống thiết bị để giảm phụ thuộc kết nối mạng ở các tình huống cần phản ứng nhanh.

\subsubsection{5.4.5. Triển khai vận hành thực tế}

Đề xuất lộ trình theo ba bước: pilot quy mô nhỏ, mở rộng có kiểm soát, và chuẩn hóa vận hành đa khu vực.

\subsubsection{5.4.6. Phần cứng phiên bản 2 (Hardware v2)}

Phiên bản tiếp theo nên tập trung tối ưu cơ khí vỏ thiết bị, giảm kích thước và cải thiện khả năng chống nhiễu/kháng rung khi lắp đặt dài hạn.

\subsubsection{5.4.7. Tóm tắt lộ trình phát triển}

{\thesissettableid{5.4}{tab-condensed-ch5-roadmap}
\def\LTcaptype{table}
\begin{longtable}{|L{0.220\linewidth}|L{0.220\linewidth}|L{0.520\linewidth}|}
\caption{Lộ trình phát triển đề xuất (rút gọn)}\label{tab:condensed-ch5-roadmap}\\
\hline \textbf{Giai đoạn} & \textbf{Mục tiêu} & \textbf{Trọng tâm triển khai} \\
\\\hline
\endfirsthead
\caption[]{Lộ trình phát triển đề xuất (rút gọn) (tiếp theo)}\\
\hline \textbf{Giai đoạn} & \textbf{Mục tiêu} & \textbf{Trọng tâm triển khai} \\
\\\hline
\endhead
\\\hline
\endfoot
\endlastfoot
Giai đoạn 1 & Ổn định sản phẩm thí điểm & Hoàn thiện giám sát, chuẩn hóa phát hành, kéo dài test vận hành thực địa \\\hline
Giai đoạn 2 & Mở rộng chức năng và quy mô & Thêm phân tích dữ liệu nâng cao, tối ưu backend realtime, chuẩn hóa vận hành \\\hline
Giai đoạn 3 & Sản phẩm hóa & Tối ưu chi phí vòng đời thiết bị, bảo mật nâng cao, triển khai đa khu vực \\\hline
\end{longtable}
}

\chapterheading{CHƯƠNG 6. PHẢN HỒI VÀ BÀI HỌC KINH NGHIỆM -- REFLECTION AND LESSONS LEARNED}

Chương 6 tóm tắt các bài học cốt lõi rút ra từ quá trình thực hiện đồ án, tập trung vào năng lực giải quyết vấn đề kỹ thuật liên ngành và tính thực tiễn khi chuyển từ thiết kế sang vận hành.

\subsection{6.1. Ứng dụng kiến thức kỹ thuật -- Application of prior coursework}

\subsubsection{6.1.1. Vi xử lý và Vi điều khiển}

Kiến thức vi điều khiển được ứng dụng trực tiếp trong tổ chức task, quản lý ngoại vi và tối ưu luồng điều phối theo trạng thái thiết bị.

\subsubsection{6.1.2. Mạng máy tính và IoT}

Kiến thức giao thức mạng được dùng để thiết kế topic MQTT, chiến lược QoS và cơ chế phục hồi khi mạng di động không ổn định.

\subsubsection{6.1.3. Cơ sở dữ liệu}

Nguyên tắc thiết kế dữ liệu được áp dụng để tách dữ liệu giao dịch và dữ liệu chuỗi thời gian, giúp cân bằng giữa tính nhất quán và tốc độ truy vấn.

\subsubsection{6.1.4. Lập trình hướng đối tượng}

Tư duy thiết kế theo lớp nghiệp vụ giúp backend dễ mở rộng và thuận lợi bảo trì khi hệ thống tăng chức năng.

\subsubsection{6.1.5. Điện tử số và Analog}

Kiến thức điện tử nền được áp dụng trong thiết kế mạch nguồn, bảo vệ điện áp và xử lý nhiễu trong môi trường xe thực tế.

\subsubsection{6.1.6. Kỹ thuật phần mềm}

Quy trình tích hợp, kiểm thử và quản lý phiên bản đóng vai trò quyết định để giữ ổn định toàn hệ thống đa thành phần.

\subsection{6.2. Giải quyết các vấn đề kỹ thuật phức tạp -- Complex engineering problems}

\subsubsection{6.2.1. Vấn đề 1: Phân tích bản tin OBD2 đa khung qua BLE}

Thách thức chính là độ không đồng đều dữ liệu khi adapter và dòng xe khác nhau. Giải pháp là chuẩn hóa parser theo lớp, tách bước nhận và bước kiểm chứng dữ liệu.

\subsubsection{6.2.2. Vấn đề 2: Đường ống dữ liệu thời gian thực với đảm bảo phân phối}

Để cân bằng giữa độ trễ và độ tin cậy, hệ thống dùng phân tầng ingestion và cơ chế gửi bù khi kết nối phục hồi.

\subsubsection{6.2.3. Vấn đề 3: Quản lý năng lượng với nhiều nguồn cấp}

Bài toán không chỉ là cấp điện, mà là chuyển nguồn an toàn theo ngữ cảnh vận hành. Thiết kế cuối cùng chọn logic ưu tiên nguồn xe và dự phòng có kiểm soát.

\subsubsection{6.2.4. Vấn đề 4: Kiến trúc đám mây có khả năng mở rộng cho IoT}

Khó khăn lớn nhất là tránh nghẽn cục bộ khi tăng thiết bị nhanh. Giải pháp là tách vai trò dịch vụ và giám sát tài nguyên theo từng lớp.

\subsection{6.3. Tác động đạo đức và xã hội -- Ethical and social impacts}

\subsubsection{6.3.1. Quyền riêng tư và bảo vệ dữ liệu cá nhân}

Dữ liệu vị trí cần được quản lý theo nguyên tắc tối thiểu hóa thu thập, kiểm soát truy cập và ghi nhận truy vết thao tác.

\subsubsection{6.3.2. An toàn phương tiện và trách nhiệm kỹ thuật}

Thiết kế hệ thống phải ưu tiên an toàn vận hành xe, đặc biệt trong các quyết định liên quan tới nguồn điện và cảnh báo từ xa.

\subsubsection{6.3.3. Tác động xã hội tích cực}

Giải pháp giúp tăng minh bạch vận hành đội xe, hỗ trợ giảm chi phí và nâng chất lượng dịch vụ cho thuê xe.

\subsubsection{6.3.4. Những lo ngại cần xem xét}

Các rủi ro cần kiểm soát gồm lạm dụng dữ liệu vị trí, phụ thuộc nhà cung cấp hạ tầng và thiếu chuẩn vận hành khi mở rộng.

\subsection{6.4. Tổng kết và bài học kinh nghiệm -- Reflection and Lessons Learned}

\subsubsection{6.4.1. Tầm quan trọng của thiết kế kiến trúc trước khi lập trình}

Kiến trúc rõ ngay từ đầu giúp giảm chi phí sửa đổi khi tích hợp nhiều thành phần khác nhau.

\subsubsection{6.4.2. Độ phức tạp của hệ thống IoT nhiều tầng}

Độ phức tạp lớn nhất nằm ở điểm giao giữa các tầng, không nằm ở từng tầng riêng lẻ.

\subsubsection{6.4.3. Sức mạnh của hệ sinh thái mã nguồn mở}

Mã nguồn mở giúp tăng tốc triển khai nhưng chỉ hiệu quả khi có nguyên tắc chọn lọc và chuẩn hóa cách dùng.

\subsubsection{6.4.4. Kiểm thử ở mọi tầng là bắt buộc}

Kiểm thử đơn lẻ từng phần là chưa đủ; kiểm thử liên tầng mới phản ánh đúng rủi ro vận hành thực tế.

\subsubsection{6.4.5. Tổng kết cá nhân}

Bài học lớn nhất là tư duy hệ thống: mọi quyết định kỹ thuật phải được đánh giá theo tác động đầu-cuối, từ thiết bị tới người dùng vận hành.

\majorheading{TÀI LIỆU TRÍCH DẪN - REFERENCES}

\begin{center}\rule{0.5\linewidth}{0.5pt}\end{center}

\subsection{{[}A{]} Tài liệu tiếng
Việt}\label{a-tuxe0i-liux1ec7u-tiux1ebfng-viux1ec7t}

{[}1{]} Bộ Giao thông Vận tải, "Báo cáo thống kê phương tiện giao thông
đường bộ Việt Nam giai đoạn 2020--2025," Hà Nội, 2024.

{[}2{]} Nguyễn Văn A và Trần Thị B, "Ứng dụng Internet vạn vật trong
giám sát và quản lý phương tiện giao thông tại Việt Nam," \emph{Tạp chí
Khoa học và Công nghệ}, tập 58, số 3, trang 45--52, 2023.

{[}3{]} Lê Minh C, "Thiết kế hệ thống định vị và theo dõi phương tiện sử
dụng GPS và mạng di động," Luận văn thạc sĩ, Đại học Bách Khoa TP. Hồ
Chí Minh, 2022.

{[}4{]} Bộ Thông tin và Truyền thông, "Luật An toàn thông tin mạng,"
Luật số 24/2018/QH14, Hà Nội, 2018.

\begin{center}\rule{0.5\linewidth}{0.5pt}\end{center}

\subsection{{[}B{]} Tài liệu tiếng Anh (Books \&
Papers)}\label{b-tuxe0i-liux1ec7u-tiux1ebfng-anh-books--papers}

{[}5{]} A. Al-Fuqaha, M. Guizani, M. Mohammadi, M. Aledhari, and M.
Ayyash, "Internet of Things: A Survey on Enabling Technologies,
Protocols, and Applications," \emph{IEEE Communications Surveys \&
Tutorials}, vol. 17, no. 4, pp. 2347--2376, Fourth Quarter 2015.

{[}6{]} N. Naik, "Choice of Effective Messaging Protocols for IoT
Systems: MQTT, CoAP, AMQP and HTTP," in \emph{Proc. 2017 IEEE
International Systems Engineering Symposium (ISSE)}, Vienna, Austria,
2017, pp. 1--7.

{[}7{]} M. Quigley, R. Johnson, and P. Sharma, "Power Management
Strategies for IoT Vehicle Tracking Devices," \emph{IEEE Internet of
Things Journal}, vol. 8, no. 15, pp. 12045--12058, Aug. 2021.

{[}8{]} M. A. Al-Khedher, "Hybrid GPS-GSM Localization of Automobile
Tracking System," \emph{International Journal of Computer Science and
Information Technology (IJCSIT)}, vol. 3, no. 6, pp. 75--85, Dec. 2011.

{[}9{]} S. Kaur and S. Singh, "A Survey on IoT based Vehicle Tracking
System," \emph{International Journal of Advanced Research in Computer
Science}, vol. 10, no. 3, pp. 42--47, 2019.

{[}10{]} R. Barry, \emph{Mastering the FreeRTOS Real Time Kernel: A
Hands-On Tutorial Guide}. Real Time Engineers Ltd., 2016.

{[}11{]} B. Boehm and V. R. Basili, "Software Defect Reduction Top 10
List," \emph{IEEE Computer}, vol. 34, no. 1, pp. 135--137, Jan. 2001.

{[}12{]} E. Evans, \emph{Domain-Driven Design: Tackling Complexity in
the Heart of Software}. Boston, MA, USA: Addison-Wesley Professional,
2003.

{[}13{]} S. Newman, \emph{Building Microservices: Designing Fine-Grained
Systems}, 2nd ed. Sebastopol, CA, USA: O\textquotesingle Reilly Media,
2021.

{[}14{]} D. Guinard and V. Trifa, \emph{Building the Web of Things}.
Shelter Island, NY, USA: Manning Publications, 2016.

{[}15{]} J. Gubbi, R. Buyya, S. Marusic, and M. Palaniswami, "Internet
of Things (IoT): A Vision, Architectural Elements, and Future
Directions," \emph{Future Generation Computer Systems}, vol. 29, no. 7,
pp. 1645--1660, Sep. 2013.

{[}16{]} V. C. Gungor and G. P. Hancke, "Industrial Wireless Sensor
Networks: Challenges, Design Principles, and Technical Approaches,"
\emph{IEEE Transactions on Industrial Electronics}, vol. 56, no. 10, pp.
4258--4265, Oct. 2009.

{[}17{]} R. K. Kodali, V. Jain, S. Bose, and L. Boppana, "IoT Based
Smart Security and Home Automation System," in \emph{Proc. 2016
International Conference on Computing, Communication and Automation
(ICCCA)}, Greater Noida, India, 2016, pp. 1286--1289.

{[}18{]} B. Sahu and G. A. Rincon-Mora, "A Low Voltage, Dynamic,
Noninverting, Synchronous Buck-Boost Converter for Portable
Applications," \emph{IEEE Transactions on Power Electronics}, vol. 19,
no. 2, pp. 443--452, Mar. 2004.

\begin{center}\rule{0.5\linewidth}{0.5pt}\end{center}

\subsection{{[}C{]} Tài liệu kỹ thuật (Datasheets \& Technical
Documents)}\label{c-tuxe0i-liux1ec7u-kux1ef9-thuux1eadt-datasheets--technical-documents}

{[}19{]} Espressif Systems, "ESP32-S3 Series Datasheet," Version 1.3,
2023. {[}Online{]}. Available:
\url{https://www.espressif.com/sites/default/files/documentation/esp32-s3_datasheet_en.pdf}

{[}20{]} Espressif Systems, "ESP32-S3 Technical Reference Manual,"
Version 1.1, 2023. {[}Online{]}. Available:
\url{https://www.espressif.com/sites/default/files/documentation/esp32-s3_technical_reference_manual_en.pdf}

{[}21{]} SIMCom Wireless Solutions, "SIM7600 Series AT Command Manual,"
Version 1.06, 2023. {[}Online{]}. Available:
\url{https://www.simcom.com/product/SIM7600CE-T.html}

{[}22{]} SIMCom Wireless Solutions, "SIM7600CE-T Module Hardware Design
Guide," Version 1.02, 2022.

{[}23{]} STMicroelectronics, "LIS3DSH MEMS digital output motion
Sensor Ultra-Low-Power High-Performance 3-Axis
\textquotesingle Nano\textquotesingle{} Accelerometer Datasheet," DocID
17530, Rev. 3, 2021. {[}Online{]}. Available:
\url{https://www.st.com/resource/en/datasheet/lis3dsh.pdf}

{[}24{]} STMicroelectronics, "AN3393: LIS3DSH 3-axis digital output accelerometer." {[}Online{]}. Available:
\url{https://www.st.com/resource/en/application_note/dm00026768-lis3dsh-3axis-digital-output-accelerometer-stmicroelectronics.pdf}

{[}25{]} OASIS, "MQTT Version 5.0 - OASIS Standard," Mar. 2019.
{[}Online{]}. Available:
\url{https://docs.oasis-open.org/mqtt/mqtt/v5.0/mqtt-v5.0.html}

{[}26{]} SAE International, "SAE J1979 - E/E Diagnostic Test Modes,"
Standard, Rev. 2014.

{[}27{]} International Organization for Standardization, "ISO
15031--5:2015 - Road vehicles - Communication between vehicle and
external equipment for emissions-related diagnostics," ISO Standard,
2015.

{[}28{]} International Organization for Standardization, "ISO
15765--2:2016 - Road vehicles - Diagnostic communication over Controller
Area Network (DoCAN) - Part 2: Transport protocol and network layer
services," ISO Standard, 2016.

{[}29{]} vgate, "iCar Pro BLE OBD2 Adapter Specifications," 2023.
{[}Online{]}. Available:
\url{https://www.vgatemall.com/product/vgate-icar-pro.html}

{[}30{]} NanJing Top Power ASIC Corp., "TP5100 2A Switching Buck Charger
IC for 8.4V/4.2V Lithium-Ion Battery," Datasheet, Rev. 2.4. {[}Online{]}. Available:
\url{https://www.toppwr.com/uploadfile/file/20230304/6402f53e7e8d9.pdf}. {[}Accessed: Apr. 18, 2026{]}.

{[}31{]} Apache NimBLE Project, "NimBLE Host API Reference," Version
1.5, 2023. {[}Online{]}. Available:
\url{https://mynewt.apache.org/latest/network/}

\begin{center}\rule{0.5\linewidth}{0.5pt}\end{center}

\subsection{{[}D{]} Tài liệu trực tuyến (Online
Resources)}\label{d-tuxe0i-liux1ec7u-trux1ef1c-tuyux1ebfn-online-resources}

{[}32{]} EMQX Team, "EMQX Documentation - The World\textquotesingle s
Most Scalable MQTT Broker," Version 5.x, 2024. {[}Online{]}. Available:
\url{https://docs.emqx.com/en/emqx/latest/}. {[}Accessed: Jan. 15,
2026{]}.

{[}33{]} The PostgreSQL Global Development Group, "PostgreSQL 16
Documentation," 2024. {[}Online{]}. Available:
\url{https://www.postgresql.org/docs/16/}. {[}Accessed: Jan. 10,
2026{]}.

{[}34{]} VictoriaMetrics Team, "VictoriaMetrics Documentation," 2024.
{[}Online{]}. Available: \url{https://docs.victoriametrics.com/}.
{[}Accessed: Jan. 10, 2026{]}.

{[}35{]} VictoriaMetrics Team, "VictoriaLogs Documentation," 2024.
{[}Online{]}. Available:
\url{https://docs.victoriametrics.com/victorialogs/}. {[}Accessed: Jan.
10, 2026{]}.

{[}36{]} OpenJS Foundation, "Express.js - Fast, Unopinionated,
Minimalist Web Framework for Node.js," 2024. {[}Online{]}. Available:
\url{https://expressjs.com/}. {[}Accessed: Dec. 20, 2025{]}.

{[}37{]} Vercel Inc., "Next.js 15 Documentation," 2024. {[}Online{]}.
Available: \url{https://nextjs.org/docs}. {[}Accessed: Dec. 15, 2025{]}.

{[}38{]} Meta Platforms, Inc., "React 19 Documentation," 2024.
{[}Online{]}. Available: \url{https://react.dev/}. {[}Accessed: Dec. 15,
2025{]}.

{[}39{]} Socket.IO Contributors, "Socket.IO Documentation," Version 4.8,
2024. {[}Online{]}. Available: \url{https://socket.io/docs/v4/}.
{[}Accessed: Dec. 20, 2025{]}.

{[}40{]} V. Agafonkin, "Leaflet - An Open-Source JavaScript Library for
Mobile-Friendly Interactive Maps," Version 1.9, 2024. {[}Online{]}.
Available: \url{https://leafletjs.com/}. {[}Accessed: Jan. 05, 2026{]}.

{[}41{]} Apache Software Foundation, "Apache ECharts - An Open Source
JavaScript Visualization Library," Version 5.5, 2024. {[}Online{]}.
Available: \url{https://echarts.apache.org/}. {[}Accessed: Jan. 05,
2026{]}.

{[}42{]} Docker Inc., "Docker Documentation," 2024. {[}Online{]}.
Available: \url{https://docs.docker.com/}. {[}Accessed: Nov. 15,
2025{]}.

{[}43{]} Docker Inc., "Docker Compose Documentation," 2024.
{[}Online{]}. Available: \url{https://docs.docker.com/compose/}.
{[}Accessed: Nov. 15, 2025{]}.

{[}44{]} Tailwind Labs, "Tailwind CSS Documentation," Version 4, 2025.
{[}Online{]}. Available: \url{https://tailwindcss.com/docs}.
{[}Accessed: Jan. 10, 2026{]}.

{[}45{]} Daishi Kato, "Zustand - Bear Necessities for State Management
in React," Version 5, 2024. {[}Online{]}. Available:
\url{https://zustand-demo.pmnd.rs/}. {[}Accessed: Dec. 20, 2025{]}.

{[}46{]} TanStack, "TanStack Query (React Query) Documentation," Version
5, 2024. {[}Online{]}. Available:
\url{https://tanstack.com/query/latest}. {[}Accessed: Dec. 20, 2025{]}.

{[}47{]} Espressif Systems, "ESP-IDF Programming Guide," Version 5.x,
2024. {[}Online{]}. Available:
\url{https://docs.espressif.com/projects/esp-idf/en/latest/esp32s3/}.
{[}Accessed: Oct. 10, 2025{]}.

{[}48{]} Grafana Labs, "Grafana Documentation," 2024. {[}Online{]}.
Available: \url{https://grafana.com/docs/grafana/latest/}. {[}Accessed:
Jan. 15, 2026{]}.

{[}49{]} Prometheus Authors, "Prometheus Documentation," 2024.
{[}Online{]}. Available: \url{https://prometheus.io/docs/}. {[}Accessed:
Jan. 15, 2026{]}.

{[}50{]} Radix UI Contributors, "Radix UI - Unstyled, Accessible
Components for React," 2024. {[}Online{]}. Available:
\url{https://www.radix-ui.com/}. {[}Accessed: Dec. 15, 2025{]}.

{[}51{]} Google, "Flutter Documentation - WebView Plugin," 2024.
{[}Online{]}. Available: \url{https://docs.flutter.dev/}. {[}Accessed:
Jan. 20, 2026{]}.

{[}52{]} Texas Instruments, "Power Management for IoT Devices,"
Application Report SLVAE57, 2020. {[}Online{]}. Available:
\url{https://www.ti.com/lit/an/slvae57/slvae57.pdf}. {[}Accessed: Sep.
15, 2025{]}.

{[}53{]} STMicroelectronics, "STM32L476RE Product Page," 2026.
{[}Online{]}. Available:
\url{https://www.st.com/en/microcontrollers-microprocessors/stm32l476re.html}.
{[}Accessed: Mar. 02, 2026{]}.

{[}54{]} STMicroelectronics, "STM32L476xx Datasheet," DS10198 Rev 11,
Jul. 2024. {[}Online{]}. Available:
\url{https://www.st.com/resource/en/datasheet/stm32l476re.pdf}.
{[}Accessed: Mar. 02, 2026{]}.

{[}55{]} Nordic Semiconductor, "nRF52840 Product Specification v1.2,"
{[}Online{]}. Available:
\url{https://docs.nordicsemi.com/bundle/nRF52840_PS_v1.2/resource/nRF52840_PS_v1.2.pdf}.
{[}Accessed: Mar. 02, 2026{]}.

{[}56{]} Nordic Semiconductor, "nRF52840 Product Page," 2026.
{[}Online{]}. Available:
\url{https://www.nordicsemi.com/Products/nRF52840/Modules}. {[}Accessed:
Mar. 02, 2026{]}.

{[}57{]} Nordic Semiconductor, "nRF52840 Documentation Hub," 2026.
{[}Online{]}. Available:
\url{https://docs.nordicsemi.com/category/nrf52840-category}.
{[}Accessed: Mar. 02, 2026{]}.

{[}58{]} SIMCom Wireless Solutions, "A7670X Series Product Page," 2026.
{[}Online{]}. Available:
\url{https://cn.simcom.com/product/A7670X.html}. {[}Accessed: Mar. 02,
2026{]}.

{[}59{]} Quectel Wireless Solutions, "EC200U Series Product Page," 2026.
{[}Online{]}. Available:
\url{https://www.quectel.com/product/lte-ec200u-series}. {[}Accessed:
Mar. 02, 2026{]}.

{[}60{]} SIMCom Wireless Solutions, "SIM7600CE Product Page," 2026.
{[}Online{]}. Available:
\url{https://en.simcom.com/product/SIM7600CE.html}. {[}Accessed: Mar.
13, 2026{]}.

{[}61{]} Quectel Wireless Solutions, "Quectel EC200U Series LTE Standard
Specification," Version 1.4, 2024. {[}Online{]}. Available:
\url{https://developer.quectel.com/en/wp-content/uploads/sites/2/2024/11/Quectel_EC200U_Series_LTE_Standard_Specification_V1.4.pdf}.
{[}Accessed: Mar. 02, 2026{]}.

{[}62{]} u-blox, "NEO-M8 Data Sheet," UBX-15031086, {[}Online{]}.
Available:
\url{https://www.u-blox.com/sites/default/files/NEO-M8-FW3_DataSheet_UBX-15031086.pdf}.
{[}Accessed: Mar. 02, 2026{]}.

{[}63{]} u-blox, "NEO-M8 Series Product Page," 2026. {[}Online{]}.
Available: \url{https://www.u-blox.com/en/product/neo-m8-series}.
{[}Accessed: Mar. 02, 2026{]}.

{[}64{]} vgate, "Vgate iCar Pro BLE Product Page," 2026. {[}Online{]}.
Available: \url{https://vgatemall.com/products-detail/i-9/}.
{[}Accessed: Mar. 02, 2026{]}.

{[}65{]} Texas Instruments, "TPS2115A Product Page," 2026. {[}Online{]}.
Available: \url{https://www.ti.com/product/TPS2115A}. {[}Accessed: Mar.
02, 2026{]}.

{[}66{]} Infineon Technologies, "IRLML6402 Product Page," 2026.
{[}Online{]}. Available: \url{https://www.infineon.com/part/IRLML6402}.
{[}Accessed: Mar. 02, 2026{]}.

{[}67{]} SONGLE Relay, "SRD-05VDC-SL-C Relay Datasheet," datasheet
mirror. {[}Online{]}. Available:
\url{https://www.alldatasheet.com/html-pdf/1132639/SONGLERELAY/SRD-05VDC-SL-C/1713/2/SRD-05VDC-SL-C.html}.
{[}Accessed: Mar. 02, 2026{]}.

\begin{center}\rule{0.5\linewidth}{0.5pt}\end{center}

\majorheading{PHỤ LỤC - APPENDICES}

\begin{center}\rule{0.5\linewidth}{0.5pt}\end{center}

\appendixsection{PHỤ LỤC 1: BÁO CÁO TÀI CHÍNH - FINANCE REPORT}

\subsection{1.1. Bảng kê chi phí linh kiện (Bill of Materials - BOM)}

{\thesissettableid{PL-1.1}{tab-condensed-app1-bom}
\def\LTcaptype{table}
\begin{longtable}{|L{0.360\linewidth}|L{0.220\linewidth}|L{0.360\linewidth}|}
\caption{Tóm tắt chi phí BOM chính}\label{tab:condensed-app1-bom}\\
\hline \textbf{Nhóm linh kiện} & \textbf{Chi phí ước tính} & \textbf{Ghi chú} \\
\\\hline
\endfirsthead
\caption[]{Tóm tắt chi phí BOM chính (tiếp theo)}\\
\hline \textbf{Nhóm linh kiện} & \textbf{Chi phí ước tính} & \textbf{Ghi chú} \\
\\\hline
\endhead
\\\hline
\endfoot
\endlastfoot
Khối xử lý và truyền thông & Mức cao nhất trong BOM & Bao gồm MCU, modem LTE/GNSS và adapter OBD2 \\\hline
Khối nguồn và bảo vệ & Mức trung bình & Bao gồm buck/LDO/sạc pin và linh kiện bảo vệ \\\hline
Khối cơ khí và phụ trợ & Mức thấp - trung bình & Vỏ hộp, dây nối, linh kiện lắp đặt \\\hline
\end{longtable}
}

\subsection{1.2. Chi phí hạ tầng đám mây (ước tính hàng tháng)}

Chi phí cloud thay đổi theo số lượng thiết bị hoạt động đồng thời. Ở quy mô thí điểm, chi phí duy trì nằm trong ngưỡng phù hợp với mô hình đồ án.

\subsection{1.3. Tổng hợp chi phí dự án}

Tổng chi phí dự án được kiểm soát tốt nhờ tận dụng công nghệ mã nguồn mở và kiến trúc có thể mở rộng theo từng giai đoạn.

{\thesissettableid{PL-1.3}{tab-condensed-app1-cost-scenarios}
\def\LTcaptype{table}
\begin{longtable}{|L{0.220\linewidth}|L{0.220\linewidth}|L{0.220\linewidth}|L{0.300\linewidth}|}
\caption{Kịch bản chi phí theo quy mô triển khai}\label{tab:condensed-app1-cost-scenarios}\\
\hline \textbf{Quy mô} & \textbf{Số thiết bị} & \textbf{Chi phí vận hành/tháng} & \textbf{Nhận xét} \\
\\\hline
\endfirsthead
\caption[]{Kịch bản chi phí theo quy mô triển khai (tiếp theo)}\\
\hline \textbf{Quy mô} & \textbf{Số thiết bị} & \textbf{Chi phí vận hành/tháng} & \textbf{Nhận xét} \\
\\\hline
\endhead
\\\hline
\endfoot
\endlastfoot
Pilot nội bộ & 5--20 & Thấp & Tập trung kiểm chứng kỹ thuật và quy trình vận hành \\\hline
Thí điểm mở rộng & 20--100 & Trung bình & Cần tăng giám sát và chuẩn hóa cảnh báo \\\hline
Triển khai dịch vụ & 100--300 & Trung bình - cao & Bắt đầu cần tối ưu ingestion và dashboard realtime \\\hline
Đội xe lớn & 300--1000 & Cao & Yêu cầu kế hoạch scale ngang theo lớp dịch vụ \\\hline
Đa khu vực & >1000 & Rất cao & Cần tách cụm vận hành, chính sách dữ liệu và quan sát tập trung \\\hline
\end{longtable}
}

\begin{center}\rule{0.5\linewidth}{0.5pt}\end{center}

\appendixsection{PHỤ LỤC 2: CÁC TIÊU CHUẨN THIẾT KẾ - STANDARDS}

\subsection{2.1. Tiêu chuẩn OBD2 (SAE J1979 / ISO 15031--5)}

Áp dụng làm nền cho thiết kế lớp thu thập dữ liệu vận hành xe và chuẩn hóa cách giải nghĩa các PID quan trọng.

\subsection{2.2. MQTT 3.1.1 trong triển khai thiết bị}

Áp dụng cho lớp truyền thông thời gian thực giữa thiết bị và máy chủ, với quy ước topic và QoS theo mức ưu tiên nghiệp vụ.

\subsection{2.3. Bảo mật thông tin (ISO 27001 - tham khảo)}

Áp dụng ở mức nguyên tắc: kiểm soát truy cập, tách quyền theo vai trò và truy vết vận hành.

\subsection{2.4. Thiết kế REST API (RFC 7231 và best practices)}

Áp dụng cho cấu trúc endpoint, mã phản hồi và chuẩn hóa định dạng lỗi trong API quản trị.

{\thesissettableid{PL-2.1}{tab-condensed-app2-standard-mapping}
\def\LTcaptype{table}
\begin{longtable}{|L{0.220\linewidth}|L{0.330\linewidth}|L{0.390\linewidth}|}
\caption{Bản đồ áp dụng chuẩn trong hệ thống}\label{tab:condensed-app2-standard-mapping}\\
\hline \textbf{Chuẩn/nguyên tắc} & \textbf{Phạm vi áp dụng} & \textbf{Cách hiện thực trong dự án} \\
\\\hline
\endfirsthead
\caption[]{Bản đồ áp dụng chuẩn trong hệ thống (tiếp theo)}\\
\hline \textbf{Chuẩn/nguyên tắc} & \textbf{Phạm vi áp dụng} & \textbf{Cách hiện thực trong dự án} \\
\\\hline
\endhead
\\\hline
\endfoot
\endlastfoot
SAE J1979 / ISO 15031 & Dữ liệu OBD2 & Chọn tập PID cốt lõi, chuẩn hóa parser và mô tả đơn vị đo \\\hline
MQTT 3.1.1 & Kênh truyền thiết bị -- cloud & Tách topic theo nghiệp vụ, gán QoS theo mức ưu tiên dữ liệu \\\hline
ISO 27001 (tham khảo) & Vận hành và bảo mật & Tách quyền truy cập, ghi log vận hành và theo dõi phiên truy cập \\\hline
RFC 7231 + REST best practices & API backend & Chuẩn hóa method/status code, phản hồi lỗi nhất quán \\\hline
Nguyên tắc least privilege & Quản trị hệ thống & Giới hạn vai trò theo chức năng vận hành và bảo trì \\\hline
\end{longtable}
}

\begin{center}\rule{0.5\linewidth}{0.5pt}\end{center}

\appendixsection{PHỤ LỤC 3: KẾ HOẠCH THỰC HIỆN - ASSIGNMENT AND TIMELINES}

\subsection{3.1. Phân chia giai đoạn dự án}

Kế hoạch được chia thành ba giai đoạn: thiết kế và dựng nền tảng, tích hợp và kiểm thử, hoàn thiện và đánh giá.

\subsection{3.2. Kế hoạch thực hiện dự án theo giai đoạn}

Mỗi giai đoạn có đầu ra kỹ thuật rõ: phiên bản nguyên mẫu, phiên bản tích hợp liên tầng và phiên bản đánh giá cuối kỳ.

\subsection{3.3. Chi tiết công việc theo giai đoạn}

Chi tiết công việc được chuẩn hóa theo tuần để theo dõi tiến độ và kiểm soát rủi ro phát sinh.

{\thesissettableid{PL-3.1}{tab-condensed-app3-milestones}
\def\LTcaptype{table}
\begin{longtable}{|L{0.180\linewidth}|L{0.220\linewidth}|L{0.260\linewidth}|L{0.280\linewidth}|}
\caption{Milestone triển khai và tiêu chí nghiệm thu}\label{tab:condensed-app3-milestones}\\
\hline \textbf{Milestone} & \textbf{Đầu ra chính} & \textbf{Tiêu chí nghiệm thu} & \textbf{Rủi ro theo dõi} \\
\\\hline
\endfirsthead
\caption[]{Milestone triển khai và tiêu chí nghiệm thu (tiếp theo)}\\
\hline \textbf{Milestone} & \textbf{Đầu ra chính} & \textbf{Tiêu chí nghiệm thu} & \textbf{Rủi ro theo dõi} \\
\\\hline
\endhead
\\\hline
\endfoot
\endlastfoot
M1: Thiết kế nền tảng & Kiến trúc và BOM chốt & Thiết kế qua review kỹ thuật nội bộ & Sai lệch linh kiện/chi phí \\\hline
M2: Dựng nguyên mẫu & Thiết bị chạy mẫu & Thu thập được telemetry cốt lõi & Nhiễu nguồn và reset cục bộ \\\hline
M3: Tích hợp cloud & Luồng dữ liệu end-to-end & Có dữ liệu realtime trên dashboard & Mất kết nối mạng di động \\\hline
M4: Kiểm thử liên tầng & Báo cáo kết quả thử nghiệm & Đạt chỉ số chính theo mục tiêu đồ án & Sai lệch giữa lab và thực địa \\\hline
M5: Hoàn thiện báo cáo & Bản rút gọn + tài liệu bổ trợ & Trình bày đủ logic thiết kế -- triển khai -- đánh giá & Thiếu nhất quán giữa số liệu và mô tả \\\hline
\end{longtable}
}

\begin{center}\rule{0.5\linewidth}{0.5pt}\end{center}

\appendixsection{PHỤ LỤC 4: TÀI LIỆU BỔ TRỢ - SUPPORTING MATERIALS}

\subsection{4.1. Kho mã nguồn của dự án}

Kho mã nguồn được tổ chức theo mô hình tách dịch vụ, hỗ trợ truy vết thay đổi theo từng thành phần hệ thống.

\subsection{4.2. Tham chiếu cơ sở dữ liệu (Lược đồ cơ sở dữ liệu)}

Tài liệu lược đồ dữ liệu được dùng để đối chiếu quan hệ giữa dữ liệu vận hành, dữ liệu sự kiện và dữ liệu cấu hình.

\subsection{4.3. Tài liệu API (API Documentation)}

Tài liệu API dùng cho vận hành dashboard, đồng bộ thiết bị và truy xuất báo cáo.

\subsection{4.4. Tài liệu cấu hình môi trường (Environment Configuration)}

Tài liệu cấu hình mô tả biến môi trường tối thiểu để triển khai backend, frontend và bridge.

\subsection{4.5. Hướng dẫn cài đặt và chạy hệ thống}

Hướng dẫn triển khai giữ dạng checklist để giảm lỗi thao tác khi dựng lại môi trường.

{\thesissettableid{PL-4.1}{tab-condensed-app4-handover-docs}
\def\LTcaptype{table}
\begin{longtable}{|L{0.230\linewidth}|L{0.260\linewidth}|L{0.220\linewidth}|L{0.250\linewidth}|}
\caption{Danh mục tài liệu bàn giao vận hành rút gọn}\label{tab:condensed-app4-handover-docs}\\
\hline \textbf{Nhóm tài liệu} & \textbf{Mục đích} & \textbf{Định dạng} & \textbf{Người sử dụng chính} \\
\\\hline
\endfirsthead
\caption[]{Danh mục tài liệu bàn giao vận hành rút gọn (tiếp theo)}\\
\hline \textbf{Nhóm tài liệu} & \textbf{Mục đích} & \textbf{Định dạng} & \textbf{Người sử dụng chính} \\
\\\hline
\endhead
\\\hline
\endfoot
\endlastfoot
Hướng dẫn dựng môi trường & Khởi tạo hệ thống ban đầu & Markdown/PDF & DevOps, backend engineer \\\hline
Checklist vận hành ngày thường & Theo dõi trạng thái hệ thống & Checklist biểu mẫu & Vận hành ca trực \\\hline
Quy trình xử lý sự cố & Phục hồi khi lỗi liên tầng & Runbook & Kỹ sư vận hành và hỗ trợ kỹ thuật \\\hline
Danh mục biến môi trường & Tránh sai cấu hình dịch vụ & Bảng cấu hình chuẩn & DevOps, bảo mật \\\hline
Tài liệu endpoint cốt lõi & Đồng bộ tích hợp liên hệ thống & API spec rút gọn & Frontend, mobile, đối tác tích hợp \\\hline
\end{longtable}
}

\subsection{4.6. Checklist vận hành chi tiết theo ca trực}

Danh mục dưới đây dùng cho kiểm tra nhanh đầu ca, giữa ca và cuối ca để hạn chế sự cố vận hành lặp lại.

{\thesissettableid{PL-4.2}{tab-condensed-app4-ops-checklist}
\def\LTcaptype{table}
\begin{longtable}{|L{0.110\linewidth}|L{0.330\linewidth}|L{0.170\linewidth}|L{0.330\linewidth}|}
\caption{Checklist vận hành chi tiết theo ca trực}\\label{tab:condensed-app4-ops-checklist}\\
\hline \textbf{STT} & \textbf{Hạng mục kiểm tra} & \textbf{Tần suất} & \textbf{Tiêu chí đạt} \\
\\\hline
\endfirsthead
\caption[]{Checklist vận hành chi tiết theo ca trực (tiếp theo)}\\
\hline \textbf{STT} & \textbf{Hạng mục kiểm tra} & \textbf{Tần suất} & \textbf{Tiêu chí đạt} \\
\\\hline
\endhead
\\\hline
\endfoot
\endlastfoot
1 & Trạng thái dịch vụ broker MQTT & Đầu ca & Không có dịch vụ dừng bất thường \\\hline
2 & Trạng thái MQTT Bridge ingestion & Đầu ca & Không backlog tăng đột biến \\\hline
3 & Trạng thái API backend & Đầu ca & Health check đạt, không lỗi nghiêm trọng \\\hline
4 & Trạng thái frontend dashboard & Đầu ca & Truy cập được, dữ liệu realtime cập nhật \\\hline
5 & Dung lượng lưu trữ PostgreSQL & Đầu ca & Dư dung lượng theo ngưỡng vận hành \\\hline
6 & Dung lượng lưu trữ metrics/logs & Đầu ca & Dư dung lượng theo chính sách retention \\\hline
7 & Tỷ lệ thiết bị online & Mỗi 2 giờ & Không giảm sâu bất thường theo vùng \\\hline
8 & Độ trễ dữ liệu realtime & Mỗi 2 giờ & Nằm trong ngưỡng cảnh báo vận hành \\\hline
9 & Tỷ lệ lỗi API 5xx & Mỗi 2 giờ & Dưới ngưỡng cảnh báo cấu hình \\\hline
10 & Tỷ lệ reconnect MQTT thiết bị & Mỗi 2 giờ & Không xuất hiện cụm lỗi bất thường \\\hline
11 & Tình trạng queue xử lý cảnh báo & Mỗi 2 giờ & Không tồn đọng kéo dài qua nhiều chu kỳ \\\hline
12 & Cảnh báo bảo mật truy cập trái phép & Liên tục & Không có phiên truy cập rủi ro cao chưa xử lý \\\hline
13 & Đồng bộ thời gian hệ thống (NTP) & Đầu ca và cuối ca & Sai lệch trong giới hạn cho phép \\\hline
14 & Kiểm tra sao lưu dữ liệu định kỳ & Cuối ca & Bản sao lưu trong ngày được tạo thành công \\\hline
15 & Kiểm tra phục hồi bản sao lưu mẫu & Hàng tuần & Khôi phục thử nghiệm thành công \\\hline
16 & Kiểm tra nhật ký lỗi ứng dụng & Cuối ca & Lỗi mới được phân loại và gán xử lý \\\hline
17 & Kiểm tra dashboard giám sát hiệu năng & Cuối ca & Không có chỉ số vượt ngưỡng kéo dài \\\hline
18 & Kiểm tra danh sách thiết bị mất liên lạc dài hạn & Cuối ca & Có ticket xử lý cho từng thiết bị bất thường \\\hline
19 & Xác nhận trạng thái ticket sự cố trong ca & Cuối ca & Không còn ticket khẩn chưa người phụ trách \\\hline
20 & Bàn giao ca vận hành & Cuối ca & Biên bản bàn giao đầy đủ sự cố và hành động \\\hline
\end{longtable}
}
\end{document}




